\section{Cote 66: algebraic cycles and the thesis subjects}
\label{sec:B}

We work with rational cycle coefficients unless a statement explicitly says otherwise. A cycle is homologically trivial when its image in the specified Weil cohomology is zero. Numerical triviality means that every complementary-dimensional intersection number is zero. We distinguish the assertion for a particular prime from uniformity in all primes.

\subsection{Question 24: numerical and homological equivalence}


\begin{theorem}[Known, independently rederived: L1 over the complex numbers]\label{B:betti}
Let $X$ be a smooth projective connected complex variety of complex dimension $4$. For every $z\in\mathrm{CH}^2(X)_{\mathbb Q}$, numerical triviality is equivalent to $\operatorname{cl}_B(z)=0$ in $H^4(X,\mathbb Q)$. Here numerical triviality is tested against every $w\in\mathrm{CH}^2(X)_{\mathbb Q}$.
\end{theorem}
\begin{proof}
Fix an ample rational divisor class $h\in H^2(X,\mathbb Q)$ and write $L(\gamma)=h\smile\gamma$. We use hard Lefschetz, the Lefschetz $(1,1)$ theorem and the Hodge--Riemann bilinear relations. Kleiman \cite[\S3, Propositions 3.1, 3.8, Corollary 3.9 and Remarks 3.4, 3.10]{BKleiman1968} already records the conclusion in complex dimension at most four. We give the complete degree-four proof; no general Lefschetz standard conjecture is assumed.

Let $A^2\subset H^4(X,\mathbb Q)$ denote the image of the cycle class map and let $\alpha\in A^2$. Define
\[
 c=\frac{\int_X\alpha h^2}{\int_Xh^4}\in\mathbb Q,
 \qquad
 \beta=(L^2:H^2(X,\mathbb Q)\longrightarrow H^6(X,\mathbb Q))^{-1}
             (L\alpha-c h^3).
\]
These expressions are defined: the denominator is positive, and hard Lefschetz makes $L^2$ an isomorphism. Since $L^2$ shifts Hodge bidegree by $(2,2)$, its inverse takes the rational $(3,3)$ class in parentheses to a rational $(1,1)$ class. The Lefschetz $(1,1)$ theorem therefore makes $\beta$ algebraic. Moreover
\[
 L^3\beta=L^2\alpha-c h^4=0,
\]
because the two top-degree classes have the same integral. Put $\gamma=\alpha-c h^2-L\beta$. Then $L\gamma=0$ and $\gamma$ is algebraic. Thus
\[
 A^2=\mathbb Qh^2\ \oplus\ L\operatorname{NS}(X)_{\mathbb Q,\mathrm{prim}}
                 \ \oplus\ A^2_{\mathrm{prim}}.
\]
The sum is direct by Lefschetz decomposition. It is orthogonal for $Q(a,b)=\int_Xab$: primitiveness gives $L\gamma=0$ and $L^3\beta=0$, which kill every cross term. The form is positive on $\mathbb Qh^2$, negative definite on the middle summand by the Hodge index theorem in degree $2$, and positive definite on real primitive $(2,2)$ classes by Hodge--Riemann. Its restriction to $A^2_{\mathrm{prim}}$ is consequently positive definite. Hence $Q|_{A^2}$ is nondegenerate.

More explicitly, if $\alpha=c h^2+L\beta+\gamma\ne0$, the algebraic class $\alpha'=c h^2-L\beta+\gamma$ satisfies
\[
 Q(\alpha,\alpha')=c^2\!\int_Xh^4-\!\int_Xh^2\beta^2+\!\int_X\gamma^2>0.
\]
A cycle representing $\alpha'$ is a numerical witness. Therefore a numerically trivial cycle has zero Betti class. The reverse implication follows from compatibility of cycle classes with products and the degree map.
\end{proof}


\begin{lemma}[Proved: invariance of numerical triviality under field extension]\label{B:num-basechange}
Let $X$ be smooth projective over a field $k$, let $K/k$ be an extension, and let $z\in\mathrm{CH}^r(X)_{\mathbb Q}$. Then $z$ is numerically trivial over $k$ if and only if its flat pullback $z_K$ is numerically trivial over $K$.
\end{lemma}
\begin{proof}
If $z_K$ is numerically trivial, testing against the pullback of every complementary cycle over $k$, and invariance of intersection degrees under field extension, proves the assertion over $k$.
For the other direction it is enough to consider one integral complementary-dimensional subvariety $W\subset X_K$. Its equations descend to a finitely generated $k$-subalgebra $R\subset K$. After localizing $R$, the closure $\mathcal W\subset X_R$ is flat over the integral scheme $S=\operatorname{Spec}R$. Here is a precise constancy argument that does not require $S$ to be smooth. Write $d=\dim X$. The rational Chern character on the smooth projective variety $X$ identifies $K_0(X)_{\mathbb Q}$ with $\bigoplus_i\mathrm{CH}^i(X)_{\mathbb Q}$. Choose a virtual bundle $v$ with $\operatorname{ch}(v)=z$, concentrated in codimension $r$. Riemann--Roch gives, for every fibre,
\[
 \chi(X_s,v_s\otimes\mathcal O_{W_s})
 =\int_{X_s}z_s\operatorname{ch}(\mathcal O_{W_s})\operatorname{td}(T_{X_s})
 =\deg(z_s\cdot[W_s]).
\]
The last equality holds because the first nonzero Chern-character term of $\mathcal O_{W_s}$ is $[W_s]$ in codimension $d-r$. Each bundle in $v$ tensored with $\mathcal O_{\mathcal W}$ is flat over $S$, so its Euler characteristic is constant in this flat projective family. This proves the required constancy, including inseparable residue extensions. The Chern-character and Riemann--Roch assertions used here are the smooth cases of \cite[\S3.3]{BFulton1975} and \cite[\S0 and Chapter III, \S1]{BBFM1975}. Choose a closed point $s\in S$; its residue field $k(s)$ is finite over $k$. For the finite projection $p:X_{k(s)}\to X$, the projection formula gives
\[
 [k(s):k]\,\deg_{k(s)}(z_{k(s)}\cdot W_s)
  =\deg_k(z\cdot p_*W_s)=0.
\]
Thus the generic intersection degree, and hence $\deg_K(z_K\cdot W)$, is zero. Rational linearity concludes the proof.
\end{proof}

\begin{theorem}[Proved: L1 for every characteristic-zero field]\label{B:charzero}
Let $k$ be any field of characteristic $0$, $X/k$ smooth projective and geometrically connected of dimension $4$, and $\ell$ any prime. For $z\in\mathrm{CH}^2(X)_{\mathbb Q}$,
\[
 z\equiv_{\mathrm{num}}0
 \quad\Longleftrightarrow\quad
 \operatorname{cl}_\ell(z)=0\text{ in }H^4_{\acute et}(X_{\bar k},\mathbb Q_\ell(2)).
\]
For a smooth projective variety without geometric connectedness the same conclusion holds componentwise after a finite extension, and descends by trace.
\end{theorem}
\begin{proof}
Choose a finitely generated subfield $k_0\subset k$ over which $X$, an ample divisor and the finitely many components and coefficients of $z$ are defined. Write these models as $X_0,z_0$. By Lemma~\ref{B:num-basechange}, numerical triviality is equivalent for $z$, $z_0$ and the base change of $z_0$ along any embedding $k_0\hookrightarrow\mathbb C$; such an embedding exists because $k_0$ is finitely generated over $\mathbb Q$.

Geometric extension invariance of smooth proper $\ell$-adic cohomology identifies the classes of $z$ and $z_0$, and the Betti--\'{e}tale comparison isomorphism
\[
 H^4_B(X_0(\mathbb C),\mathbb Q)\otimes_{\mathbb Q}\mathbb Q_\ell
 \simeq H^4_{\acute et}(X_{0,\mathbb C},\mathbb Q_\ell)
\]
identifies their cycle classes (with the conventional Tate twist on the right). The finite-coefficient comparison is Milne \cite[Theorem 21.1]{BMilne2013}; taking the inverse limit over $\ell^n$ and tensoring with $\mathbb Q_\ell$ gives this displayed comparison. Tensoring a rational vector space with $\mathbb Q_\ell$ is faithful. The equivalence now follows from Theorem~\ref{B:betti}.

For the last sentence, a finite extension separates the geometric components. Numerical triviality restricts to each component; apply the already proved assertion there. Conversely, vanishing on all components is vanishing of the direct-sum cohomology class. Finite pullback followed by trace multiplies cycle classes by the nonzero extension degree, so descent is faithful with rational coefficients.
\end{proof}

\begin{remark}[Proved: scope and an integral caveat]\label{B:integral}
The statement concerns rational Betti or $\ell$-adic cohomology. It must not be silently replaced by integral cohomology: a nonzero torsion cycle class is numerically trivial. We make no assertion here that every numerically trivial cycle on every fourfold has zero integral class.
\end{remark}

\begin{proposition}[Refuted inference: the Ancona theorem quoted in the prompt]\label{B:ancona-scope}
Ancona's Theorem~1.3 proves positivity for primitive classes \emph{modulo numerical equivalence}. Its combination with the Lefschetz standard conjecture alone does not give the asserted all-primes homological--numerical comparison. The conclusion stated in Ancona's Theorem~1.6 is comparison for infinitely many primes. This is a correction to the proposed proof, not a counterexample to standard conjecture~$D$.
\end{proposition}
\begin{proof}
These are the actual statements of \cite[Theorems 1.3, 1.6; Remark 3.21]{BAncona2021}. The logical distinction can be seen in linear algebra: let $V=\mathbb Qe\oplus\mathbb Qf$ with $q(e,e)=1$ and $q(e,f)=q(f,f)=0$. The induced form on $V/\operatorname{rad}(q)$ is positive definite, although $\operatorname{rad}(q)=\mathbb Qf\ne0$. Positivity on the quotient cannot establish injectivity into it. Here $V$ models algebraic classes modulo homological equivalence, and the radical models homological classes that are numerically trivial. This algebraic obstruction is explicitly distinguished in Broe \cite[Remark 1.3]{BBroe2026}. The matrix calculation is included in the finite checks.
\end{proof}

\begin{theorem}[Known, with derivation: abelian fourfolds in every characteristic]\label{B:abelian}
Let $A$ be an abelian variety of dimension $4$ over any field $k$ and let $\ell\ne\operatorname{char}k$. Then numerical and $\ell$-adic homological equivalence agree on rational cycles on $A$, in particular on its $2$-cycles. This conclusion uses Broe's 2026 preprint, not only the sources remembered in the prompt.
\end{theorem}
\begin{proof}[Derivation from the opened sources]
Broe \cite[Theorem 1.1]{BBroe2026} proves surjectivity of the codimension-two cycle map onto Frobenius-invariant $H^4_{\acute et}(A_{\bar k},\mathbb Q_\ell(2))$ for every abelian fourfold over a finite field. The divisor case is Tate's theorem; multiplication by the square of an ample divisor carries its algebraic generators to codimension three and is an isomorphism by hard Lefschetz. Degrees zero and four are generated by the fundamental class and a closed point. Thus the full Tate conjecture holds for this $A$.

An abelian variety is of abelian type in Kahn's Definition~1(a). Kahn's Theorem~1, equivalently Theorem~1.10 with the other factor $\operatorname{Spec}k$, now says that rational and numerical equivalence agree with rational coefficients \cite{BKahn2003}. Since rational equivalence implies $\ell$-adic homological equivalence and homological equivalence implies numerical equivalence, all three relations agree over the finite field, for every $\ell$ different from its characteristic.

For arbitrary $k$, descend $A$, a polarization and a tested cycle to a finitely generated field and spread out to an abelian scheme over an integral scheme of finite type over $\mathbb Z[1/\ell]$. A closed fibre is an abelian fourfold over a finite field. Ancona's Corollary~3.19 transfers the equality of the two equivalence relations from this fibre to the generic fibre. For clarity, its mechanism is the following. Lieberman's Lefschetz theorem for abelian varieties, as recorded in Kleiman \cite[Theorem 2A11]{BKleiman1968}, makes the algebraic Lefschetz decomposition available. The later original account also gives this input in Kleiman \cite[Theorem~4-1 and Proposition~4-3(2)(c), pp.~14--16]{BKleiman1994}; its Proposition~5-1, p.~17, formulates the positivity implication on homological algebraic classes. All degrees except primitive codimension two are reduced to divisors. On the special fibre, the primitive pairing on homological algebraic classes is positive definite, because equality with numerical equivalence has just been established and Ancona's Theorem~3.18 supplies positivity there. Smooth proper base change embeds the generic-fibre space of primitive algebraic classes into this space, preserving the pairing. Its restriction remains positive definite, so a primitive generic class orthogonal to all generic algebraic classes is zero, by testing against itself. The other Lefschetz pieces are nondegenerate by the divisor theorem. Extension invariance and finite descent then give the assertion over $k$.

This is also the argument in Broe's Corollaries~1.2 and~8.2; the source is an opened arXiv preprint dated 2026, rather than a claim of independent proof of its Tate theorem.
\end{proof}

\begin{theorem}[Proved: the exact positivity criterion and a further class of fourfolds]\label{B:criterion}
Fix a Weil cohomology theory and let $X$ be a smooth projective fourfold. Suppose its rational algebraic cohomology in degree four admits the algebraic orthogonal Lefschetz decomposition
\[
 A^2(X)=\mathbb Qh^2\oplus L A^1_{\mathrm{prim}}(X)\oplus A^2_{\mathrm{prim}}(X),
\]
and suppose the divisor summand has the usual nondegenerate Hodge index form. Assume also that the primitive intersection form is positive semidefinite. Then numerical and homological equivalence agree in codimension two if and only if the primitive form is positive definite. In particular, the argument applies whenever the Lefschetz standard conjecture and the Hodge standard conjecture \emph{on homological algebraic classes} both hold.

More explicitly, the \emph{class of applicability of the positivity argument} is the class $\mathscr C_H$ of polarized fourfolds with the displayed algebraic decomposition, nondegenerate divisor block and positive definite primitive block. Within the stated semidefinite setting, membership in $\mathscr C_H$ is necessary and sufficient for this argument to prove the comparison. This class includes all characteristic-zero fourfolds for Betti and $\ell$-adic cohomology, all abelian fourfolds using Theorem~\ref{B:abelian}, and smooth projective fourfolds admitting a finite surjective morphism from an abelian fourfold. This determines applicability of the argument; it is not a claim that varieties outside it fail standard conjecture~$D$.

Independently of a positivity formulation, $\ell$-adic comparison in codimension two holds for every smooth projective fourfold $X$ admitting a dominant generically finite morphism $f:A\to X$ from an abelian fourfold, for every prime $\ell\ne\operatorname{char}k$.
\end{theorem}
\begin{proof}
The radical of the intersection form on $A^2(X)$ consists exactly of numerically trivial homological classes. The first two orthogonal summands are nondegenerate, so this radical is the radical of the primitive form. A positive semidefinite rational form has zero radical exactly when it is positive definite: diagonalize it over $\mathbb Q$, or use that a vector with zero square is orthogonal to every vector by applying nonnegativity to $q(v+tw)$ for all $t\in\mathbb Q$. This proves the criterion and identifies precisely the extra injectivity missing from positivity on the numerical quotient.

The characteristic-zero inclusion in $\mathscr C_H$ was proved in Theorems~\ref{B:betti} and~\ref{B:charzero}; the abelian inclusion is the positive-definite primitive pairing established in the proof of Theorem~\ref{B:abelian}. If $f:A\to X$ is finite surjective, an ample class $h$ on $X$ pulls back to an ample class on $A$. Pullback takes primitive classes to primitive classes and is injective by $f_*f^*=(\deg f)\operatorname{id}$. The equality $q_A(f^*\alpha,f^*\alpha)=(\deg f)q_X(\alpha,\alpha)$ gives primitive positivity on $X$. The algebraic decomposition descends by decomposing $f^*\alpha$ on $A$ and applying $(\deg f)^{-1}f_*$: the projection formula makes $f_*$ commute with multiplication by the chosen ample classes, and hence preserve the required primitive conditions. The divisor Hodge index theorem supplies the divisor block. Thus $X$ belongs to $\mathscr C_H$ as asserted.

For the second assertion, let $z$ be numerically trivial on $X$. For every complementary cycle $w$ on $A$, the projection formula gives $\deg(f^*z\cdot w)=\deg(z\cdot f_*w)=0$. Thus $f^*z$ is numerically trivial on $A$ and Theorem~\ref{B:abelian} gives $f^*\operatorname{cl}(z)=0$. The cohomological push--pull identity $f_*f^*=(\deg f)\operatorname{id}$, valid for this proper generically finite morphism between smooth projective varieties of equal dimension, and the nonzero rational scalar $\deg f$, imply $\operatorname{cl}(z)=0$. The converse holds for every smooth projective variety by compatibility of intersections with cycle classes.
\end{proof}

\subsection{Question 24, first question: the meaning of \texorpdfstring{$\tau$}{tau}}

\begin{theorem}[Refuted: numerical triviality need not imply $\tau$-triviality]\label{B:tau}
The original image of cote~66, scan~4, reads $\tau$-equivalence; this is the standard relation defined in Kleiman \cite[\S1.1, p.~362]{BKleiman1968}: a nonzero integer multiple is algebraically equivalent to zero. There is a smooth projective connected complex threefold $X$ and an integral one-cycle $z$ which is homologically, hence numerically, trivial, but for which $mz$ is not algebraically trivial for any integer $m\ne0$.
\end{theorem}
\begin{proof}
Take a sufficiently general smooth quintic hypersurface $X\subset\mathbb P^4_{\mathbb C}$ to which Ceresa--Collino \cite[\S III, (3.3)]{BCeresaCollino1983} applies, and choose two distinct lines $\ell_1,\ell_2\subset X$. The cited statement asserts linear independence of the $2875$ line classes modulo algebraic equivalence, so in particular $m(\ell_1-\ell_2)$ is not algebraically trivial for any nonzero integer $m$. This is a theorem for a general member, in the paper's stated sense of avoiding a countable union of proper analytic exceptional loci; we do not replace it by a claim for every quintic or for the Fermat quintic.

Set $z=[\ell_1]-[\ell_2]$. The weak Lefschetz theorem gives $H^2(X,\mathbb Q)=\mathbb Qh$; Poincar\'e duality gives $\dim H^4(X,\mathbb Q)=1$. Both lines have intersection number $1$ with $h$, so their degree-four classes agree. Thus $\operatorname{cl}_B(z)=0$, and $z$ is numerically trivial. The hypotheses requested on the variety hold because a smooth hypersurface in projective space is smooth, projective and connected, with dimension $3$.

The historical margin attributes the negative answer to Griffiths; Ceresa--Collino's introduction and \S I explicitly identify Griffiths's inversion theorem and \S14 of his 1969 paper as the earlier proof, and provide the sharper line-independence statement used here. We cite the primary proof actually opened. The original scan has been checked, resolving the transcription's $\tau/t$ uncertainty in favour of $\tau$. The same construction also refutes the stronger assertion with algebraic equivalence itself.
\end{proof}

\subsection{Question 25: an explicit square on a product of elliptic curves}

The checked image, cote~66, scan~5, asks: ``Le produit de deux cycles alg. \'equiv. \`a $0$ est-il rat. \'equiv. \`a $0$ ? (\textnormal{``pb du carr\'e''}).'' It immediately continues: ``Un $0$-cycle qui est alb. \'equiv. \`a $0$ est-il rationnellement \'equivalent \`a $0$ ?'' These are two distinct assertions; the same construction below refutes both. ``Alb.'' means the kernel of the Albanese map on degree-zero cycles, not algebraic equivalence.

\begin{lemma}[Proved: the trace-form consequence of Mumford's argument]\label{B:trace}
Let $S$ be a smooth projective complex surface and $\omega\in H^0(S,\Omega_S^2)$. If $T$ is a smooth irreducible complex variety and $F,G:T\to\operatorname{Sym}^n S$ are morphisms whose effective zero-cycles $F(t),G(t)$ are rationally equivalent for every $t\in T(\mathbb C)$, then their induced trace forms satisfy $\omega_F=\omega_G\in H^0(T,\Omega_T^2)$.
\end{lemma}
\begin{proof}
Mumford \cite[\S1 and \S2, Lemmas 2--3, pp.~197--203]{BMumford1969} constructs the trace form functorially, proves additivity under addition of effective cycles, and expresses the rational-equivalence relation by countably many closed parameter spaces carrying families of rational curves after auxiliary effective cycles are added. Apply his Lemma~3 to $(F,G):T\to\operatorname{Sym}^nS\times\operatorname{Sym}^nS$. Since the complex numbers are uncountable, some one of these parameter spaces dominates $T$. After taking an irreducible dominating component and resolving it, there exist a smooth variety $T'$, a dominant morphism $p:T'\to T$, a morphism $C:T'\to\operatorname{Sym}^mS$ and a morphism
\[
 H:T'\times\mathbb P^1\longrightarrow\operatorname{Sym}^{n+m}S
\]
with $H(t,0)=F(p(t))+C(t)$ and $H(t,\infty)=G(p(t))+C(t)$. If resolution or a restriction is needed, the argument takes place over a dense open subset and then extends by regularity.

Every regular two-form on $T'\times\mathbb P^1$ is pulled back from $T'$, since $H^0(\mathbb P^1,\mathcal O)=\mathbb C$ and $H^0(\mathbb P^1,\Omega^1)=0$. Therefore the restrictions of $\omega_H$ at $0$ and $\infty$ coincide. Functoriality and additivity give $p^*\omega_F+\omega_C=p^*\omega_G+\omega_C$. Pullback of rational differential forms by a dominant morphism of characteristic-zero varieties is injective, so $\omega_F=\omega_G$ on a dense open subset, and hence everywhere. This is precisely the comparison form of Mumford's proof; no constancy of either family is being assumed.
\end{proof}

\begin{theorem}[Refuted, with a complete counterexample to both questions in 25]\label{B:square}
Let $E/\mathbb C$ be the elliptic curve $y^2=x^3-x$, with origin $O$ at infinity, and put $S=E\times E$. Choose $u,v\in\mathbb C$ algebraically independent over $\overline{\mathbb Q}$, choose roots $a^2=u^3-u$ and $b^2=v^3-v$, and set $P=(u,a)$, $Q=(v,b)$. Then
\[
 z=[(P,Q)]-[(P,O)]-[(O,Q)]+[(O,O)]\in\mathrm{CH}_0(S)
\]
has infinite order and lies in the Albanese kernel. The divisor
\[
 D=\operatorname{pr}_1^*([P]-[O])+\operatorname{pr}_2^*([Q]-[O])\in\mathrm{CH}^1(S)
\]
is algebraically trivial, but $D^2=2z\ne0$ in $\mathrm{CH}_0(S)$.
\end{theorem}
\begin{proof}
The cubic $x^3-x$ has distinct roots, so its smooth projective completion is an elliptic curve; $S$ is a smooth projective connected surface. The chosen pair $(P,Q)$ is a geometric generic point of $E\times E$ over $\overline{\mathbb Q}$, since the projection to the two $x$-coordinates is finite and $u,v$ are algebraically independent.

Let $\eta=dx/y$, a nonzero regular differential on $E$, and $\omega=\operatorname{pr}_1^*\eta\wedge\operatorname{pr}_2^*\eta$, a nonzero regular two-form on $S$. Define morphisms
\[
 F(R,T)=[(R,T)]+[(O,O)],\qquad
 G(R,T)=[(R,O)]+[(O,T)]
\]
from $S$ to $\operatorname{Sym}^2S$. The trace form of $F$ is $\omega$, whereas that of $G$ is zero because its two moving summands factor through curves.

Suppose $mz=0$ for some integer $m>0$. Mumford's Lemma~3 gives a countable union of closed algebraic loci in $S$ on which $mF$ and $mG$ are rationally equivalent. These loci are defined over $\overline{\mathbb Q}$: the construction uses the equations of $S,F,G$, symmetric powers, Chow varieties, bounded-degree maps from $\mathbb P^1$, and projections, all of which are defined over that field. The generic point $(P,Q)$ belongs to one of those loci, so that locus must be all of $S$. Consequently $mF(s)\sim_{\mathrm{rat}}mG(s)$ for every complex point $s$. Lemma~\ref{B:trace} and additivity of trace forms then give $m\omega=0$, a contradiction. Thus $z$ has infinite order. The use of a generic pair is essential; the theorem does not assert nonvanishing for arbitrary pairs of nontorsion points.

The Albanese of $S$ is $E\times E$, with the degree-zero cycle map obtained by summing the points with multiplicities, as in Bloch \cite[Lecture~1, Proposition (1.1)]{BBloch2010}. The four terms of $z$ sum to $(P,Q)-(P,O)-(O,Q)+(O,O)=0$, so $z$ lies in its kernel. Each degree-zero divisor $[P]-[O]$ and $[Q]-[O]$ belongs to $\operatorname{Pic}^0(E)$, hence is algebraically trivial. Pullback and addition preserve algebraic equivalence, so $D$ is algebraically trivial. In the Chow ring of $S$, the squares of the two pullback summands vanish because $\mathrm{CH}^2(E)=0$, while their product is exactly $z$. Bilinearity gives $D^2=2z$, which is nonzero because $z$ has infinite order. This verifies both requested failures with integral cycles as well as with rational coefficients.
\end{proof}

\begin{proposition}[Proved: the correct product statement on $E\times E$]\label{B:alb-kernel}
For any complex elliptic curve $E$, the Albanese kernel $T(E\times E)$ is generated, as an abelian group, by the exterior products
\[
 ([P]-[O])\times([Q]-[O])
\quad (P,Q\in E(\mathbb C)).
\]
Equivalently, the product map from $\mathrm{CH}_0(E)^0\otimes_{\mathbb Z}\mathrm{CH}_0(E)^0$ onto $T(E\times E)$ is surjective. Products of algebraically trivial cycles are algebraically, and thus homologically and numerically, trivial; rational triviality is the false additional conclusion in question~25.
\end{proposition}
\begin{proof}
For every point $(P,Q)$, subtracting the two coordinate-axis points and adding the origin gives the displayed exterior product. Thus for any degree-zero cycle $w=\sum n_i[(P_i,Q_i)]$ we have
\[
 w=\sum n_i z(P_i,Q_i)
  +i_{1*}\!\left(\sum n_i[P_i]\right)
  +i_{2*}\!\left(\sum n_i[Q_i]\right),
\]
where $i_1(P)=(P,O)$ and $i_2(Q)=(O,Q)$; the omitted origin coefficient is zero because $\sum n_i=0$. If $w$ is Albanese-trivial, both degree-zero divisors in parentheses have zero class in the Jacobian $E$, hence are principal by Abel's theorem. Their pushforwards are rationally trivial, proving generation. Every generator has zero Albanese image, so equality holds. Finally, intersecting or externally multiplying a family which connects an algebraically trivial cycle to zero with any fixed cycle gives a family connecting the product to zero. Algebraic equivalence therefore has the asserted ideal property, and cycle-class and degree compatibility give the remaining implications.
\end{proof}

\subsection{Subject 21: the closed-pair criterion and the affine comparison}

\begin{theorem}[Known, with derivation: the affine answer to subject 21]\label{B:henselian}
Let $A$ be a commutative ring and $I\subset A$ an ideal, with $X=\operatorname{Spec}A$ and $Z=\operatorname{Spec}(A/I)$. The following are equivalent:
\begin{enumerate}
\item $(A,I)$ is a henselian pair;
\item $\Gamma(X,\mathcal F)\to\Gamma(Z,\mathcal F|_Z)$ is bijective for every sheaf of sets on $X_{\acute et}$.
\end{enumerate}
Under these equivalent conditions, for every torsion abelian \'etale sheaf $\mathcal F$ and every integer $q\ge0$, restriction is an isomorphism
\[
 H^q_{\acute et}(X,\mathcal F)\longrightarrow H^q_{\acute et}(Z,\mathcal F|_Z).
\]
No noetherian hypothesis or invertibility of the torsion orders is needed. In particular the assertion holds if $A\simeq\varprojlim_n A/I^n$, and for the henselization $A^h$ of $A$ along $I$.
\end{theorem}
\begin{proof}[Derivation and identification of the manuscript's cases]
The implication (1)$\Rightarrow$(2) is Stacks Project Lemma~59.82.6, tag~09ZH. Conversely, apply (2) to the sheaf represented by each \'etale morphism $\operatorname{Spec}B\to X$. Every $A$-algebra map $B\to A/I$ then lifts to $B\to A$. This is condition~(2) of Stacks Project Lemma~15.11.6, which is equivalent to $(A,I)$ being henselian. Thus the comparison in degree zero is sufficient when it is required for \emph{every} \'etale sheaf of sets.

The higher-cohomology conclusion is Gabber's affine proper base-change theorem, in the precise form and proof of Stacks Project Theorem~59.82.7, tag~09ZI. We have opened this proof, and the primary publication is \cite{BGabber1994}. Its torsion hypothesis is part of the theorem; the unrestricted word ``sheaf'' in the manuscript cannot remove it.

For the complete case, $1+a$ is a unit for every $a\in I$, by the convergent formal geometric series, so $I$ lies in the Jacobson radical. A coprime monic factorization modulo $I$ lifts successively modulo $I^n$, and completeness gives a factorization over $A$; this is precisely Stacks Project Lemma~15.11.4. Consequently $(A,I)$ is henselian and the theorem applies. The \'etale localization named in the manuscript is realized by the filtered colimit $A^h$ of \'etale $A$-algebras inducing the prescribed isomorphism modulo $I$. Stacks Project Lemma~15.12.1 gives $(A^h,IA^h)$ henselian and $A^h/IA^h\simeq A/I$. Applying Gabber to that pair is exactly the requested cohomology comparison after localization.
\end{proof}

\begin{theorem}[Known, with derivation: the general degree-zero criterion]\label{B:henselian-general}
Let $X$ be a quasi-compact quasi-separated scheme and $i:Z\hookrightarrow X$ a closed immersion. The following conditions are equivalent:
\begin{enumerate}
\item $H^0(X,\mathcal F)\to H^0(Z,i^*\mathcal F)$ is bijective for every \'etale sheaf of sets $\mathcal F$;
\item for every finite $X$-scheme $X'$, restriction induces a bijection between the clopen subsets of $X'$ and those of $X'\times_X Z$.
\end{enumerate}
Condition~(2) is the geometric definition of a \emph{henselian pair} $(X,Z)$. If these conditions hold, then
\[
 \bigl(B,J\bigr)=\bigl(\Gamma(X,\mathcal O_X),
 \ker[\Gamma(X,\mathcal O_X)\to\Gamma(Z,\mathcal O_Z)]\bigr)
\]
is a henselian ring pair. If, additionally, $X$ is proper over a noetherian ring $A$ and $Z=X\times_A A/I$ for some ideal $I\subset A$, the comparison is an isomorphism in every degree for every torsion abelian \'etale sheaf on $X$.
\end{theorem}
\begin{proof}[Derivation and the proper reduction]
Hall--Rydh \cite[Proposition 5.4 and Definition 5.5]{BHallRydh2023} prove the equivalence, even for algebraic stacks. Here the closed subscheme $Z$ is again quasi-compact and quasi-separated, so their hypotheses apply. To see the finite-cover test explicitly, take the finite direct image of the constant two-element sheaf: its sections classify clopen subsets, and finite base change identifies its restriction. Conversely, their Propositions 5.2--5.3 express every sheaf as a filtered colimit of constructible sheaves, each admitting an injection into a finite product of such finite direct images with finite constant fibres. For an injection $\mathcal F\hookrightarrow\mathcal G$, the equalizer diagram $\mathcal F\to\mathcal G\rightrightarrows\mathcal G\amalg_{\mathcal F}\mathcal G$ yields surjectivity from injectivity and the comparison for $\mathcal G$. This proves why condition~(2) suffices.

For the ring assertion, let $B\to C$ be \'etale with a $B$-map $C\to B/J$. The inclusion $B/J\hookrightarrow\Gamma(Z,\mathcal O_Z)$ gives a section over $Z$ of $X\times_{\operatorname{Spec}B}\operatorname{Spec}C\to X$. Condition~(1) lifts it to a section over $X$, hence a map $C\to B$. Its reduction agrees with the prescribed map because $B/J$ injects into $\Gamma(Z,\mathcal O_Z)$. The \'etale section-lifting criterion, Stacks Lemma 15.11.6, proves henselianity. This is also Pippi's Lemma~8 \cite{BPippi2017}.

For the final assertion, the map of pairs $(A,I)\to(B,J)$ factors through the henselization $(A^h,IA^h)$, giving $g:X\to\operatorname{Spec}A^h$. The composite of $g$ with the separated map $\operatorname{Spec}A^h\to\operatorname{Spec}A$ is proper; the graph of $g$ into $X\times_A A^h$ is closed, so $g$ is proper. Moreover $Z=X\times_{A^h}A^h/IA^h$. Proper base change identifies the restriction of $Rg_*\mathcal F$ to $\operatorname{Spec}(A/I)$ with $Rg_{Z*}i^*\mathcal F$. Applying Theorem~\ref{B:henselian} to each torsion cohomology sheaf $R^qg_*\mathcal F$, the two Leray spectral sequences have identical $E_2$ pages and hence identical abutments. They are first-quadrant spectral sequences, so comparison degree by degree is valid without a uniform cohomological-dimension assumption. This is precisely Pippi's Corollary~1 and Theorem~3, with the reduction displayed.
\end{proof}

The general first sentence of subject~21 does not assume affineness, properness, or quasi-compactness. The theorem identifies its universal degree-zero hypothesis exactly in the stated quasi-compact quasi-separated setting; that finiteness hypothesis is not present on the manuscript page and is not discarded. The implication from universal degree-zero comparison to the finite-clopen criterion uses only finite base change and remains valid without that hypothesis; the converse in that wider generality is not supplied by the cited theorem. The higher-cohomology implication with torsion coefficients for an arbitrary nonproper henselian pair is refuted by Theorem~\ref{C:subject21-torsion-counterexample} below; the affine theorem retains its stated scope. Without the torsion restriction, Proposition~\ref{B:nontorsion-rigidity} below refutes the implication even for an affine henselian pair. Pippi's introduction explicitly formulates the torsion implication as the question of SGA~4, Expos\'e XII, 6.13; his Theorem~3 settles the proper case just proved.

\begin{proposition}[Proved: the exact supported-cohomology obstruction]\label{B:H1-gate}
Let $i:Z\hookrightarrow X$ be any closed immersion of schemes, and let $j:U=X\setminus Z\hookrightarrow X$ be its complement. Assume that $H^0(X,\mathcal F)\to H^0(Z,i^*\mathcal F)$ is bijective for every \'etale sheaf of sets $\mathcal F$. Then restriction is fully faithful on finite \'etale covers and is injective on the nonabelian pointed sets $H^1(X,G)\to H^1(Z,i^*G)$ for every \'etale sheaf of groups $G$. Moreover the following conditions are equivalent:
\begin{enumerate}
\item restriction is an isomorphism in every degree for every torsion abelian \'etale sheaf on $X$;
\item $H^q(X,j_!\mathcal G)=0$ for every torsion abelian \'etale sheaf $\mathcal G$ on $U$ and every $q\ge2$.
\end{enumerate}
The analogous vanishings in degrees zero and one already follow from the hypothesis. In particular the full torsion comparison holds whenever the torsion \'etale cohomological dimension of $X$ is at most one. None of these assertions needs quasi-compactness.
\end{proposition}
\begin{proof}
If two $G$-torsors become isomorphic over $Z$, their sheaf of equivariant isomorphisms has a section over $Z$. Formation of this sheaf commutes with inverse image: this can be checked after the covers trivializing both torsors, where it is the sheaf $G$. The assumed $H^0$ comparison gives a section over $X$, proving injectivity on $H^1$. For finite \'etale covers $V,W$, their internal Hom is locally a finite product of the sheaf represented by $W$, so it too commutes with inverse image. The same $H^0$ comparison proves the bijection on all morphisms and hence full faithfulness.

Extension by zero of an abelian \'etale sheaf along $j$ is exact and has zero restriction to $Z$. The hypothesis therefore gives $H^0(X,j_!\mathcal G)=0$, and the just-proved degree-one injectivity gives $H^1(X,j_!\mathcal G)=0$. For any abelian sheaf $\mathcal F$ on $X$, the localization sequence (Stacks Lemma~59.70.8, tag~095L) is
\[
 0\longrightarrow j_!j^*\mathcal F\longrightarrow\mathcal F
 \longrightarrow i_*i^*\mathcal F\longrightarrow0.
\]
Exactness is checked on geometric stalks: the first term is the stalk of $\mathcal F$ on $U$ and zero on $Z$, while the last is zero on $U$ and the stalk of $\mathcal F$ on $Z$. The same stalk description, Stacks Lemma~59.46.3, makes $i_*$ exact for a closed immersion; its exact left adjoint $i^*$ ensures that it preserves injectives. Hence $H^q(X,i_*\mathcal H)=H^q(Z,\mathcal H)$. The long exact sequence proves (2)$\Rightarrow$(1). Conversely apply (1) to $\mathcal F=j_!\mathcal G$, whose restriction is zero. If $X$ has torsion cohomological dimension at most one, condition~(2) is automatic. This identifies the unproved step for arbitrary pairs as a specific family of supported cohomology groups, with degrees zero and one already eliminated.
\end{proof}

\begin{proposition}[Known, with complete derivation: a cyclic-torsor test for the degree-two obstruction]\label{C:torsion-degree-two-reduction}
Let $X$ be a quasi-compact quasi-separated scheme, let
$i:Z\hookrightarrow X$ be a closed immersion, and let
$j:U=X\setminus Z\hookrightarrow X$ be its open complement.
Assume that
\[
 H^0(X,\mathcal F)\longrightarrow H^0(Z,i^{-1}\mathcal F)
\]
is bijective for every \'etale sheaf of sets $\mathcal F$ on $X$.
The following are equivalent:
\begin{enumerate}
\item $H^2(X,j_!\mathcal G)=0$ for every torsion abelian \'etale
sheaf $\mathcal G$ on $U$.
\item For every torsion abelian \'etale sheaf $\mathcal F$ on $X$,
restriction on $H^1$ is surjective, and hence bijective.
\item For every finite morphism $f:Y\to X$, every prime $\ell$, and
every integer $\nu>0$, every torsor under the constant group
$\mathbb Z/\ell^\nu\mathbb Z$ on $Y\times_XZ$ is the restriction of
such a torsor on $Y$.
\end{enumerate}
In condition~(3), it suffices to test finite morphisms of finite
presentation. If these conditions hold, restriction on $H^2$ is
injective for every torsion abelian sheaf on $X$.
\end{proposition}
\begin{proof}
The equivalence between the first-degree comparison and the finite
constant-coefficient test is the $n=1$ case of SGA~4, Expos\'e~XII,
Proposition~6.5(iii), with the transfer argument also stated in
Corollary~6.6 \cite[Expos\'e XII, 6.5--6.6; original pp.~103--111]{BSGA4III1973}. We give the details, including the dimension shift
relating that comparison to the displayed supported-cohomology group.

The degree-zero hypothesis gives injectivity of
$H^1(X,\mathcal H)\to H^1(Z,i^{-1}\mathcal H)$ for every abelian
sheaf $\mathcal H$, by the torsor-isomorphism argument of
Proposition~\ref{B:H1-gate}. The localization short exact sequence
\[
 0\longrightarrow j_!j^{-1}\mathcal F\longrightarrow\mathcal F
 \longrightarrow i_*i^{-1}\mathcal F\longrightarrow0
\]
therefore proves $(1)\Rightarrow(2)$.

We first record a torsion-preserving acyclic embedding, to justify
the converse without presuming that restriction preserves injectives.
Embed a torsion abelian sheaf $\mathcal K$ into an injective abelian
sheaf $\mathcal I$ on $X_{\acute et}$. For every positive integer $m$,
multiplication by $m$ on $\mathcal I$ is surjective on sections over
each \'etale $V\to X$: apply injectivity to multiplication by $m$ on
the free abelian sheaf generated by $V$, which is a monomorphism.
Thus
\[
 0\longrightarrow\mathcal I[m]\longrightarrow\mathcal I
 \xrightarrow{m}\mathcal I\longrightarrow0
\]
is exact and its map on global sections is surjective. Its long exact
sequence gives $H^q(X,\mathcal I[m])=0$ for all $q>0$.
The torsion subsheaf
\[
 \mathcal T=\varinjlim_{m\mid m'}\mathcal I[m]
\]
contains $\mathcal K$. Since \'etale cohomology on a qcqs scheme
commutes with filtered colimits (Stacks Lemma~59.51.4, tag~03Q5),
$H^q(X,\mathcal T)=0$ for all $q>0$. Both $\mathcal T$ and
$\mathcal Q=\mathcal T/\mathcal K$ are torsion sheaves.

Assume~(2) and apply this construction to $\mathcal K=j_!\mathcal G$.
Since $i^{-1}\mathcal K=0$, exactness of inverse image gives
$i^{-1}\mathcal Q\simeq i^{-1}\mathcal T$.
Condition~(2) for $\mathcal T$ gives
$H^1(Z,i^{-1}\mathcal T)=0$. Injectivity of first-degree restriction
for $\mathcal Q$ then gives $H^1(X,\mathcal Q)=0$. The long exact
sequence for $0\to\mathcal K\to\mathcal T\to\mathcal Q\to0$ gives
$H^2(X,\mathcal K)=0$, proving $(2)\Rightarrow(1)$.

For a finite morphism $f:Y\to X$, finite direct image is exact and
commutes with arbitrary base change. Consequently, for any finite
abelian group $M$,
\[
 \begin{aligned}
 H^1(X,f_*\underline M)&=H^1(Y,\underline M),\\
 H^1(Z,i^{-1}f_*\underline M)
 &=H^1(Y\times_X Z,\underline M).
 \end{aligned}
\]
These are Stacks Proposition~59.55.2 and Lemma~59.55.3.
Taking $M=\mathbb Z/\ell^\nu\mathbb Z$ proves $(2)\Rightarrow(3)$.

Assume~(3), even only for finite morphisms of finite presentation.
First let $\mathcal F$ be a constructible abelian sheaf. By Stacks
Lemma~59.74.4, tag~09Z7, there is a monomorphism
\[
 \mathcal F\lhook\joinrel\longrightarrow
 \mathcal E=\bigoplus_{a=1}^{r}f_{a*}\underline{M_a},
\]
where each $f_a$ is finite and finitely presented and each $M_a$ is a
finite abelian group. Finite abelian groups are finite direct sums
of cyclic prime-power groups, so~(3) and finite base change give
surjectivity of first-degree restriction for $\mathcal E$.
Put $\mathcal Q=\mathcal E/\mathcal F$.

Take $\alpha\in H^1(Z,i^{-1}\mathcal F)$. Its image in
$H^1(Z,i^{-1}\mathcal E)$ lifts to $\beta\in H^1(X,\mathcal E)$.
The image of $\beta$ in $H^1(X,\mathcal Q)$ restricts to zero, so is
zero by the already established first-degree injectivity for
$\mathcal Q$. Hence $\beta$ lifts to $\gamma\in H^1(X,\mathcal F)$.
The difference $\alpha-i^{-1}\gamma$ maps to zero in
$H^1(Z,i^{-1}\mathcal E)$ and therefore equals the connecting image
of some $q\in H^0(Z,i^{-1}\mathcal Q)$. Lift $q$ to
$\widetilde q\in H^0(X,\mathcal Q)$ using the degree-zero hypothesis.
Then $\gamma+\delta(\widetilde q)$ restricts to $\alpha$.
This proves surjectivity for every constructible $\mathcal F$.

Every torsion abelian \'etale sheaf on a qcqs scheme is a filtered
colimit of constructible abelian sheaves (Stacks Lemma~59.73.2,
tag~03SA). Inverse image preserves these colimits, and both $X$ and
its closed subscheme $Z$ are qcqs, so their $H^1$ commute with them.
Thus the preceding surjectivity holds for every torsion abelian
$\mathcal F$, proving $(3)\Rightarrow(2)$.

Finally, use an acyclic torsion embedding
$0\to\mathcal F\to\mathcal T\to\mathcal Q\to0$ as above.
Condition~(2) gives $H^1(Z,i^{-1}\mathcal T)=0$.
If $\xi\in H^2(X,\mathcal F)$ restricts to zero, choose
$\eta\in H^1(X,\mathcal Q)$ with $\delta\eta=\xi$.
The restriction of $\eta$ lies in the image of
$H^1(Z,i^{-1}\mathcal T)=0$, so is zero; first-degree injectivity
for $\mathcal Q$ gives $\eta=0$, and hence $\xi=0$.
\end{proof}

The proposition does not prove that its equivalent conditions follow
from the universal degree-zero hypothesis. It identifies the first
torsion obstruction as extension of cyclic torsors on
closed fibres of finite covers. Full faithfulness of finite \'etale
restriction, already proved from degree zero, supplies uniqueness
of an extension when it exists; it does not supply existence.
Theorem~\ref{C:subject21-torsion-counterexample} below proves actual
failure of existence in the unrestricted qcqs setting.


\begin{lemma}[Proved: a flat finite-type henselian pair over a trait is proper]\label{C:henselian-flat-trait-proper}
Let $R$ be a henselian discrete valuation ring, with fraction field
$K$, residue field $k$, generic point $\eta$ and closed point $s$.
Let $X$ be a flat separated scheme of finite type over $R$.
Assume that restriction
\[
 H^0(X,\mathcal F)\longrightarrow H^0(X_s,\mathcal F|_{X_s})
\]
is bijective for every \'etale sheaf of sets $\mathcal F$ on $X$.
Then $X$ is proper over $R$. Neither excellence nor completeness of
$R$ is required. The same conclusion holds under the weaker hypothesis
that every closed point of $X$ lies in $X_s$.
\end{lemma}
\begin{proof}
The hypothesis implies that every open subset of $X$ containing
$X_s$ is all of $X$: the sheaf represented by such an open has a
section over $X_s$, and a section of an open immersion over $X$
exists only when the open immersion is an isomorphism.
In particular every closed point of $X$ lies in $X_s$; this weaker
condition is all the rest of the proof uses.

By Nagata compactification (Stacks Theorem~38.33.8, tag~0F41),
choose an open immersion $X\hookrightarrow\overline X$ with
$\overline X$ proper over $R$. Replace $\overline X$ by the
scheme-theoretic closure of $X$. Flatness of $X$ over the DVR implies
that $X_\eta$ is dense in $X$; consequently $X_\eta$ is dense in
$\overline X$. In particular every irreducible component of
$\overline X$ dominates $\operatorname{Spec}R$.

Suppose that the boundary $D=\overline X\setminus X$ is nonempty.
It is proper over $R$ with nonempty closed image, so $D_s$ is nonempty.
Choose a closed point $z$ of $D_s$; it is also closed in
$\overline X_s$. Some generic point $\xi$ of an irreducible component
of $\overline X$ specializes to $z$, and $\xi\in X_\eta$.
The density theorem for closed points in fibres, Stacks
Lemma~37.52.4 (tag~053U), applied to this specialization from the
generic to the closed fibre, provides a point $y\in X_\eta$ that is
closed in $\overline X_\eta$ and specializes to $z$. Indeed the closed
generic-fibre points specializing to $z$ are dense in the generic
fibre of $\overline{\{\xi\}}$, and $X_\eta$ meets this fibre in a
nonempty open subset.

Let $C$ be the reduced closure of $y$ in $\overline X$. It is integral
and proper over $R$, and $C_\eta=\operatorname{Spec}\kappa(y)$ with
$[\kappa(y):K]<\infty$. The special fibre of $C$ is zero-dimensional.
To check this last assertion without a hidden normality or excellence
assumption, take any affine open $\operatorname{Spec}A\subset C$
meeting the special fibre. The domain $A$ contains $R$, is of finite
type over $R$, and $A\otimes_RK$ has dimension zero. Stacks
Lemma~10.125.9 (tag~00QK), the dimension-of-fibres theorem over a
valuation ring, gives
\[
 \dim(A\otimes_R k)=\dim(A\otimes_RK)=0.
\]
Thus $C\to\operatorname{Spec}R$ is proper with finite fibres, and is
finite by Stacks Lemma~37.44.1 (tag~02LS).

Write $C=\operatorname{Spec}A_C$. Since $R$ is henselian local, a
finite $R$-algebra is a finite product of local rings
(Stacks Lemma~10.153.4, tag~04GH). Since $A_C$
is a domain, there is only one factor. Therefore $C$ has a unique
closed point, necessarily $z$. Its generic fibre consists of $y$,
and its special fibre is supported at $z$. As $z\notin X$ and $y\in X$,
\[
 C\cap X=\{y\}
\]
is a nonempty closed subset of $X$ disjoint from $X_s$. Its complement
is a proper open subset of $X$ containing $X_s$, contradicting the
first paragraph; equivalently, $y$ is a closed point outside $X_s$.
Hence $D$ is empty and $X=\overline X$ is proper.
\end{proof}

\begin{theorem}[Proved: torsion comparison for every separated finite-type pair over a henselian trait]\label{C:henselian-trait-comparison}
Let $R$ be a henselian discrete valuation ring, let $X$ be a separated
scheme of finite type over $R$, and put $Z=X_s$. Assume the universal
degree-zero comparison for $(X,Z)$. Then, for every torsion abelian
\'etale sheaf $\mathcal F$ on $X$ and every $q\ge0$,
\[
 H^q(X,\mathcal F)\xrightarrow{\sim}H^q(Z,i^{-1}\mathcal F).
\]
No flatness or properness hypothesis on $X$ is imposed.
\end{theorem}
\begin{proof}
Let $Y\subset X$ be the reduced closure of the generic fibre.
Every irreducible component of $Y$ dominates the trait, and $Y$ is
reduced. Multiplication by a uniformizer is consequently injective
on every local ring of $Y$: an element killed by the uniformizer
vanishes in every minimal-prime quotient and therefore is zero.
The same holds for every nonzero element of $R$, so $Y$ is
torsion-free, hence flat, over the DVR. It is separated and of
finite type over $R$.

The pair $(Y,Y_s)$ inherits the universal degree-zero property.
Indeed, if $a:Y\hookrightarrow X$ and $\mathcal H$ is an \'etale
sheaf of sets on $Y$, apply the hypothesis on $X$ to $a_*\mathcal H$.
Base change for the finite morphism $a$ identifies its restriction
to $Z$ with the pushforward of $\mathcal H|_{Y_s}$, giving precisely
the required comparison on $Y$. Lemma~\ref{C:henselian-flat-trait-proper}
therefore makes $Y$ proper over $R$.

Set $W=Y\times_X Z=Y_s$ and write $c:W\hookrightarrow X$.
The underlying topological space of $X$ is the union of $Y$ and $Z$,
because the base trait has only its generic and closed points.
For every abelian \'etale sheaf $\mathcal F$ there is a short exact
sequence
\[
 0\longrightarrow\mathcal F
 \longrightarrow a_*a^{-1}\mathcal F\oplus i_*i^{-1}\mathcal F
 \longrightarrow c_*c^{-1}\mathcal F\longrightarrow0,
\]
where the last map is the difference of restrictions. At a geometric
point in just one of $Y,Z$ this is the identity sequence for the
corresponding stalk; at a point of $W$ it is
$0\to F\to F\oplus F\to F\to0$. Hence the sequence is exact
without any reducedness assumption on $X$, $Z$, or $W$.

Closed direct image is exact and preserves injectives, so the long
exact cohomology sequence is
\[
 \cdots\longrightarrow H^{q-1}(Y,\mathcal F|_Y)
       \oplus H^{q-1}(Z,\mathcal F|_Z)
 \longrightarrow H^{q-1}(W,\mathcal F|_W)
 \longrightarrow H^q(X,\mathcal F)
\]
\[
 \longrightarrow H^q(Y,\mathcal F|_Y)\oplus H^q(Z,\mathcal F|_Z)
 \longrightarrow H^q(W,\mathcal F|_W)\longrightarrow\cdots.
\]
Proper base change over the henselian affine pair $(R,\mathfrak m_R)$,
as derived in Theorem~\ref{B:henselian-general}, makes
$H^r(Y,\mathcal F|_Y)\to H^r(W,\mathcal F|_W)$ an isomorphism for
every $r$. Thus the first displayed map onto $H^{q-1}(W,\mathcal F|_W)$
is surjective, and the kernel of the last displayed difference map
projects isomorphically onto $H^q(Z,\mathcal F|_Z)$. Exactness yields
the asserted isomorphism, including degree zero. If $Y$ is empty,
the same sequence simply says that the étale sheaf theory is
unchanged by the closed immersion $Z\hookrightarrow X$ with the
same underlying space.
\end{proof}

\begin{corollary}[Known reduction when the closed fibre is proper]\label{C:henselian-proper-closed-fibre}
Let $(A,I)$ be a henselian ring pair and let $X$ be separated and of
finite type over $A$. Put $Z=X\times_A A/I$. If $Z$ is proper over
$A/I$ and $(X,Z)$ satisfies the universal degree-zero comparison,
then $X$ is proper over $A$. If $A$ is noetherian, the all-degree
torsion comparison of Theorem~\ref{B:henselian-general} follows.
\end{corollary}
\begin{proof}
Stacks Lemma~37.53.9 (tag~0CT9), applied to the whole open-and-closed
subscheme $Z$ of the closed fibre, supplies an open-and-closed
$X'\subset X$ proper over $A$ with $X'\times_A A/I=Z$.
The open immersion $X'\hookrightarrow X$ has a section over $Z$.
The universal degree-zero hypothesis extends this section over $X$,
forcing $X'=X$. The final assertion is the stated proper theorem.
\end{proof}

These positive results impose explicit additional hypotheses absent
from subject~21. The arbitrary qcqs implication is refuted by
Theorem~\ref{C:subject21-torsion-counterexample} below; its nonseparated,
nonnoetherian construction lies outside the positive classes.

\begin{theorem}[Proved: properness over every positive-dimensional henselian noetherian local base]\label{B:henselian-local-flat-proper}
Let $(R,\mathfrak m,k)$ be a henselian noetherian local ring with
$\dim R>0$. Let $X$ be flat, separated and of finite type over $R$.
If every point closed in $X$ belongs to the closed fibre $X_k$,
then $X$ is proper over $R$. In particular, universal degree-zero
comparison for $(X,X_k)$ implies properness and, for every torsion
abelian \'etale sheaf $\mathcal F$ and every $q\geq0$,
\[
 H^q(X,\mathcal F)\xrightarrow{\sim}
 H^q(X_k,\mathcal F|_{X_k}).
\]
No excellence, Nagata, reducedness, irreducibility, completeness,
or dimension-one hypothesis is imposed on $R$.
\end{theorem}
\begin{proof}
Choose a prime $\mathfrak p\subsetneq\mathfrak m$ with
$\dim(R/\mathfrak p)=1$, by taking the penultimate member of a
saturated maximal chain of primes. Put $R_0=R/\mathfrak p$.
This is a henselian noetherian local domain: it is a finite local
$R$-algebra, so Stacks tag~04GH applies. Let $V$ be its integral
closure in its fraction field. Krull--Akizuki, tag~00PG, and its
normalization consequence, tag~09IG, show that $V$ is a
one-dimensional Dedekind domain, integral over $R_0$, with finite
residue field extensions. Finiteness of $V$ as an $R_0$-module
is not assumed or asserted.

Write $V=\bigcup_\alpha V_\alpha$ as the directed union of its
finite $R_0$-subalgebras. Each $V_\alpha$ is a domain finite over
a henselian local ring, hence is itself henselian local by
tag~04GH. The transition inclusions are integral, so their unique
maximal ideals contract to one another. Their union therefore
makes $V$ a local ring: an element outside that union is already
a unit at a finite stage. The ring $V$ is also henselian.
Indeed, a monic polynomial and a lift of a simple residue root
are defined at one finite stage; since the transition maps are
local, the derivative is already a unit there. The henselian
property of that stage supplies the root in $V$. Consequently
$V$ is a henselian DVR, by the noetherian normal local
dimension-one criterion, tag~00PD. Write $l$ for its residue
field; $l/k$ is finite. These are the exact primary normalization
and henselian statements used here \cite[Tags 00PG, 09IG, 04GH,
00PD]{BStacks}.

The base change $X_V$ is flat, separated and of finite type over
$V$. The projection $X_V\to X$ is integral, hence closed
\cite[Lemmas 29.45.6--7]{BStacks}. A closed point $y$ of $X_V$
therefore has image a closed point $x$ of $X$, which by hypothesis
lies over $\mathfrak m$. The only prime of $V$ over
$\mathfrak m$ is its closed point: primes over a maximal ideal
under an integral map are maximal, and $V$ is local. Thus all
closed points of $X_V$ lie in its closed fibre. The weaker
closed-point version of Lemma~\ref{C:henselian-flat-trait-proper}
makes $X_V$ proper over $V$.

It follows that $X_k\times_k l$ is proper over $l$. Properness
descends through the faithfully flat field extension $k\to l$
\cite[Tag 02L1]{BStacks}, so $X_k$ is proper. Stacks tag~0CT9,
applied to the entire proper closed fibre, supplies a clopen
subscheme $W\subset X$ proper over $R$ with $W_k=X_k$. If
$X\setminus W$ were nonempty, its noetherian underlying space
would have a closed point. Since this complement is closed in
$X$, that point would be closed in $X$ and outside $X_k$, a
contradiction. Hence $W=X$.

Universal degree-zero comparison implies the closed-point
hypothesis by testing the open subterminal sheaf represented by
$X\setminus\{x\}$ for a putative closed point outside $X_k$.
The all-degree torsion consequence is the proper henselian
comparison of Theorem~\ref{B:henselian-general}.
\end{proof}

The positive-dimensional and flat hypotheses in the properness
claim cannot be omitted. Over an artinian local $R$, the scheme
$\mathbb A^1_R$ satisfies universal degree-zero comparison with
its closed fibre, because that fibre is a nilpotent thickening
equivalence for the \'etale topology, but it is not proper. Over
a positive-dimensional $R$, the nonflat scheme $\mathbb A^1_k$
is supported on its closed fibre and is again nonproper. In
higher base dimension the reduced horizontal closure need not
be flat. The following independent curve argument treats the
nonflat case without making that inference.


% Curve-selection idea proposed by /root/cycles_b; local-curve argument
% supplied by /root/algebra_c; /root removed catenarity and locality using0CT9.
% Independently checked against primary sources by /root/algebra_c.
% This settles a restricted class, not arbitrary qcqs henselian pairs.
\begin{lemma}[Proved: a local curve avoiding a prescribed equation]\label{C:local-curve-avoidance}
Let $(A,\mathfrak n)$ be a noetherian local domain and let
$0\ne a\in\mathfrak n$. There is a prime $\mathfrak p\subset A$
such that $a\notin\mathfrak p$ and $\dim(A/\mathfrak p)=1$.
\end{lemma}
\begin{proof}
Choose a maximal ideal of the nonzero ring $A[1/a]$ and let
$\mathfrak p$ be its contraction to $A$. The local noetherian domain
$D=A/\mathfrak p$ has a nonzero nonunit $\bar a$, and
$D[1/\bar a]$ is a field. Thus every nonzero prime of $D$ contains
$\bar a$.

We verify directly that $\dim D=1$. If its maximal ideal had height
at least two, let $\mathfrak q_1,\ldots,\mathfrak q_r$ be the
finitely many primes minimal over $(\bar a)$. Each has height one,
by the principal ideal theorem and the domain hypothesis
\cite[Tag~00KV]{BStacks}. Prime avoidance
\cite[Tag~00DS]{BStacks} gives an element $b$ of the maximal ideal outside their
union. A prime $\mathfrak q$ minimal over $(b)$ also has height one.
It is nonzero, so it contains $\bar a$ and hence one of the
$\mathfrak q_i$. Equality follows from height one, contradicting
$b\notin\mathfrak q_i$. Therefore $\dim D\le1$. The nonzero
nonunit $\bar a$ shows that $D$ is not a field, so its dimension is
exactly one.
\end{proof}

\begin{theorem}[Proved: horizontal properness over every noetherian henselian ring pair]\label{C:henselian-horizontal-proper}
Let $(A,I)$ be a henselian ring pair with $A$ noetherian.
Let $X$ be an integral separated scheme of finite type over $A$,
whose generic point lies outside $X_I=X\times_A A/I$.
If every point closed in $X$ belongs to $X_I$, then $X$ is proper
over $A$. Locality of $A$, flatness, catenarity, excellence,
normality, and a finite normalization hypothesis are not required.
\end{theorem}
\begin{proof}
By Nagata compactification \cite[Theorem~38.33.8, Tag~0F41]{BStacks},
embed $X$ as an open subscheme of a proper $A$-scheme
$\overline X$. Taking the reduced closure of $X$, we may and do
assume that $\overline X$ is integral and that $X$ is dense in it.
Suppose the closed boundary $D=\overline X\setminus X$ is nonempty.
Its image in $\operatorname{Spec}A$ is a nonempty closed subset,
so contains a maximal ideal $\mathfrak m$. Since
$I\subseteq\operatorname{Jac}(A)$, this maximal ideal contains
$I$. Choose a closed point $z$ of the fibre of $D$ over
$\mathfrak m$. It is closed in $\overline X$ and lies in
$\overline X_I$.

Put $B=\mathcal O_{\overline X,z}$. This is a noetherian local
domain. The boundary ideal in $B$ contains a nonzero element $d$,
because the generic point of $\overline X$ lies in $X$. There is
also an element $r\in I$ with nonzero image in $B$:
otherwise the generic point of $X$ would lie in $X_I$.
The product $a=dr$ is a nonzero element
of the maximal ideal of $B$. Apply
Lemma~\ref{C:local-curve-avoidance} to obtain a prime
$\mathfrak q\subset B$ avoiding $a$ with $\dim(B/\mathfrak q)=1$.
Let $y$ be the corresponding generalization of $z$ in
$\overline X$. Since $d$ and $r$ are both absent from
$\mathfrak q$, we have $y\in X\setminus X_I$.

Let $C$ be the reduced closure of $y$ in $\overline X$.
Then $C$ is integral and proper over $A$, its generic point is $y$,
and
\[
 \mathcal O_{C,z}=B/\mathfrak q,
 \qquad \dim\mathcal O_{C,z}=1.
\]
The nonzero element $r$ remains nonzero in the one-dimensional
local domain $\mathcal O_{C,z}$. Hence
\[
 \mathcal O_{C_I,z}
 =\mathcal O_{C,z}/I\mathcal O_{C,z}
\]
is zero-dimensional: every prime of $\mathcal O_{C,z}$ containing
$r$ is its maximal ideal. Since $z$ is closed in the noetherian
scheme $C_I$, this also shows that $\{z\}$ is open in $C_I$.
Indeed no other point of $C_I$ generalizes to $z$, so $\{z\}$
is an irreducible component and is disjoint from the finitely
many other irreducible components. Let $E\subset C_I$ be the
induced open-and-closed subscheme supported at $z$. It is proper
over $A/I$, because it is a closed subscheme of the proper
$A/I$-scheme $C_I$.

We apply the lifting theorem for proper open-and-closed subschemes
\cite[Lemma~37.53.9, Tag~0CT9]{BStacks} to $C/A$ and the proper
open-and-closed subscheme $E\subset C_I$. It gives an
open-and-closed subscheme $W\subset C$ with $W_I=E$.
The integral scheme $C$ is connected, and $W$ is nonempty,
so $W=C$. Consequently $C_I=E$ is supported at the single
point $z$ and is zero-dimensional.

The quasi-finite locus of $C\to\operatorname{Spec}A$ is open
\cite[Lemma~29.57.2, Tag~01TI]{BStacks} and contains the whole closed
subscheme $C_I$ \cite[Lemma~29.30.5, Tag~0397]{BStacks}: its only
point $z$ is isolated in its fibre over $\mathfrak m$.
Its closed complement is proper over $A$. If that complement
were nonempty, its nonempty closed image in
$\operatorname{Spec}A$ would contain a maximal ideal and
therefore meet $V(I)$, contrary to its disjointness from $C_I$.
Thus $C\to\operatorname{Spec}A$ is proper and quasi-finite,
hence finite \cite[Lemma~37.44.1, Tag~02LS]{BStacks}.

Write $C=\operatorname{Spec}D_C$, where $D_C$ is a finite
$A$-algebra and a domain. Every maximal ideal of $D_C$ contracts
to a maximal ideal of $A$, so $ID_C\subseteq\operatorname{Jac}(D_C)$.
As $\operatorname{Spec}(D_C/ID_C)=C_I$ has only the point $z$,
the ring $D_C$ is local with unique closed point $z$. Since a local
scheme has dimension equal to the dimension of its local ring
at its closed point, $\dim C=\dim\mathcal O_{C,z}=1$.
Its only other point is consequently the generic point $y$.
Therefore
\[
 C\cap X=\{y\}
\]
is a nonempty closed subset of $X$ outside $X_I$. This contradicts
the closed-point hypothesis. The boundary was therefore empty,
which proves properness of $X$.
\end{proof}

\begin{theorem}[Proved: all-degree torsion comparison over every noetherian henselian ring pair]\label{C:henselian-noetherian-pair-comparison}
Let $(A,I)$ be a henselian ring pair with $A$ noetherian,
let $X$ be any separated scheme of finite
type over $A$, and set $Z=X\times_A A/I$. Suppose every closed point of $X$
lies in $Z$. Then for every torsion abelian \'etale sheaf
$\mathcal F$ on $X$ and every $q\ge0$,
\[
 H^q(X,\mathcal F)\xrightarrow{\sim}
 H^q(Z,i^{-1}\mathcal F).
\]
In particular this holds if $(X,Z)$ satisfies the universal
degree-zero comparison for \'etale sheaves of sets. Neither
flatness nor properness of $X$ is required. No locality,
catenarity, excellence, completeness, or dimension hypothesis is
imposed on $A$.
\end{theorem}
\begin{proof}
There are finitely many irreducible components of $X$.
Give each its reduced structure. Every component whose generic
point lies outside $Z$ satisfies the hypotheses of
Theorem~\ref{C:henselian-horizontal-proper}: its closed points are
closed in $X$ and therefore lie in $Z$. Such a component is proper
over $A$. Let $Y\subset X$ be their reduced union; allow $Y$ to be
empty. This finite union is proper over $A$. For completeness,
the disjoint union of the components maps finitely and
surjectively onto $Y$, and is proper over $A$. After arbitrary
base change, the inverse image of a closed subset of $Y$ has
closed image in the base; surjectivity identifies that image
with the image of the original closed subset. Thus $Y$ is
universally closed, and it is already separated and of finite
type, proving properness directly.

Every remaining irreducible component of $X$ is contained in $Z$,
since $Z$ contains its generic point. Hence
$|X|=|Y|\cup|Z|$. Set $W=Y\times_XZ=Y_I$, and denote the three
closed immersions into $X$ by $a:Y\hookrightarrow X$,
$i:Z\hookrightarrow X$ and $c:W\hookrightarrow X$. There is a
short exact sequence of abelian \'etale sheaves
\[
 0\longrightarrow\mathcal F
 \longrightarrow a_*a^{-1}\mathcal F\oplus i_*i^{-1}\mathcal F
 \longrightarrow c_*c^{-1}\mathcal F\longrightarrow0,
\]
with the difference of restrictions as its last map. Exactness
is checked at geometric points: it is the identity sequence on
the part belonging to just one closed subset, and
$0\to F\to F\oplus F\to F\to0$ on their intersection.
Thus no reducedness of $X$, $Z$ or $W$ is required.

Closed direct images are exact and preserve injectives.
Proper henselian torsion comparison for
$f:Y\to\operatorname{Spec}A$ identifies
$H^r(Y,\mathcal F|_Y)$ with $H^r(W,\mathcal F|_W)$ for every $r$.
Explicitly, proper base change identifies the restrictions of
the torsion sheaves $R^qf_*(\mathcal F|_Y)$ with the corresponding
closed-fibre direct images. Affine Gabber comparison for the
given henselian pair $(A,I)$ then identifies the two Leray
spectral sequences, as in the proper reduction in
Theorem~\ref{B:henselian-general}. This argument requires no
additional degree-zero hypothesis on $Y$.
In the cohomology sequence of the displayed closed-cover sequence,
the difference map from
$H^{q-1}(Y,\mathcal F|_Y)\oplus H^{q-1}(Z,\mathcal F|_Z)$
onto $H^{q-1}(W,\mathcal F|_W)$ is consequently surjective.
The kernel of the next difference map projects isomorphically
onto $H^q(Z,\mathcal F|_Z)$. Exactness proves the claimed
restriction isomorphism, including degree zero.

Finally, universal degree-zero comparison implies the stated
closed-point condition: if $x\in X\setminus Z$ were closed,
the open subterminal sheaf represented by $X\setminus\{x\}$
would have a section on $Z$ and no section on $X$.
\end{proof}

\begin{corollary}[Proved: the degree-zero criterion in the finite-type henselian-base class]\label{C:henselian-noetherian-closedpoint-criterion}
Let $(A,I)$ be a noetherian henselian ring pair, let $X$ be
separated and of finite type over $A$, and put $Z=X\times_A A/I$.
Then the following conditions are equivalent:
\begin{enumerate}
\item every point closed in $X$ belongs to $Z$;
\item restriction $H^0(X,\mathcal H)\to H^0(Z,i^{-1}\mathcal H)$
is bijective for every \'etale sheaf of sets $\mathcal H$.
\end{enumerate}
Either condition therefore gives the all-degree torsion comparison
of Theorem~\ref{C:henselian-noetherian-pair-comparison}.
\end{corollary}
\begin{proof}
The implication (2)$\Rightarrow$(1) was proved above using the
complement of a closed point. Suppose (1). Let $X'\to X$ be
finite, and put $Z'=X'\times_XZ$. A finite morphism sends closed
points to closed points, so every closed point of $X'$ belongs
to $Z'$. The scheme $X'$ is noetherian, separated and of finite
type over $A$. By the horizontal-component proof above,
$|X'|=|Y'|\cup|Z'|$, with $Y'$ proper over $A$.

Let $T\subset Z'$ be clopen. Its restriction to
$Y'_I=Y'\times_{X'}Z'$ is a proper clopen subscheme over $A/I$.
The proper-clopen lifting theorem
\cite[Lemma~37.53.9, Tag~0CT9]{BStacks} lifts it to a clopen subscheme
$D\subset Y'$. On the closed cover $Y'\cup Z'$, the subsets
$D$ and $T$ agree on the intersection. Their union and its
complement are both closed in $X'$: the complement is
$(Y'\setminus D)\cup(Z'\setminus T)$, and the agreement makes
the two displayed unions disjoint. Thus $D\cup T$ is a clopen
extension of $T$ to $X'$. It is unique, since any nonempty
clopen symmetric difference would contain a point closed in
the noetherian scheme $X'$, hence a point of $Z'$.

We have proved the finite-cover clopen criterion. Hall--Rydh's
Proposition~5.4 \cite{BHallRydh2023}, with the qcqs hypotheses satisfied here, gives
(2), as stated in Theorem~\ref{B:henselian-general}. No separate
proper-base-change assertion for arbitrary sheaves of sets
is used.
\end{proof}

The argument uses the isolated closed-fibre component of a
proper local-curve closure, and henselian lifting of that
component. It does not use a dimension formula or assert that
a horizontal closure is flat over a higher-dimensional base.
It establishes comparison for all separated finite-type models over
noetherian henselian ring pairs. The following construction refutes
the implication for
arbitrary qcqs pairs outside this class.


\subsubsection{A nonnoetherian torsion counterexample from Nagata's corrected example}

\begin{lemma}[Proved: the corrected Nagata input]\label{C:Nagata-input}
There are a normal local domain $(R,\mathfrak m)$, a nonzero
$d\in\mathfrak m$, and $a,b\in\operatorname{Frac}(R)$ such that
\[
 R\text{ is henselian},\qquad R[1/d]=\operatorname{Frac}(R),
 \qquad ab=d,
 \qquad a,a^{-1},b,b^{-1}\notin R.
\]
\end{lemma}
\begin{proof}
Here is a precise specialization of Nagata's construction. Let $K$ be
an algebraic closure of $\mathbb Q(t)$, equipped with an extension $w$
of the $t$-adic valuation, normalized by $w(t)=1$. Its value group is
$\mathbb Q$: the valuation of an algebraic element is rational, since
in an algebraic equation two terms must attain the least valuation,
and $K$ contains every root of $t$. Put $\alpha=\sqrt2$ and let $x$
be transcendental over $K$. For $\alpha\leq\lambda\leq2\alpha$,
define the Gauss valuation
\[
 v_\lambda\left(\sum_i c_i x^i\right)
   =\min_i\{w(c_i)+i\lambda\}.
\]
For every $e\in K$ with $\alpha<w(e)<2\alpha$, define the translated
Gauss valuation
\[
 v_e\left(\sum_i c_i(x+e)^i\right)
   =\min_i\{w(c_i)+2i\alpha\}.
\]
Extend these valuations to $K(x)$ by the quotient rule, and set
\[
 N=\bigcap_{\alpha\leq\lambda\leq2\alpha}\mathcal O_{v_\lambda}
   \ \cap\!
   \bigcap_{\alpha<w(e)<2\alpha}\mathcal O_{v_e}.
\]
Nagata's Theorem~1, with the replacement proof in his 1955
correction, proves that $N$ is completely integrally closed and
has exactly one nonzero prime ideal
\cite[Theorem~1, pp.~29--31]{CNagata1952}
\cite[pp.~209--212, replacing p.~30 line~14 through p.~31 line~7]{CNagata1955}.
The correction is essential here: the uncorrected 1952 proof is not
being used. More explicitly, its proof establishes that each nonzero
element $c\in N$ either has value zero for every displayed valuation,
or has all its values in an interval $[\epsilon_c,M_c]$ with
$0<\epsilon_c\leq M_c<\infty$. The former elements are units. For
two nonunits $c,h$, an integer $q$ with $q\epsilon_c\geq M_h$
gives $c^q/h\in N$. Thus every nonzero prime contains every nonunit,
which gives the asserted unique nonzero prime. Complete integral
closedness also follows directly from the rank-one valuations:
if $u z^q\in N$ for all $q\geq0$, with $u\ne0$, then each valuation
of $z$ is nonnegative.

Choose a fourth root $d=t^{1/4}\in K$, and put
\[
 a=\frac{t^2}{x},\qquad b=t^{-7/4}x.
\]
Every displayed valuation of $d$ is $1/4$, so $d$ is a nonzero
nonunit of $N$, and $ab=d$. The endpoint values are
\[
\begin{array}{c|cc}
 &v_\alpha&v_{2\alpha}\\ \hline
 a&2-\sqrt2>0&2-2\sqrt2<0\\
 b&\sqrt2-7/4<0&2\sqrt2-7/4>0.
\end{array}
\]
Consequently $a,a^{-1},b,b^{-1}$ all lie outside $N$.
Moreover $\operatorname{Frac}(N)=K(x)$: $x\in N$, and for every
$c\in K$ some $d^q c$ has nonnegative $w$-value and hence belongs
to $N$. Since $N$ has only its zero and maximal prime, inverting
$d$ gives this fraction field.

Let $R=N^h$. Normality of henselization for arbitrary local rings
is Stacks Project, tag~06DI \cite{BStacks}. To check the Goldman property without
any noetherian assumption, express $R$ as the filtered colimit of
local pointed \'etale neighbourhoods of $N$. Each stage is a normal
local domain, and after inverting $d$ it is a domain essentially
\'etale over $K(x)$, hence a finite separable field. The local flat
transition maps are injective, so $R[1/d]$ is a field and is therefore
$\operatorname{Frac}(R)$. Finally, faithful flatness preserves all
four exclusions. Indeed, if a fraction $u/v$ of $N$ belongs to $R$,
then $u\in vR\cap N=vN$, so $u/v\in N$. This proves the lemma.
\end{proof}

\begin{lemma}[Proved: a normal Goldman local ring with two surviving branches]\label{C:Goldman-branches}
For $R,d,a,b$ as in Lemma~\ref{C:Nagata-input}, there is a normal
local domain $(B,\mathfrak n)$ containing $R[a,b]$ such that
\[
 B[1/d]=\operatorname{Frac}(B),\qquad d\in\mathfrak n,
\]
and two primes $\mathfrak p,\mathfrak q\subsetneq\mathfrak n$ satisfy
\[
 b\in\mathfrak p,\quad a\notin\mathfrak p,
 \qquad a\in\mathfrak q,\quad b\notin\mathfrak q.
\]
In particular $V(d)\cap D(a+b)$ is nonempty and disconnected in
$\operatorname{Spec}B$.
\end{lemma}
\begin{proof}
Write $L=\operatorname{Frac}(R)$ and $k=R/\mathfrak m$. Let $C$ be
the integral scheme-theoretic closure of the $L$-point $(a,b)$ in
$\mathbf P^1_R\times_R\mathbf P^1_R$. It is proper and of finite
type over $R$, even though $R$ need not be noetherian. In homogeneous
coordinates $([A_0:A_1],[B_0:B_1])$ it lies in
\[
 H:\quad A_0B_0=d A_1B_1.
\]
The special fibre $H_k$ is the reduced union of two projective lines
$L_a=\{B_0=0\}$ and $L_b=\{A_0=0\}$, meeting at the single point
$o=([0:1],[0:1])$.

First consider the integral graph closure $C_a\subset\mathbf P^1_R$
of the fraction $a$. Its special fibre must have support all of
$\mathbf P^1_k$. Otherwise that fibre is a finite set. The
quasi-finite locus of $C_a\to\operatorname{Spec}R$ then contains
the special fibre. Its closed complement is proper over the local
base and misses the closed point, so it is empty. Thus $C_a\to
\operatorname{Spec}R$ is proper and quasi-finite, hence finite
(Stacks Project, tags~01TI and~02LS). It is birational, and $R$ is
normal, so it is an isomorphism (tag~0AB1). The resulting section of
$\mathbf P^1_R$ has its closed point in one of the two standard
affine charts. Since an open subset of a local spectrum containing
the closed point is the whole spectrum, this would put either
$a$ or $a^{-1}$ in $R$, a contradiction. The identical argument
applies to $C_b$.

The two projections of $C$ have images $C_a,C_b$: each projection
is proper and its image contains the generic point of the respective
integral graph closure. Thus on the special fibre both projections
are surjective onto $\mathbf P^1_k$. A closed subset of $H_k$ with
this property contains both $L_a$ and $L_b$. Consequently $C_k$
has support $H_k$, and in fact $C_k=H_k$ as schemes, since $C_k$
is a closed subscheme of the reduced scheme $H_k$ with full support.

Let $\nu:\widetilde C\to C$ be the normalization. Its existence
and integrality require no finiteness assertion: $C$ is integral
(Stacks Project, Section~29.55). By tag~0817(1), the integral
quasi-coherent algebra $\nu_*\mathcal O_{\widetilde C}$ is the
directed union of its finite quasi-coherent
$\mathcal O_C$-subalgebras $\mathcal A_\gamma$. Put
$C_\gamma=\operatorname{Spec}_C\mathcal A_\gamma$. Every
$C_\gamma$ is integral, finite and birational over $C$, and hence
proper and of finite type over $R$. In particular no finite
normalization hypothesis is being imposed. The transition maps
$C_\delta\to C_\gamma$ are finite and the inverse limit is
$\widetilde C$.

The special fibre $(C_\gamma)_k$ is connected. Indeed, any
nonempty proper clopen part would be proper over $k$ and would
lift, by Stacks Project tag~0CT9 for the henselian pair
$(R,\mathfrak m)$, to a nonempty proper clopen subset of the
integral scheme $C_\gamma$. This is impossible.
The fibre $(C_\gamma)_k$ is finite over the nodal curve $H_k$,
and is therefore noetherian of dimension one. A zero-dimensional
irreducible component would be a closed singleton disjoint from
the finitely many other irreducible components; it would be clopen.
There are curve components above both $L_a$ and $L_b$, by
surjectivity of the finite morphism, so such a singleton cannot
occur. Each irreducible component therefore maps onto exactly one
of $L_a,L_b$.

Let $E_\gamma$ be the finite set of triples $(p_\gamma,q_\gamma,
z_\gamma)$ where $p_\gamma,q_\gamma$ are generic points of
components mapping onto $L_a,L_b$, respectively, and
$z_\gamma$ specializes from both and lies above $o$.
This set is nonempty. Otherwise the union of components mapping
onto $L_a$ would be disjoint from the union mapping onto $L_b$,
giving a disconnection. The finiteness follows because there are
finitely many components and the fibre above $o$ is finite.
Finite transition morphisms carry these triples to triples at
earlier stages: generic points over the two axis generic points
remain such generic points, and specialization is preserved.
The inverse limit of this directed system of finite nonempty sets
is nonempty. For completeness, in the product of their finite
discrete compact spaces every finite collection of compatibility
conditions is satisfied by choosing a triple at a common later
stage and mapping it down. Compactness then satisfies all the
conditions. Surjectivity of the transition maps on triples is not
required.

Choose a compatible family of triples. All its points lie over
the affine chart $C\cap\mathbf A^2_R=\operatorname{Spec}R[a,b]$.
Let $T$ be the integral closure of $R[a,b]$ in $L$. Restricting the
finite subalgebras to this chart gives a directed union $T$; the
compatible prime ideals give primes $\mathfrak p_\infty,
\mathfrak q_\infty\subset\mathfrak n_\infty$ of $T$.
Their contractions to $R[a,b]$ are the generic points of the two
affine axes and their node, respectively. In particular
\[
 b\in\mathfrak p_\infty,\quad a\notin\mathfrak p_\infty,
 \quad a\in\mathfrak q_\infty,\quad b\notin\mathfrak q_\infty,
 \quad a,b\in\mathfrak n_\infty.
\]
The contractions ensure the same strict containments after setting
$B=T_{\mathfrak n_\infty}$. This is a normal local domain. It is
an overring of $R$ inside $L$, so $B[1/d]=L$ since $R[1/d]=L$.
The element $d=ab$ is nonzero and belongs to the displayed primes.

Finally put $h=a+b$. Both $\mathfrak p$ and $\mathfrak q$ avoid
$h$, while $h\in\mathfrak n$. In $B_h/(d)$ the element
\[
 e=\frac{a}{a+b}
 \quad\text{satisfies}\quad
 e^2-e=-\frac{ab}{(a+b)^2}=0.
\]
It has value $1$ at $\mathfrak p$ and $0$ at $\mathfrak q$.
It is therefore a nontrivial idempotent, proving nonemptiness and
disconnectedness of $V(d)\cap D(h)$.
\end{proof}

\begin{theorem}[Refuted: unrestricted torsion comparison]\label{C:subject21-torsion-counterexample}
There exist a quasi-compact quasi-separated scheme $X$ and a
principal closed subscheme $i:Z\hookrightarrow X$ for which
\[
 H^0(X,\mathcal F)\xrightarrow{\ \sim\ }
 H^0(Z,i^{-1}\mathcal F)
 \quad\text{for every \'etale sheaf of sets }\mathcal F,
\]
but, for every integer $n>1$,
\[
 H^1(X,\underline{\mathbb Z/n\mathbb Z})=0,
 \qquad H^1(Z,\underline{\mathbb Z/n\mathbb Z})\ne0.
\]
In particular, degree-zero comparison for all \'etale sheaves of
sets does not imply higher torsion comparison for arbitrary
scheme pairs.
\end{theorem}
\begin{proof}
Let $B,d,a,b$ be given by Lemma~\ref{C:Goldman-branches}, and let
$A=B^{\mathrm{sh}}$. This is a normal local domain by tag~06DI.
The local \'etale-stage argument in Lemma~\ref{C:Nagata-input}
shows $A[1/d]=\operatorname{Frac}(A)$. Faithful flatness of
$B\to A$ preserves the nonemptiness of both clopen parts of
$\operatorname{Spec}(B_h/(d))$, where $h=a+b$. Hence
\[
 T_0=\operatorname{Spec}(A_h/(d))
\]
is nonempty and disconnected.

Glue two copies of $S=\operatorname{Spec}A$ by the identity on
$U=D(h)$, obtaining $X$. Let $Z=V(d)\subset X$. The intersection
of the two affine charts is principal affine, so $X$ is
quasi-compact and quasi-separated. Its open complement of $Z$
is the single generic point $\eta=\operatorname{Spec}A[1/d]$.

We verify the degree-zero assertion using the finite-clopen
criterion of Hall--Rydh, Proposition~5.4 \cite{BHallRydh2023}.
Let $Y\to X$ be any
finite morphism and $E\subset Y_Z$ any clopen subset. On each
chart $Y_i=Y\times_XS_i$, the henselian ring pair $(A,dA)$ gives
a unique clopen extension $E_i$ of $E|_{(Y_i)_Z}$. Their symmetric
difference on $Y_U$ is a clopen subset $D$ missing $(Y_U)_Z$.
If $D$ is nonempty, take $y\in D$. Its image in $U$ is $\eta$.
The reduced closure $W$ of $y$ in the affine finite chart $Y_1$
is finite and integral over $S$. It dominates $S$ and is
therefore surjective. The nonempty open $W_U$ is irreducible,
with generic point $y$, and its inclusion in the closed subset
$D\subset Y_U$ follows by specialization. But $W_U$ meets
$(Y_U)_Z$, since $W\to S$ is surjective and $T_0$ is nonempty.
This is a contradiction. Thus $E_1,E_2$ agree and glue uniquely.
Every finite $Y\to X$ therefore has the same clopen subsets as
$Y_Z$, proving the asserted comparison for all \'etale sheaves
of sets.

Set $\Lambda=\mathbb Z/n\mathbb Z$. The schemes $S$ and
$\operatorname{Spec}(A/(d))$ are strictly henselian local, so
their constant finite-group torsors are trivial; their $H^1$
groups with coefficients $\underline\Lambda$ vanish. Their
$H^0$ groups equal $\Lambda$. The scheme $U$ is nonempty and
integral, so $H^0(U,\underline\Lambda)=\Lambda$. The two
Mayer--Vietoris sequences therefore give
\[
\begin{split}
 H^1(X,\underline\Lambda)&=
  \operatorname{coker}(\Lambda\oplus\Lambda
    \xrightarrow{(r,s)\mapsto r-s}\Lambda)=0,\\
 H^1(Z,\underline\Lambda)&=
  H^0(T_0,\underline\Lambda)/\underline\Lambda\ne0.
\end{split}
\]
The last nonzero class can be represented by the locally constant
function equal to $1$ on one nonempty clopen part defined by the
idempotent $e=a/(a+b)$ and $0$ on its complement. It cannot be
the difference of two constant sections on the two local charts.
For $j:\eta\hookrightarrow X$, the localization sequence also
injects this nonzero group into
$H^2(X,j_!\underline\Lambda)$, directly realizing the previously
isolated degree-two obstruction.
\end{proof}

\paragraph{Scope and checks.}
The counterexample is nonseparated: its two chart embeddings
$S\to X$ agree on the dense proper open $U$, but send the closed
point to different points. A separated target would force these
embeddings to agree. It is also necessarily nonnoetherian.
The ring $A$ has a nonzero prime strictly below its maximal ideal
(a lift of either branch prime, avoiding $h$), so it has dimension
at least two; a noetherian local Goldman domain has dimension at
most one, as proved in Lemma~\ref{C:local-curve-avoidance}.
The construction does not contradict affine Gabber rigidity,
proper-pair rigidity, or the separated finite-type theorem over
noetherian henselian bases. No localization of a henselian pair
has been assumed to remain henselian, and no normalization has
been assumed finite.


\begin{proposition}[Known, with derivation: an affine closed fibre forces affineness]\label{C:henselian-affine-closed-fibre}
Let $X$ be a quasi-separated scheme and let $i:Z\hookrightarrow X$
be an affine closed subscheme. Suppose that
\[
 H^0(X,\mathcal F)\longrightarrow H^0(Z,i^{-1}\mathcal F)
\]
is bijective for every \'etale sheaf of sets on $X$. Then $X$ is
affine. Consequently restriction induces an isomorphism on $H^q$
for every $q\ge0$ and every torsion abelian \'etale sheaf on $X$.
No quasi-compactness, noetherianity, or finite-presentation assumption
on the original pair is required beyond the stated quasi-separatedness
of $X$ and affineness of $Z$.
\end{proposition}
\begin{proof}
Apply Alper--Hall--Halpern-Leistner--Rydh's affine \'etale
neighbourhood theorem \cite[Theorem~1.2 and its proof in Section~3]{C_AHHR2024}
to the closed immersion $Z\hookrightarrow X$ and the identity
$Z\to Z$. It gives an affine scheme $W$ and an \'etale morphism
$f:W\to X$ with $W\times_X Z\simeq Z$ over $Z$.
The inverse image of the \'etale sheaf represented by $W\to X$ is
the sheaf represented by $W\times_X Z\to Z$. The latter has its
identity section. Universal degree-zero comparison therefore supplies
a section $s:X\to W$ of $f$.

The morphism $f$ is separated: its relative diagonal is the base
change of the closed absolute diagonal of the affine scheme $W$
along $W\times_XW\to W\times_{\operatorname{Spec}\mathbb Z}W$.
A section of a separated morphism is a closed immersion. Hence
$s$ realizes $X$ as a closed subscheme of the affine scheme $W$,
so $X$ is affine. Write $X=\operatorname{Spec}A$ and
$Z=\operatorname{Spec}(A/I)$. The degree-zero hypothesis makes
$(A,I)$ a henselian ring pair, by Theorem~\ref{B:henselian-general}.
Gabber's affine theorem, Theorem~\ref{B:henselian}, proves the
torsion comparison in every degree.
\end{proof}

The neighbourhood theorem is applied in its scheme form, with
trivial stabilizers. Its statement does not impose the characteristic
or stabilizer conditions used in the later theorems on linearly
fundamental stacks. Its proof first replaces the ambient space by a
quasi-compact neighbourhood of the affine closed subscheme, and
then uses approximation; this does not add an unmentioned
quasi-compactness assumption to the proposition.

\begin{lemma}[Proved: affine henselizations give a cohomological cover]\label{B:henselized-affine-cover}
Let $X$ be qcqs, let $i:Z\hookrightarrow X$ be closed, and assume
universal degree-zero comparison for every \'etale sheaf of sets.
Choose finitely many affine open subschemes
$U_a=\operatorname{Spec}A_a$ whose union contains $Z$, and write
$Z\cap U_a=\operatorname{Spec}(A_a/I_a)$. Then
\[
 W=\coprod_a\operatorname{Spec}(A_a)^h_{I_a}\longrightarrow X
\]
is an fpqc cover and a $v$-cover, with each displayed affine term
a henselian pair along its inverse image of $Z$. Torsion \'etale
cohomology satisfies descent along this cover. This statement
does not assert that its \v Cech intersections are henselian pairs.
\end{lemma}
\begin{proof}
Every nonempty closed subset $T\subset X$ meets $Z$: otherwise
the open subterminal sheaf represented by $X\setminus T$ would
have a section on $Z$ but no section on $X$. In particular every
point $x$ specializes to some $z\in Z$, and the union of the
$U_a$ is all of $X$.

Henselization is ind-\'etale, flat, and an isomorphism modulo the
specified ideal \cite[Tag 0AGU]{BStacks}. Choose $a$ with
$z\in U_a$. Since open subsets are stable under generalization,
$x\in U_a$. The henselization has a point over $z$, and flatness
lifts its generalization to $x$ \cite[Tag 03HV]{BStacks}.
Therefore $W\to X$ is surjective. Each henselization is affine
over $U_a$, and each affine open $U_a\to X$ is quasi-compact
because $X$ is quasi-separated. The finite disjoint union is
thus flat and quasi-compact, proving that it is an fpqc cover.

Here is also the required valuative verification. Base change
along $\operatorname{Spec}V\to X$, where $V$ is a valuation
ring with fraction field $K$. Choose a point $y$ of $W_V$ over
the closed point of $V$. Flatness lifts the generic
generalization to a point $\eta'$ specializing to $y$. Let $T$
be the reduced closure of $\eta'$ in $W_V$. Its local ring at
$y$ is a local domain with fraction field $\kappa(\eta')$ and
contains $V$ injectively. A valuation ring $V'$ of this fraction
field dominating $\mathcal O_{T,y}$ exists by
\cite[Tag 00IA]{BStacks}. The induced extension $V\to V'$ is
local and torsion-free, hence faithfully flat. The map
$\operatorname{Spec}V'\to T\to W_V$ is the required lift in
the definition of a $v$-cover.

Bhatt--Mathew \cite[Lemma 5.1 and Proposition 5.2]{BBhattMathew2021}
give descent for torsion \'etale cohomology along $v$-covers;
their Theorem~5.4 proves the stronger arc-descent statement.
Their affine-base formulation applies after restriction to
affine opens of $X$, with the restricted coefficient sheaf;
ordinary \'etale descent glues these equivalences over $X$.
The complete cited proofs were opened.
\end{proof}

\begin{proposition}[Refuted: automatic henselianity of the intersections in that cover]\label{B:henselized-Cech-obstruction}
Let $k$ be a characteristic-zero field and put
\[
 A=k[[x,t]],\qquad X=\operatorname{Spec}A,\qquad Z=V(t),
 \qquad B=(A[1/x])^h_{(t)}.
\]
The affine henselian pair $(X,Z)$ has universal degree-zero
comparison. Nevertheless, with $C=B\otimes_A B$, the map
\[
 H^0(\operatorname{Spec}C,\{0,1\})\longrightarrow
 H^0(\operatorname{Spec}(C/tC),\{0,1\})
\]
is not injective. Thus a \v Cech self-intersection in the
henselization cover associated with $\{X,D(x)\}$ fails even the
degree-zero hypothesis. Henselizing that intersection along $t$
does not produce an arc-cover of the original intersection.
This refutes an automatic termwise descent argument, not torsion
comparison for general henselian scheme pairs.
\end{proposition}
\begin{proof}
The ring $A$ is $t$-adically complete. The quotient
$A[1/x]/tA[1/x]=k((x))$ is a field, so $B$ is local. Its
universal property first inverts every element outside $(t)$,
and identifies it with the ordinary henselization of the DVR
$A_{(t)}$. Hence $B$ is a henselian DVR and its fraction field
is algebraic separable over $K=\operatorname{Frac}A$.
Hensel's lemma gives a unique $u\in B$ such that
\[
 u^2-u=t/x,\qquad u\equiv0\pmod t.
\]
Set $u_1=u\otimes1$, $u_2=1\otimes u$ and $d=2u_1-1$ in $C$.
The element $d$ is a unit, and subtraction of the two quadratic
equations gives
\[
 (u_1-u_2)^2=(2u_1-1)(u_1-u_2).
\]
Consequently
\[
 e=(u_1-u_2)/d
\]
is an idempotent. Both $u_i$ vanish modulo $t$, so $e$ maps to
zero in $C/tC$.

It remains to prove $e\ne0$. The discriminant
\[
 1+4t/x=(x+4t)/x
\]
is not a square in $K$. Indeed, changing coordinates from
$(x,t)$ to $(x+4t,t)$ identifies the localization of $A$ at
$(x+4t)$ with a DVR in which this fraction has valuation one.
Thus $K(u)/K$ is a separable quadratic extension with
nontrivial automorphism $u\mapsto1-u$. Embed
$\operatorname{Frac}B$ into an algebraic closure of $K$ and
extend this automorphism to that closure. The resulting two
$K$-embeddings of $B$ give a homomorphism $C\to\overline K$
under which $u_1\mapsto u$, $u_2\mapsto1-u$, and $e\mapsto1$.
Therefore $e\ne0$. The multiplication map $C\to B$ sends $e$
to zero, so $e\ne1$ as well.

The nonempty clopen component $\operatorname{Spec}(eC)$ is
disjoint from $V(t)$. The maps to the constant two-point scheme
specified by $0$ and $e$ are different on $\operatorname{Spec}C$
but become equal on its closed fibre, proving the failure of
injectivity. The cover $X\amalg\operatorname{Spec}B\to X$
is exactly the cover of Lemma~\ref{B:henselized-affine-cover}
for the legitimate, though redundant, affine cover
$\{X,D(x)\}$. Its first \v Cech term has
$\operatorname{Spec}C$ as a clopen summand.

Finally, the image of $e$ in $C^h_{(t)}$ is zero: in a henselian
pair an idempotent reducing to zero is zero. Thus
$\operatorname{Spec}C^h_{(t)}\to\operatorname{Spec}C$ misses
the entire component $\operatorname{Spec}(eC)$. A map from a
field at a point of that component cannot lift after any field
extension, so this is neither a $v$-cover nor an arc-cover.
Replacing the matching term by its henselization therefore
changes the descent diagram without a covering justification.
\end{proof}

The program \texttt{code/B\_henselized\_cech.py} checks seven
exact polynomial identities and specializations, with output in
\texttt{logs/B\_henselized\_cech.json}; an independent rerun
passed. These finite identities do not replace the valuative
and cohomological proofs. Theorem~\ref{C:subject21-torsion-counterexample}
shows that no descent argument can prove comparison for all such
pairs. Any successful positive comparison requires additional
hypotheses, such as those in the proved affine, proper, or separated
finite-type noetherian henselian-base theorems.

\begin{proposition}[Refuted even for constant integer coefficients; explicit published-type counterexample]\label{B:constant-integer-rigidity}
Set
\[
 f=y^2-x^2(x+1),\qquad
 A=\varprojlim_{n\geq1}\mathbb C[x,y]/(f^n),\qquad
 X=\operatorname{Spec}A,\quad Z=V(f)\subset X.
\]
Then $(A,fA)$ is a noetherian henselian pair, $X$ is regular and
integral, and restriction on $H^0$ is bijective for every \'etale
sheaf of sets. Nevertheless, for the constant abelian sheaf
$\underline{\mathbb Z}$,
\[
 H^1_{\acute et}(X,\underline{\mathbb Z})=0,
 \qquad
 H^1_{\acute et}(Z,\underline{\mathbb Z})\simeq\mathbb Z.
\]
Thus the unrestricted-coefficient implication in subject~21 is
false even for a constant, finitely generated, torsion-free
coefficient group. For this very same pair, the comparison remains
an isomorphism in all degrees for every torsion abelian \'etale
sheaf, by Theorem~\ref{B:henselian}.
\end{proposition}
\begin{proof}
This is a fully specified version of the nodal-cubic construction
in Alper, Hall, Halpern-Leistner, and Rydh
\cite[Example~3.1]{C_AHHR2024}. We give the complete cohomology
calculation, rather than relying on the source's description of an
infinite cyclic torsor.

The ring $\mathbb C[x,y]$ is noetherian, regular, and excellent.
Its $f$-adic completion $A$ is noetherian and $f$-adically complete,
and $A/fA=\mathbb C[x,y]/(f)$. The completion map is regular by
Stacks Lemma~15.51.14, tag~0AH2, since a finite-type algebra over a
field is a G-ring. Regularity of $A$ follows by applying Stacks
Lemma~10.112.8, tag~031E, to its local rings: the base local rings
and the fibre local rings are regular and the map is flat.
Completeness makes $(A,fA)$ henselian, by Stacks
Lemma~15.11.4, tag~0ALJ. Theorem~\ref{B:henselian} gives the
universal $H^0$ comparison.

The polynomial $f$ is irreducible over $\mathbb C(x)$ as a
quadratic in $y$: $x^2(x+1)$ has odd valuation at $x=-1$ and is
not a square. Gauss's lemma makes $\mathbb C[x,y]/(f)$ a domain.
Hence $Z$ is connected. Idempotents lift uniquely in a henselian
pair, so $A$ has no nontrivial idempotent. A noetherian regular
ring is normal, and its finitely many irreducible components are
pairwise disjoint open-and-closed subsets: at a meeting point its
local ring would not be a domain. Connectedness therefore makes
$X$ integral.

We use the following elementary fact with $S=X$ and later with
$S=\mathbb A^1_{\mathbb C}$: for an integral normal noetherian
scheme $S$ with function field $K$, every
$\underline{\mathbb Z}$-torsor is trivial. First, its restriction
to $\operatorname{Spec}K$ is trivial. Indeed the equivalence with
discrete Galois modules (Stacks Lemmas~59.59.1--59.59.2) identifies
its class with a continuous homomorphism $G_K\to\mathbb Z$;
compactness gives a finite image, and $\mathbb Z$ has no nonzero
finite subgroup. Choose such a generic section.

Take an \'etale cover trivializing the torsor, and refine it by
connected noetherian normal schemes $V$. Each $V$ is integral,
and its generic fibre over $S$ is the spectrum of a finite
separable extension of $K$. Relative to the chosen trivialization,
the generic section is therefore one integer. That integer
extends uniquely to the constant section on $V$. On an overlap,
the extensions agree: each nonempty connected component of the
overlap is integral and \'etale over $S$, hence dominates a
nonempty open subset of $S$, and the two locally constant sections
agree at its generic point. They consequently agree everywhere
on that component. The sheaf property glues the sections over
$S$. This proves the stated fact and gives
$H^1_{\acute et}(X,\underline{\mathbb Z})=0$.

The normalization of $Z$ is the finite map
\[
 \nu:\mathbb A^1_{\mathbb C}\longrightarrow Z,
 \qquad t\longmapsto (t^2-1,t(t^2-1)).
\]
The equation $t^2=x+1$ proves finiteness; the equality $t=y/x$
away from $x=0$ proves birationality; and $\mathbb C[t]$ is
normal. The sole singular point is $o=(0,0)$, whose preimages
are $t=1$ and $t=-1$. The degree-two tangent cone
$y^2-x^2=(y-x)(y+x)$ has two distinct lines, so this is an
ordinary node. The normalization is an isomorphism away from
$o$.

Let $a:\{o\}\hookrightarrow Z$. Difference of values at $t=1$
and $t=-1$ defines a short exact sequence of \'etale abelian
sheaves
\[
 0\longrightarrow\underline{\mathbb Z}_Z
 \longrightarrow\nu_*\underline{\mathbb Z}_{\mathbb A^1}
 \longrightarrow a_*\underline{\mathbb Z}_{\{o\}}
 \longrightarrow0.
\]
Off $o$ the first map is an isomorphism. At a geometric point
over $o$, finite base change identifies the sequence with
$0\to\mathbb Z\to\mathbb Z^2\to\mathbb Z\to0$, with the
diagonal and difference maps. Thus exactness holds on every
geometric stalk. Finite direct image is exact and computes
cohomology (Stacks Proposition~59.55.2). The preceding normal
scheme argument gives $H^1(\mathbb A^1,\underline{\mathbb Z})=0$.
The global difference map from
$H^0(\mathbb A^1,\underline{\mathbb Z})=\mathbb Z$ to
$H^0(\{o\},\underline{\mathbb Z})=\mathbb Z$ is zero.
The long exact sequence therefore gives the claimed isomorphism
$H^1(Z,\underline{\mathbb Z})\simeq\mathbb Z$.

Inverse image of a constant sheaf is constant, so these are exactly
the two groups in the restriction comparison for
$\mathcal F=\underline{\mathbb Z}_X$. The calculation exhibits
the failure in its specified sheaf category. The normalization,
singular-point, fibre, and tangent-cone identities are independently
checked by \texttt{code/ROOT\_nodal\_constant\_checks.py}; the
cohomological argument is a proof, not a finite computation.
\end{proof}

\begin{proposition}[Refuted: degree zero for one coefficient sheaf does not suffice]\label{B:H0-counterexample}
Even for an affine smooth complex scheme, an isomorphism in degree zero for a fixed torsion \'etale sheaf does not imply isomorphisms in higher degrees. Take $X=\operatorname{Spec}\mathbb C[t,t^{-1}]$, $Z=\{t=1\}$ and $\mathcal F=\mu_n$ for any integer $n>1$.
\end{proposition}
\begin{proof}
Both schemes are connected and all $n$th roots of unity lie in $\mathbb C$, so restriction on $H^0$ is the identity of $\mu_n(\mathbb C)$. The Kummer sequence and $\operatorname{Pic}(\mathbb C[t,t^{-1}])=0$ give
\[
 H^1_{\acute et}(X,\mu_n)
 =\mathbb C[t,t^{-1}]^*/\mathbb C[t,t^{-1}]^{*n}
 =\mathbb Z/n\mathbb Z,
\]
where the generator is the unit $t$, or equivalently the finite \'etale cover $s^n=t$. Its derivative $ns^{n-1}$ is invertible, so it is indeed \'etale. On $Z$, Kummer gives $H^1(Z,\mu_n)=0$, since $\mathbb C$ is algebraically closed. Thus the degree-one comparison has nonzero kernel. For $n=2$, the residual exponent calculation modulo $2$ is included in the exact finite checks; the cohomological identification is proved here.
\end{proof}

\begin{proposition}[Refuted: unrestricted coefficients even under universal degree zero]\label{B:nontorsion-rigidity}
Let $R=\mathbb C[[t]]$ and let $B=(R[x])^h_{(t)}$ be the henselization along the ideal $(t)$. Set $X=\operatorname{Spec}B$ and $i:Z=\operatorname{Spec}(B/tB)\hookrightarrow X$, so $Z\simeq\mathbb A^1_{\mathbb C}$. Restriction is bijective on $H^0$ for every \'etale sheaf of sets on $X$. Nevertheless, for the additive \'etale structure sheaf $\mathcal F=\mathcal O_X$,
\[
 H^1_{\acute et}(X,\mathcal F)=0,
 \qquad H^1_{\acute et}(Z,i^*\mathcal F)\ne0.
\]
Here $i^*\mathcal F$ means inverse image as an abelian \'etale sheaf, not the tensor-product structure sheaf $\mathcal O_Z$. Thus the literally unrestricted-sheaf reading of subject~21 is false, even for an affine henselian pair satisfying its strongest degree-zero hypothesis. The torsion theorem remains the valid replacement.
\end{proposition}
\begin{proof}
The pair $(B,tB)$ is henselian by construction, so Theorem~\ref{B:henselian} gives the asserted universal $H^0$ comparison. The affine vanishing $H^1_{\acute et}(X,\mathcal O_X)=0$ is the quasi-coherent comparison theorem, Stacks Theorem~59.22.4, tag~03P2.

For $g=x,x-1,x(x-1)$, put
\[
 E_g=(R[x,1/g])^h_{(t)},\qquad U_g=D(g)\subset Z.
\]
The map $\operatorname{Spec}E_g\to X$ is ind-\'etale, and its closed fibre is $U_g$. Inverse image of the small \'etale structure sheaf under an ind-\'etale map is the structure sheaf: this follows by writing the map as a filtered limit of \'etale affine maps and using finite presentation of \'etale neighborhoods, or by the corresponding strict-local-ring stalks. Applying the degree-zero henselian comparison to $(E_g,tE_g)$ therefore gives, compatibly with restriction,
\[
 \Gamma(U_g,i^*\mathcal O_X)=E_g.
\]
The two opens $U_x,U_{x-1}$ cover $Z$. Their Mayer--Vietoris sequence embeds
\[
 E_{x(x-1)}/(E_x+E_{x-1})
 \hookrightarrow H^1_{\acute et}(Z,i^*\mathcal O_X).
\]

We use the explicit class of Gongyo--Kato--Uchiba \cite[Definition 1.3 and Theorem 3.2, pp.~2, 5--7]{BGongyoKatoUchiba2026}. In $E_{x(x-1)}$, Hensel's lemma gives the unique root
\[
 w^2+w=\frac{t}{x(x-1)},\quad w\equiv0\pmod t,
 \qquad f=1+2w,\quad f^2=1+\frac{4t}{x(x-1)}.
\]
The cited theorem says exactly that this class is nonzero in the displayed quotient. Its essential obstruction can also be checked directly. All $E_g$ are noetherian, $t$-torsion-free, $t$-adically separated domains, with reductions $\mathbb C[x,1/g]$; their restriction maps are injective. Indeed separatedness follows from $t\in\operatorname{Jac}(E_g)$, and reduction modulo $t$, followed by repeated division by $t$, proves both the domain and injection assertions. Their elements are algebraic over $K(x)$, where $K=\mathbb C((t))$.

Suppose $f=u+v$ with $u\in E_x$, $v\in E_{x-1}$. The elementary quadratic decomposition lemma \cite[Lemma 3.1]{BGongyoKatoUchiba2026} implies $f\in K(x)(u)$ or $f\in K(x)(v)$. To verify that lemma, if both containments fail, the finite field $L=K(x)(u,v)$ has two involutions fixing $u$ and $v$, respectively, and negating $f$. Their product fixes $f$ and sends $u$ to $u-2f$ (or to $u+2f$, according to order). Its iterates are distinct in characteristic zero, contradicting finiteness of $\operatorname{Aut}_{K(x)}L$.

Hensel's lemma gives $\alpha\in tR$ with $\alpha^2-\alpha+4t=0$; put $\beta=1-\alpha$. Then
\[
 f^2=\frac{(x-\alpha)(x-\beta)}{x(x-1)},
 \qquad \alpha\beta=4t\ne0,\quad \alpha-\beta\in R^*.
\]
The universal property of henselization gives injections $E_{x-1}\hookrightarrow R[[x]]$ and $E_x\hookrightarrow R[[x-1]]$: their reductions modulo $t$ are injective, and the preceding separatedness argument applies. Since $R$ is complete, substitution $x=\alpha+z$ in the first, and $x=\beta+z$ in the second, gives field embeddings into $K((z))$. In either case the displayed rational function has valuation one, because its numerator has a simple zero while its denominator has nonzero constant term in $K$. It cannot be a square. Thus $f$ lies in neither field, a contradiction. Its nonzero quotient class proves the desired nonvanishing. The finite program verifies the cover identity, quadratic identity, factorization and the nilpotent-matrix identity underlying the involution calculation. The valuation and cohomology arguments are proofs, not finite computations.
\end{proof}

\begin{theorem}[Known, with derivation: completion of a punctured henselian pair]\label{B:Fujiwara}
Let $(R,I)$ be a henselian pair, with $I$ finitely generated, let $\widehat R=\varprojlim R/I^n$, and put $U=\operatorname{Spec}R\setminus V(I)$ and $\widehat U=\operatorname{Spec}\widehat R\setminus V(I\widehat R)$. For every torsion abelian \'etale sheaf $\mathcal G$ on $U$, pullback induces a quasi-isomorphism
\[
 R\Gamma(U,\mathcal G)\simeq R\Gamma(\widehat U,\mathcal G|_{\widehat U}).
\]
This is the Fujiwara--Gabber punctured comparison. It is distinct from comparing an affine scheme with its closed fibre in Theorem~\ref{B:henselian}.
\end{theorem}
\begin{proof}[Derivation from the opened proof]
Extend $\mathcal G$ by zero to $\mathcal G_1$ on $\operatorname{Spec}R$, and on quasi-compact quasi-separated $R$-schemes define
\[
 \mathcal F(Y)=R\Gamma(Y_{\acute et},\mathcal G_1|_Y).
\]
Bhatt--Mathew \cite[Theorem 6.11 and its proof]{BBhattMathew2021} applies their arc-descent theorem and the formal-gluing square to this functor. The required finitary property is the limit property of torsion \'etale cohomology. The required rigidity on affine rings is exactly Gabber's theorem in Theorem~\ref{B:henselian}. Their Corollary 6.10 says that a rigid finitary arc-sheaf with values in the nonnegative derived category is invariant under this passage from $U$ to $\widehat U$. Indeed, the formal-gluing square is homotopy cartesian, while its comparison on henselian affine terms is an equivalence by rigidity; the remaining comparison is therefore an equivalence. Applying this with the specified $\mathcal F$ proves the formula.

The opened source identifies Fujiwara \cite[Corollary 6.6.4]{BFujiwara1995} as the noetherian predecessor and Gabber's extension in \cite[Expos\'e XX, \S4.4]{BILO2014}. The original Fujiwara PDF was not obtained; the complete proof used here is the opened Bhatt--Mathew proof. Consequently R9's attribution must distinguish these two comparison theorems, rather than attribute the identical statement indiscriminately to both papers.
\end{proof}

\subsection{An explicit specialization obstruction in characteristic five}

\begin{proposition}[Refuted: invariance of the full fundamental group in subject 16]\label{B:specialization}
Let $k=\overline{\mathbb F}_5$, $R=k[[t]]$, and let $\mathcal E/R$ be the projective Weierstrass curve with equation
\[
 y^2z=x^3+t xz^2+z^3.
\]
It is smooth and proper with geometrically connected fibres. If $\bar\eta$ is a geometric generic point and $0$ the closed point, then
\[
 \dim_{\mathbb F_5}H^1_{\acute et}(\mathcal E_{\bar\eta},\mathbb F_5)=1,
 \qquad
 \dim_{\mathbb F_5}H^1_{\acute et}(\mathcal E_0,\mathbb F_5)=0.
\]
Thus the specialization homomorphism of full \'etale fundamental groups cannot be an isomorphism. The general valid replacement for a smooth proper family with geometrically connected fibres is surjectivity, with an isomorphism on maximal prime-to-$p$ quotients when the specialization has residue characteristic $p$.
\end{proposition}
\begin{proof}
The discriminant is $\Delta=-16(4t^3+27)$, whose reduction modulo $t$ is $3\in k^*$. The Weierstrass model is therefore a smooth projective genus-one curve over $R$, with the section $[0:1:0]$; all its geometric fibres are connected.

Over a perfect field of characteristic $5$, write $f(x)=x^3+a x+b$ and $\eta=dx/y$ on the elliptic curve $y^2=f(x)$. Since $y^5=y f(x)^2$, we have $\eta=f(x)^2dx/y^5$. The Cartier operator satisfies $C(g^5\omega)=gC(\omega)$ and
\[
 C\left(\sum_j c_jx^jdx\right)
   =\sum_m c_{5m+4}^{1/5}x^m dx.
\]
Only the coefficient of $x^4$ in $f(x)^2$, namely $2a$, contributes. Hence $C(\eta)=(2a)^{1/5}\eta$. By Serre duality, Frobenius on the one-dimensional space $H^1(E,\mathcal O_E)$ vanishes exactly when $a=0$. For $\mathcal E_{\bar\eta}$ it is nonzero; for $\mathcal E_0$ it is zero.

The Artin--Schreier sequence of \'etale sheaves
\[
 0\longrightarrow\mathbb F_5\longrightarrow\mathcal O_E
   \xrightarrow{F-1}\mathcal O_E\longrightarrow0
\]
and surjectivity of $u\mapsto u^5-u$ on an algebraically closed field identify $H^1_{\acute et}(E,\mathbb F_5)$ with the kernel of $F-1$ on $H^1(E,\mathcal O_E)$. In a basis this map is $u\mapsto cu^5-u$. If $c=0$ the kernel is zero. If $c\ne0$, the separable polynomial $cu^5-u$ has exactly five roots in the algebraically closed field, forming a one-dimensional $\mathbb F_5$-space. This proves the two dimensions. Since $H^1_{\acute et}(E,\mathbb F_5)=\operatorname{Hom}_{\mathrm{cont}}(\pi_1(E),\mathbb Z/5)$, isomorphic fundamental groups would give equal dimensions, a contradiction.

The replacement is Stacks Project Theorem 58.30.3, tag 0C0R, whose proof was opened: prime-to-$p$ covers extend across the specialization by tame ramification, while arbitrary finite covers need not do so. The displayed family supplies the missing explicit failure. The coefficient $2t$, the nonzero discriminant and independent finite-field point counts are checked by the attached program. The argument does not identify this family with Artin's particular unpublished counterexample mentioned in the margin.
\end{proof}

\begin{corollary}[Refuted: the naive same-form $p$-primary Euler formula]\label{B:p-Euler}
In characteristic $5$, the compactly supported Euler characteristic with constant $\mathbb F_5$ coefficients is not determined by the genus alone, even for smooth projective curves. Explicitly, over $\overline{\mathbb F}_5$,
\[
 E_{\mathrm{ord}}:y^2=x^3-x,
 \qquad E_{\mathrm{ss}}:y^2=x^3+1
\]
have genus $1$, whereas $\chi(E_{\mathrm{ord}},\mathbb F_5)=0$ and $\chi(E_{\mathrm{ss}},\mathbb F_5)=1$. Thus one cannot answer subject 3 by replacing $\ell\ne p$ by $\ell=p$ in the usual genus-and-Swan formula without additional invariants or hypotheses.
\end{corollary}
\begin{proof}
Both cubics have nonzero discriminant in characteristic $5$. The preceding Cartier calculation gives respectively $a=-1$ and $a=0$, hence $H^1$ dimensions $1$ and $0$. In the same Artin--Schreier sequence, $F-1$ on $H^1(E,\mathcal O_E)$ is surjective: it is $-1$ when $c=0$, and is the surjective polynomial $cu^5-u$ on an algebraically closed field when $c\ne0$. Since coherent cohomology vanishes above degree one on a curve, $H^2_{\acute et}(E,\mathbb F_5)=0$, and higher groups vanish as well. Each connected curve has $H^0=\mathbb F_5$, proving the two Euler characteristics. The source obstruction in Pink's introduction is thereby instantiated with two fully specified curves, rather than assumed from the word ``ordinary'' alone.
\end{proof}

\subsection{The twenty-eight subjects: exact endpoints and derivations}
\label{B:catalogue}

The manuscript is a list of research programmes, not twenty-eight uniformly quantified conjectures. In this catalogue ``known'' refers to the fully stated theorem that follows, with its actual hypotheses. It does not assert that every possible interpretation of a programme has been exhausted. The numbers below are the manuscript's numbers. All primary texts named as opened in the accompanying source audit were read at the stated locations. Original sources which could not be obtained are identified explicitly.

\subsubsection{1. Flat topology and its crystalline connection}
\begin{proposition}[Known: comparison for coefficients coming from the \'etale site]\label{B:subject1}
For a scheme $X$, an abelian sheaf $F$ on its small \'etale site, and every $q\ge0$, the canonical change-of-topology map gives
\[
 H^q_{\acute et}(X,F)\simeq H^q_{\mathrm{fppf}}(X,F_{\mathrm{fppf}}).
\]
Here $F_{\mathrm{fppf}}$ is the inverse-image sheaf on the big flat site; this is not an assertion about arbitrary flat sheaves. A precise crystalline companion is the following. Let $(A,I,\gamma)$ be a divided-power ring, $A$ noetherian and $p$-adically complete, with $p$ nilpotent in $A/I$. If $Y/A$ is proper smooth, $X=Y_{A/I}$, and $\mathcal E$ is a finite-type quasi-coherent crystal on $X/A$, then it determines a coherent module with integrable connection $(E_Y,\nabla)$ and
\[
 R\Gamma(\operatorname{Cris}(X/A),\mathcal E)
 \simeq R\Gamma(Y,E_Y\otimes\Omega^\bullet_{Y/A}).
\]
\end{proposition}
\begin{proof}[Derivation]
Stacks Project Lemmas 59.100.7--8, tags 0DDK and 0DDU, prove that the derived direct image of the inverse-image \'etale complex is the original complex. Applying derived global sections gives the first formula, with no invertibility assumption on $F$. For the second, Remark 60.24.14, tag 07MZ, applies the crystalline Poincar\'e comparison modulo $p^e$ for each $e$. The compatible evaluations of the crystal algebraize to $E_Y$ by proper formal existence; the connection algebraizes with them. Both displayed complexes are the derived inverse limits of their reductions, by crystalline completeness on the left and the theorem on formal functions on the right. Taking those limits proves the formula. These comparisons supply concrete formalism and the requested connection with crystals, without conflating the two topologies.
\end{proof}

\subsubsection{2. Local and global flat duality}
\begin{proposition}[Known: flat duality including non-\'{e}tale finite group schemes]\label{B:subject2}
Let $K$ be a local function field of characteristic $p$, and $N/K$ a finite commutative group scheme with Cartier dual $N^D$. With the canonical locally compact topologies, cup product gives perfect Pontryagin pairings
\[
 H^r_{\mathrm{fppf}}(K,N^D)\times H^{2-r}_{\mathrm{fppf}}(K,N)
 \longrightarrow\mathbb Q/\mathbb Z.
\]
Let $C$ be a smooth projective geometrically connected curve over a finite field and $N/C$ finite flat commutative. There are perfect pairings of finite groups
\[
 H^r_{\mathrm{fppf}}(C,N^D)\times H^{3-r}_{\mathrm{fppf}}(C,N)
 \longrightarrow\mathbb Q/\mathbb Z.
\]
For an open curve the second factor is compactly supported flat cohomology, with the topologies specified in the source.
\end{proposition}
\begin{proof}[Derivation]
These are Milne \cite[III, Theorems 6.10 and 8.2, and Corollary 8.9]{BMilne2006}. Cartier evaluation $N^D\times N\to\mathbb G_m$ gives the cup product. Locally it is followed by the invariant $H^2(K,\mathbb G_m)\simeq\mathbb Q/\mathbb Z$. Globally compact-support trace is in degree three, which explains the shift from $2-r$ to $3-r$. For $C$ proper, compact support is ordinary cohomology, and the finiteness assertion gives the displayed finite perfect pairings. Theorems for finite flat coefficients apply in particular to $\mu_p$ and its dual $\mathbb Z/p$; replacing the flat site by the \'etale site would lose this part of subject~2.
\end{proof}

\subsubsection{3. The characteristic-primary Euler formula}
\begin{proposition}[Known under stated hypotheses; naive extrapolation refuted]\label{B:subject3}
Let $X/k$ be a smooth projective connected curve over an algebraically closed field of characteristic $p>0$. Let $\mathbb F$ be a finite field of characteristic $p$, or a local field of residue characteristic $p$, and let $F$ be a constructible \'etale sheaf of finite-dimensional $\mathbb F$-vector spaces. Assume there is a finite irreducible Galois cover $Y\to X$ with $Y$ ordinary and the generic monodromy of the pullback of $F$ pro-$p$. Then
\[
 \chi(X,F)=(1-g_X)\operatorname{rank}F+
 \sum_{x\in X(k)}\bigl(\dim F_x-\operatorname{LT}_x(F_{\bar\eta})\bigr).
\]
Here $\operatorname{LT}$ is Pink's local term, specified below. Without ordinariness there is no formula obtained merely by using the usual prime-to-$p$ genus term: Corollary~\ref{B:p-Euler} gives two genus-one counterexamples. Independently, a finite \'etale Galois cover $V\to U$ of separated finite-type $k$-schemes with $p$-group Galois group satisfies
\[
 \chi_c(V,\mathbb Q_p)=(\deg V/U)\chi_c(U,\mathbb Q_p).
\]
\end{proposition}
\begin{proof}[Source and derivation]
The first formula is Pink \cite[Theorem 0.2; Definitions 5.2--5.5]{BPink2000}. For a finite coefficient field $\mathbb F$ of degree $r$ over $\mathbb F_p$, semisimplify the local inertia representation after each embedding $\iota:\mathbb F\hookrightarrow k$. Its simple constituents are tame characters $\chi$, represented by the rational numbers $a_\chi\in\mathbb Z_{(p)}\cap[0,1)$ under the chosen uniformizer convention. If their multiplicities are $m_{\iota,\chi}$, then
\[
 \operatorname{LT}(M)=\frac1r\sum_{\iota,\chi}m_{\iota,\chi}(1-a_\chi).
\]
For local-field coefficients one reduces an invariant lattice modulo the maximal ideal; the source proves independence of this choice. Thus the formula is explicit, rather than a name for an unspecified ramification correction. Pink's proof passes to the ordinary covering, compares the relevant \'etale and coherent invariants, and applies equivariant coherent Riemann--Roch; the ordinariness hypothesis is the point where the comparison is used.

The second assertion is Crew \cite[Theorem 1.5 and Corollary 1.7]{BCrew1984}. The theorem identifies the virtual compactly supported representation of a free $p$-group action as an integral multiple of its regular representation. Equivalently the associated perfect $\mathbb Z_p[G]$-complex has Euler class a multiple of the free-module class. Taking total dimension gives $|G|$ times that multiple, while taking invariants gives the Euler characteristic of the quotient. These two evaluations give the displayed equation. It addresses an exact part of the programme without supplying the false unrestricted $p$-primary genus formula.
\end{proof}

\subsubsection{4. Higher-dimensional Euler and Swan classes}
\begin{theorem}[Known: the higher-dimensional conductor formula]\label{B:subject4}
Let $U$ be a connected smooth separated scheme of finite type over a perfect field $k$, and let $\ell$ be invertible in $k$. For a lisse $\overline{\mathbb Q}_\ell$-sheaf $F$ there is a rational Swan zero-cycle class on the boundary, in the inverse limit of the Chow groups of compactification boundaries, such that
\[
 \chi_c(U_{\bar k},F)=\operatorname{rank}(F)\chi_c(U_{\bar k},\overline{\mathbb Q}_\ell)
                   -\deg\operatorname{Sw}(F).
\]
For a smooth curve its degree is the sum of the usual Swan conductors. In particular a finite \'etale cover $V\to X$ of degree $d$ of smooth proper $k$-schemes satisfies $\chi(V_{\bar k},\mathbb Q_\ell)=d\chi(X_{\bar k},\mathbb Q_\ell)$.
\end{theorem}
\begin{proof}[Derivation]
Kato--Saito \cite[Definition 4.2.8 and Theorem 4.2.9]{BKatoSaito2008} construct the class and prove the formula, including its independence of a coefficient lattice. Their construction attaches localized intersection classes to nonidentity ramification operators and weights them by character values; the trace computation in Corollary 3.4.6 and the conductor identities of \S4 give the asserted degree. When $U=X$ is proper the boundary class is zero. For a finite \'etale $f:V\to X$, the sheaf $f_*\mathbb Q_\ell$ is lisse of rank $d$, and $R^if_*\mathbb Q_\ell=0$ for $i>0$ because geometric fibres are finite discrete sets. The Leray identification $H^*(X,f_*\mathbb Q_\ell)=H^*(V,\mathbb Q_\ell)$ now proves the last assertion. Smoothness and prime-to-characteristic coefficients are explicit here; this is the precise higher-dimensional endpoint used, not a blanket assertion about every singular coefficient problem.
\end{proof}

\subsubsection{5. Euler characteristics with an operator algebra}
\begin{proposition}[Proved reduction to the conductor formula]\label{B:subject5}
Let $U$ be a smooth connected curve over an algebraically closed field of characteristic $p$, let $C$ be its smooth projective completion, and let $A/U$ be an abelian scheme. Fix $\ell\ne p$ and a finite-dimensional semisimple $\mathbb Q$-algebra $R$ acting on $A$. Put $V=T_\ell(A)\otimes\overline{\mathbb Q}_\ell$. After decomposing the split algebra $R\otimes\overline{\mathbb Q}_\ell$, write
\[
 V=\bigoplus_\rho W_\rho\otimes M_\rho,
\]
where $W_\rho$ are its distinct simple modules and $M_\rho$ are lisse multiplicity sheaves. In its representation Grothendieck group,
\[
 \sum_i(-1)^i[H_c^i(U,V)]
 =\chi_c(U,\overline{\mathbb Q}_\ell)[V_{\bar\eta}]
 -\sum_{x\in C\setminus U}\sum_\rho\operatorname{Sw}_x(M_\rho)[W_\rho].
\]
\end{proposition}
\begin{proof}
The operator action commutes with monodromy, so the central idempotents and matrix units of the split semisimple algebra give the displayed decomposition into constant operator modules tensored with local systems. Tensoring by a finite-dimensional constant vector space and taking these direct summands commute with compactly supported cohomology. Consequently the left-hand side is $\sum_\rho\chi_c(U,M_\rho)[W_\rho]$. Apply the curve case of Theorem~\ref{B:subject4} separately to each $M_\rho$ and sum. This proves the formula coefficient by coefficient in a basis of the representation group, hence proves equality there. It makes the operator refinement precise. Integral operator lattices, nonsemisimple coefficient algebras and characteristic-primary coefficients are additional questions, not consequences of this rational decomposition.
\end{proof}

\subsubsection{6. Explicit Lefschetz formulas}
\begin{proposition}[Known: discrete and coherent fixed-point formulas]\label{B:subject6}
For a separated finite-type $X/\mathbb F_q$, a constructible $\mathbb Q_\ell$-sheaf $F$ with $\ell\nmid q$, and geometric Frobenius $\Phi$,
\[
 \sum_i(-1)^i\operatorname{Tr}(\Phi\mid H_c^i(X_{\overline{\mathbb F}_q},F))
 =\sum_{x\in X(\mathbb F_q)}\operatorname{Tr}(\Phi_x\mid F_{\bar x}).
\]
For a compact complex manifold $M$, a holomorphic $f:M\to M$ with only fixed points satisfying $\det_{\mathbb C}(1-df_x)\ne0$, a holomorphic vector bundle $E$ and a lift $u:f^*E\to E$,
\[
 \sum_q(-1)^q\operatorname{Tr}(f,u\mid H^q(M,E))
 =\sum_{f(x)=x}\frac{\operatorname{Tr}(u_x)}{\det_{\mathbb C}(1-df_x)}.
\]
\end{proposition}
\begin{proof}[Derivation]
For the first, the curve case is \cite[Theorems 29.4 and 29.15]{BMilne2013}. Constructible coefficients on a curve reduce by a finite stratification to the stated lisse smooth case and finite sets; nilpotents do not change the \'etale site. Here is the dimension induction giving the stated generality. The two sides are additive for an open--closed decomposition, by the compact-support localization triangle. Stratify $X$ into affine integral pieces. For a piece of dimension $d>1$, Noether normalization followed by projection gives a morphism to $\mathbb A^{d-1}$ all of whose fibres have dimension at most one. The constructible complexes $Rf_!F$ have stalks $R\Gamma_c(X_{\bar y},F)$ by compact-support base change. Apply the dimension-$(d-1)$ formula to their cohomology sheaves; the fibrewise curve formula converts each stalk's alternating trace into the sum over rational points of that fibre. The compact-support Leray spectral sequence identifies the resulting global alternating trace with that on $X$. This proves the first formula by induction. The local Frobenius acts on the indicated stalk and compact support is necessary for nonproper $X$.

The second follows from Atiyah--Bott \cite[Theorem A, pp.~377--378]{BAtiyahBott1967} applied to the Dolbeault elliptic complex of $E$. Its local numerator is
\[
 \operatorname{Tr}(u_x)\sum_q(-1)^q\operatorname{Tr}(\wedge^q\overline{df_x})
 =\operatorname{Tr}(u_x)\det_{\mathbb C}(1-\overline{df_x}).
\]
The real denominator is $|\det_{\mathbb R}(1-df_x)|=|\det_{\mathbb C}(1-df_x)|^2$. Cancelling the conjugate determinant gives the stated fraction. The simple-fixed-point hypothesis ensures the denominator is nonzero. These are explicit realizations of both coefficient regimes on the sheet; general fixed loci require their own local terms.
\end{proof}

\subsubsection{7. Global duality on arithmetic schemes}
\begin{theorem}[Known: duality by a cycle complex]\label{B:subject7}
Let $B$ be the spectrum of the ring of integers of a number field and let $f:X\to B$ be separated of finite type. Let $\mathbb Z_X^c$ denote Bloch's dualizing cycle complex. For a torsion \'etale sheaf $F$ on $X$ there is a natural duality
\[
 R\operatorname{Hom}_X(F,\mathbb Z_X^c)
 \simeq R\operatorname{Hom}_{\mathbb Z}(R\Gamma_c(X,F),\mathbb Z)[-1].
\]
For constructible $F$ it induces perfect finite pairings
\[
 \operatorname{Ext}^{2-i}_X(F,\mathbb Z_X^c)\times H_c^i(X,F)
 \longrightarrow\mathbb Q/\mathbb Z.
\]
At real places compact support is understood with the source's Tate modification.
\end{theorem}
\begin{proof}[Derivation]
This is Geisser \cite[Theorem 7.8]{BGeisser2010}. It applies to $X$ itself, with no smoothness or properness imposed beyond the stated separated finite-type hypothesis. For a finite group $M$, $\operatorname{Hom}(M,\mathbb Z)=0$ and $\operatorname{Ext}^1(M,\mathbb Z)=\operatorname{Hom}(M,\mathbb Q/\mathbb Z)$, by the injective resolution $\mathbb Z\to\mathbb Q\to\mathbb Q/\mathbb Z$. Applying this to the finite cohomology groups in the duality complex yields the degree $2-i$ and the displayed perfect pairings. Taking $B=\operatorname{Spec}\mathbb Z$ answers the arithmetic-base direction explicitly. Omitting the real-place modification would change the theorem.
\end{proof}

\subsubsection{8. Differential operators, connections and crystals}
\begin{proposition}[Known, with an explicit duality consequence]\label{B:subject8}
Let $M$ be a complex manifold. Holomorphic finite-rank vector bundles with integrable connection are equivalent to finite-dimensional complex local systems, compatibly with tensor products, duals and internal Hom. If $M$ is compact of complex dimension $d$, the associated de Rham groups have perfect pairings
\[
 H^i_{\mathrm{dR}}(M,E)\times H^{2d-i}_{\mathrm{dR}}(M,E^\vee)
 \longrightarrow\mathbb C.
\]
For a smooth scheme in characteristic zero an integrable connection is the same as a left module structure over the sheaf of differential operators on the underlying vector bundle. In the divided-power characteristic-$p$ situation of subject~1, the crystal--connection correspondence additionally requires topological quasi-nilpotence.
\end{proposition}
\begin{proof}[Derivation]
Deligne \cite[I, Theorem 2.17 and Propositions 2.18--2.19]{BDeligne1970} gives the equivalence and identifies the de Rham complex with a resolution of the horizontal local system $L$. Hence its hypercohomology is $H^i(M,L)$. Orientation and Poincar\'e duality on the compact real $2d$-manifold, with local coefficients paired by evaluation $L\otimes L^\vee\to\mathbb C$, give the formula. In characteristic zero, define the action of a vector field $v$ by $\nabla_v$. The Leibniz rule is precisely its commutator relation with functions, and integrability gives $[\nabla_v,\nabla_w]=\nabla_{[v,w]}$. The defining relations of differential operators therefore extend this action uniquely, and restriction recovers the connection.

For crystals, Stacks Lemma 60.17.4, in Section 60.17, identifies quasi-coherent crystals with complete modules carrying an integrable topologically quasi-nilpotent connection on the completed divided-power envelope. The Taylor comparison is $m\mapsto\sum_K\nabla^K(m)\xi^{[K]}$: quasi-nilpotence gives convergence and integrability makes the operators commute, so the binomial identity gives the cocycle condition. This is why an arbitrary integrable characteristic-$p$ connection cannot replace a crystal. Subject~1 records the resulting proper comparison theorem with all base assumptions.
\end{proof}

\subsubsection{9. Analytic duality}
\begin{theorem}[Known: duality with the necessary topologies]\label{B:subject9}
Let $X$ be a complex analytic space, countable at infinity and of bounded dimension. There is a dualizing complex $K_X$ and trace such that, for a coherent sheaf $F$, the separated quotients of
\[
 H^p(X,F)\quad\hbox{and}\quad\operatorname{Ext}^{-p}_c(X;F,K_X)
\]
are perfect strong topological duals in the canonical quotient-Fr\'echet and quotient-DFS topologies. Similarly $H_c^p(X,F)$ is dual to $\operatorname{Ext}^{-p}(X;F,K_X)$ after separation. When $X$ is compact, these are the usual finite-dimensional perfect pairings.
\end{theorem}
\begin{proof}[Derivation]
Ramis--Ruget \cite[Theorem 1, \S3, p.~81; Theorem 2, \S4]{BRamisRuget1970} construct these pairings. They first obtain the smooth local trace from integration of distribution-valued forms, then define the dualizing complex on a singular embedding by derived Hom and prove compatibility under a change of embedding. Those compatibilities glue the local complexes and traces. Applying the theorem to the coherent $F$ gives the displayed groups with the indicated shifts. For compact $X$, coherent cohomology is finite dimensional, compact support is ordinary support, and the separated quotients equal the original vector spaces. Thus the modern answer does not silently omit the topological separation issue on noncompact spaces.
\end{proof}

\subsubsection{10. Riemann--Roch for coherent and perfect complexes}
\begin{proposition}[Known: algebraic and analytic Riemann--Roch]\label{B:subject10}
For a quasi-projective scheme $X$ over a field, there is a natural transformation
$\tau_X:G_0(X)\to\mathrm{CH}_*(X)_{\mathbb Q}$, covariant for proper morphisms, satisfying
\[
 \tau_X(E\otimes F)=\operatorname{ch}(E)\cap\tau_X(F),
 \qquad \tau_X(\mathcal O_X)=\operatorname{td}(T_X)\cap[X]
\]
when respectively $E$ is a vector bundle and $X$ is smooth. More generally, for a projective local-complete-intersection morphism $f:X\to Y$ of constant virtual relative dimension, where $Y$ is quasi-compact with an ample line bundle, SGA~6 gives
\[
 \operatorname{ch}(f_!x)=f_{\mathrm{gr},*}
       \bigl(\operatorname{ch}(x)\operatorname{td}(T_f)\bigr)
\]
in the rational associated graded of $K$-theory, for perfect-complex classes $x$.
For complex analytic spaces there is a proper-covariant transformation from coherent $G_0$ to topological $K$-homology, agreeing with the usual bundle class under Poincar\'e identification on a complex manifold.
\end{proposition}
\begin{proof}[Derivation]
The first transformation and identities are Baum--Fulton--MacPherson \cite[Introduction theorem; Chapters I.6 and II.3]{BBFM1975}. Proper naturality immediately gives $\tau_Y(\sum_i(-1)^iR^if_*F)=f_*\tau_X(F)$; for $Y$ a point this is the numerical Euler formula. Bounded complexes with coherent cohomology are handled by their alternating class in $G_0$. The second displayed theorem is exactly SGA~6 \cite[Expos\'e VIII, Theorem 3.6, internal p.~20]{BSGA6}; the original scan was inspected, including the ample and local-complete-intersection hypotheses. Its $T_f$ is the dual class of the relative cotangent complex, so this statement does not presuppose a smooth tangent bundle.

For the analytic assertion, Levy \cite[Introduction, Riemann--Roch theorem, pp.~149--150]{BLevy1987} constructs $\alpha_X:G_0^{\mathrm{hol}}(X)\to K_0^{\mathrm{top}}(X)$ with $\alpha_Yf_!=f_*\alpha_X$ for every proper holomorphic $f$. Grauert's theorem makes $f_!$ well defined. On smooth spaces its bundle normalization is the Dolbeault orientation; applying the homological Chern character produces the familiar Todd-class form. An unbounded pseudo-coherent complex has no automatic finite alternating Euler class; a boundedness or perfectness hypothesis is therefore retained where the formula uses one.
\end{proof}

\subsubsection{11. Piecewise structures and the relative N\'eron construction}
\begin{theorem}[Known: the ordinary abelian N\'eron-model functor]\label{B:subject11}
Let $R$ be a discrete valuation ring, $K$ its fraction field, and $A/K$ an abelian variety. The functor on smooth $R$-schemes
\[
 T\longmapsto\operatorname{Hom}_K(T_K,A)
\]
is represented by a smooth separated group scheme $\mathcal A/R$ of finite type. Its generic fibre is $A$ and its representing property characterizes it up to unique isomorphism.
\end{theorem}
\begin{proof}[Derivation]
Bosch--L\"utkebohmert--Raynaud \cite[\S1.3, Theorem 2; \S1.4, Theorem 3]{BBLR1990} provide the N\'eron model and its mapping property. Restricting that property to smooth test schemes gives exactly the displayed functor. Uniqueness follows by extending the identity on $A$ in both directions and applying uniqueness of extension to the two composites. Likewise multiplication, inverse and identity extend from the generic fibre because products of smooth $R$-schemes are smooth; their group identities follow from the same uniqueness. The handwritten term ``sch\'emas morcel\'es'' is visually clear but is not defined in cote~66. This is a theorem about the ordinary smooth-test N\'eron-model functor. It does not by itself identify a historical functor defined on another category through finite constructible partitions. The following theorems give the directly relevant finite-partition and relative-cohomology comparisons; subject~12 independently asks about N\'eron models.
\end{proof}

\paragraph{The original partition construction and the scope of the answer.}
Original cote~43, folios~91,93,95--96 (PDF pages one higher), gives a
complete discrete valuation ring with perfect residue field, finite
surjective descent on the target, and a functor
$\mathcal N:\mathcal C(S)\to\mathcal C'(k)$ explicitly named after
N\'eron. Its formula is a colimit over finite constructible partitions,
applying the valuation-ring construction to each piece before evaluating
the sheaf \cite[folio~95]{A567_GrothendieckCote43}.
The next folio proposes a relative-cohomology interpretation but marks
injective preservation for verification and crosses out an argument.
We do not attribute a completed proof to that proposal. Folio~85
separately defines a localization at integral radicial surjections;
this does not identify it with the finite-piecewise category below.

EGA~I connects its constructible-topology scheme to N\'eron reduction
\cite[I, 7.2.14, pp.~338--339]{ROOT_EGAI1971}; the precise relation to
Olivier's universal absolutely flat construction is recorded by
\cite[8.3, p.~581]{ROOT_Ferrand2003}.
We answer the catalogue's expressly permitted strongest-proved-statement
alternative by the complete category and group-object comparisons,
finite and complete Greenberg realizations, exact infinite obstruction,
and perfect-test relative cohomology below. Every test category is
specified. We do not assert the lexical identity of every occurrence of
``morcel\'es'', define the still unverified historical ``flasque''
qualification by guesswork, or identify all historical and modern sites.

% Audited C-owned proof fragment, integrated in subject11 on 2026-11-19.
% Historical scope: this defines a precise finite-type category; it does not
% identify the 1964 expression `schemas morceles'.
\paragraph{A precise finite-type category of piecewise morphisms.}

Let $k$ be an algebraically closed field.  A constructible $k$-set means a
constructible subset of the $k$-points of a separated reduced finite-type
$k$-scheme.  A map is \emph{piecewise regular} if its source admits a finite
partition into locally closed reduced $k$-subschemes on each of which the map
is induced by a $k$-morphism.  Let $\mathcal P_k$ be this category.  Finite
products and their usual universal properties exist in $\mathcal P_k$.
Every object can, by using a finite affine stratification and disjoint tags,
be encoded as a constructible subset of an affine space.  The choice of
encoding does not change the category up to equivalence.

Write $\mathcal C_k$ for the category with the same objects but all maps
whose graphs are constructible.  In positive characteristic these categories
need not have the same morphisms.  Neither category records nilpotents.
An algebraic group in this subsection is a separated reduced finite-type
group scheme over $k$, and is not assumed affine or connected.

\begin{lemma}[Proved: piecewise morphisms and graphs]\label{C:piecewise-graph-category}
A piecewise regular map has constructible graph.  Conversely, in
characteristic zero every map with constructible graph is piecewise regular.
In arbitrary characteristic, a constructible graph whose projection to the
source is bijective on all scheme points and induces isomorphisms on their
residue fields is piecewise isomorphic over the source, and hence defines a
piecewise regular map.  Thus the last graph convention, interpreted on
finite-type reduced schemes and finite stratifications, gives
$\mathcal P_k$.
\end{lemma}

\begin{proof}
The graph of a morphism to a separated scheme is closed in the product of
its source with the target; a finite union over locally closed source
strata is constructible.  Composition is piecewise regular: refine the
source strata by the inverse images of strata for the second map, then
refine these constructible sets into locally closed strata.

For the converse in characteristic zero, use noetherian induction on the
closure of the source.  It suffices to find a dense open regularity locus
on each irreducible component.  Restrict to an irreducible source $V$ and
to affine source and target charts meeting the generic graph.  The graph
has one point over each point of its source.  The same uniqueness holds
after algebraically closed field extension: a failure would define a
nonempty constructible incidence set over $k$, which would already have a
$k$-point.  Consequently the dominant irreducible component of the graph
closure is generically finite over $V$ with purely inseparable function
field extension.  There is only one such component, since two dominant
components would give two points in a geometric generic fibre.  In
characteristic zero its function field is $k(V)$.  A birational morphism
of varieties is an isomorphism on suitable dense opens.  Shrinking further
puts that open of the graph closure inside the actual graph and excludes
the images of its remaining strata.  Its inverse then supplies the
required morphism on an open of $V$.  The complement has smaller dimension,
so induction gives a finite stratification.

Under the stated residue-field condition the same argument works in any
characteristic: the generic extension is already trivial.  Conversely, a
finite graph stratification by sections identifies each graph stratum
with its source stratum, so the projection is bijective on all scheme
points with identical residue fields.  The underlying facts about
radicial maps and birational open isomorphisms are precisely Stacks,
Lemmas 29.10.2 and 29.51.5 (tags 01S4 and 0BAC)\cite[tags 01S4 and 0BAC]{BStacks}.
\end{proof}

\begin{lemma}[Proved: regularity of piecewise homomorphisms]\label{C:piecewise-group-fullness}
Every piecewise regular homomorphism $f:G(k)\to H(k)$ between algebraic
$k$-groups is induced by a unique homomorphism of algebraic groups.
\end{lemma}

\begin{proof}
Remove the other irreducible components from one irreducible component
of $G$, and then choose a nonempty open regularity stratum for $f$.
This gives a $V$ open in all of $G$ on which $f$ is regular.
Choose $v\in V(k)$.  For each $g\in G(k)$ the open
$g v^{-1}V$ contains $g$.  On this open the homomorphism identity gives
\[
 f(z)=f(gv^{-1})\,f(vg^{-1}z).
\]
The right-hand side is a morphism.  These morphisms agree on overlaps,
since they agree on all $k$-points and the schemes are reduced.  They glue
to a morphism on $G$.  The group identities hold as scheme morphisms by
the same reducedness argument, since products of reduced finite-type
schemes over the algebraically closed field $k$ are reduced.
Uniqueness also follows from density of
closed points.  No restriction on the characteristic is used here.
\end{proof}

\begin{theorem}[Proved: algebraization in the specified finite-type category]\label{C:piecewise-group-algebraization}
The functor from algebraic $k$-groups to group objects in $\mathcal P_k$ is
an equivalence of categories.  In particular, if $\operatorname{char}k=0$,
it is an equivalence with group objects in $\mathcal C_k$.
\end{theorem}

\begin{proof}
Full faithfulness is the preceding lemma.  We prove essential surjectivity.
Let $D$ be a constructible set with group operations that are piecewise
regular.  Encode $D$ in an affine space and put $d=\dim D$.  If $d=0$,
then $D$ is finite and the constant finite group scheme proves the claim.
Assume $d>0$.

\emph{Step 1: a normal finite-index subgroup with one generic component.}
Choose disjoint constructible subsets $C_1,\ldots,C_m$ of $D$, each dense
open in one of the distinct $d$-dimensional irreducible components of
$\overline D$, such that the complement of their union has dimension less
than $d$.  Two constructible subsets are identified here if their
symmetric difference has dimension less than $d$.  Modulo this relation,
the $C_i$ are the atoms of the finite Boolean algebra of constructible
subsets of $D$.

A piecewise regular bijection with piecewise regular inverse preserves
dimension and these generic atoms.  Hence left translation gives an
action of $D$ on $\{C_1,\ldots,C_m\}$.  The action is constructible: for
fixed $i,j$ its relevant fibre is
\[
 \{g\in D:\dim(gC_i\cap C_j)=d\}.
\]
Constructibility follows by stratifying the incidence correspondence and
applying constructibility of fibre dimension (Stacks, Lemma 37.30.3,
tag 05F9)\cite[tag 05F9]{BStacks}.  This applies to piecewise regular maps by a finite
stratification into ordinary morphisms.

The action is transitive.  Indeed, the piecewise regular map
\[
 C_j\times C_i\longrightarrow D,\qquad (b,a)\longmapsto ba^{-1}
\]
has a source of dimension $2d$.  Its fibre over $g$ is isomorphic as a
constructible set to $\{a\in C_i:ga\in C_j\}$.  If every fibre had
dimension at most $d-1$, the dimension-of-fibres inequality would bound
the source dimension by $2d-1$, a contradiction.  A fibre of dimension
$d$ shows that some $g$ carries the generic atom of $C_i$ to that of $C_j$.

Let $D^0$ be the stabilizer of $C_1$.  It is constructible, has index $m$,
and has dimension $d$.  If it had $s$ generic atoms, each of its $m$
cosets would have $s$ atoms, and the disjoint coset partition would give
$m=ms$.  Thus $s=1$.  A constructible subgroup of finite index in $D^0$
must therefore be all of $D^0$.  Applying this to
$D^0\cap gD^0g^{-1}$ shows that $D^0$ is normal in $D$.

\emph{Step 2: obtain a strict birational group law.}
Let $V$ be the unique $d$-dimensional generic component of $D^0$.
Choose a dense open $V_0\subset V$ contained in $D^0$.  Every translation
of $D^0$ preserves its unique generic atom.  In particular the product
of two independent generic points lies in that atom: for each fixed
$a$, the bad set of $b$ for which $ab$ lies outside $V_0$ has dimension
less than $d$, by translation invariance; the fibre-dimension inequality
therefore bounds the total bad set in $D^0\times D^0$ by $2d-1$.
The multiplication,
left-division and right-division maps are piecewise regular, so the three
pairwise projections of the generic multiplication graph induce
birational maps to $V\times V$.  More explicitly,
$(a,b)\mapsto(a,ab)$ and $(a,b)\mapsto(ab,b)$ have the piecewise regular
inverses $(a,c)\mapsto(a,a^{-1}c)$ and
$(c,b)\mapsto(cb^{-1},b)$.  On dense opens these are inverse rational
maps, and all three coordinates lie in $V_0$.

Take a dense open $W_0$ of the closure of this graph which is contained
in the actual multiplication graph, lies in $V_0^3$, and on which all
three pairwise projections are open immersions.  Denote their dense open
images in $V^2$ by $U_{12},U_{13},U_{23}$.  The projections of each
$U_{ij}$ onto either factor are dense opens of $V$.  Let $X\subset V_0$
be their intersection, and set $W=W_0\cap X^3$.

We check the fibre-density condition, which cannot be omitted.  Fix any
one coordinate $a\in X$.  The fibre of $W_0$ at this coordinate maps
isomorphically, by either remaining coordinate, onto a nonempty open of
$V$: the two relevant pairwise projections are open immersions, and
$a$ belongs to both projection images by the choice of $X$.  Thus it
is an irreducible open variety, birational to $V$.  Requiring the other
two coordinates to lie in $X$ removes only proper closed subsets of
this fibre.  The resulting fibre is nonempty, and both of its coordinate
images are dense opens of $X$.  It follows that every pairwise projection
of $W$ is an open immersion into $X^2$ with dense fibres over either
factor.  Associativity on its domain is inherited from $D^0$.
Consequently $(X,W)$ satisfies exactly Definition 2.1 of
Edixhoven--Romagny\cite[Definition 2.1, pp.~4--5]{CEdixhovenRomagny2013}, not merely a weaker generic associativity condition.

\emph{Step 3: algebraize and identify the entire connected part.}
Apply Edixhoven--Romagny, Theorem 2.11, whose complete proof is in their
Section 2\cite[Theorem 2.11 and its proof, pp.~4--7]{CEdixhovenRomagny2013}.  It constructs an algebraic group $A$ containing $X$ as a
dense open.  Its underlying abstract group is the subgroup of
$\operatorname{Bir}(X)$ generated by the rational left translations
$\lambda_x$ for $x\in X(k)$.  It is separated and of finite type, since
$X$ is of finite type.  It is irreducible because $X$ is irreducible and
dense.

For every $h\in D^0$, left translation by $h$ gives a birational
self-map of $V$, hence an element $\rho(h)$ of
$\operatorname{Bir}(X)$.  The map $\rho$ is a homomorphism and is
injective: if $ha=a$ on a nonempty open, cancellation gives $h=1$.
Moreover $D^0=X X^{-1}$.  To see this, for any $h$ choose
$y\in X\cap h^{-1}X$, which is nonempty since both sets have the unique
generic atom, and write $h=(hy)y^{-1}$.  Thus $\rho(D^0)$ is precisely
the subgroup generated by the $\lambda_x$, namely $A(k)$.

This isomorphism $\rho:D^0\to A(k)$ is constructible: its graph is the
image of the constructible map
\[
 X\times X\longrightarrow D^0\times A,
 \qquad (x,y)\longmapsto
 \bigl(xy^{-1},\lambda_x\lambda_y^{-1}\bigr).
\]
We must also check piecewise regularity in arbitrary characteristic.
The translates of the dense open $X\subset A$ cover $A$, and
quasi-compactness of the finite-type scheme $A$ supplies a finite
subcover $\rho(h_i)X$.  Pulling this cover back through the already
established abstract group isomorphism gives $D^0=\bigcup_i h_iX$.
On $h_iX$, the map $\rho$ is the composite of the
piecewise regular translation $h_iX\to X$ with the regular map
$x\mapsto\rho(h_i)\lambda_x$.  Hence $\rho$ is piecewise regular.
The corresponding opens $\rho(h_i)X$ cover $A$, and on each of them the
inverse is a regular chart followed by a piecewise regular translation
in $D^0$.  Therefore $\rho$ is an isomorphism in $\mathcal P_k$.

\emph{Step 4: restore the finitely many components.}
Choose left coset representatives $r_i$ for $D/D^0$.  On each coset
$r_iD^0$ transport the variety structure of $A$ by left translation.
For $r_ir_jD^0=r_\ell D^0$ put
$c_{ij}=r_\ell^{-1}r_ir_j\in D^0$.  In the transported coordinates,
multiplication is
\[
 (r_i a)(r_j b)=r_\ell c_{ij}(r_j^{-1}ar_j)b.
\]
Conjugation by $r_j$ is a piecewise regular automorphism of $D^0$;
after transport to $A$ it is regular by
Lemma~\ref{C:piecewise-group-fullness}.  Thus multiplication is regular
on every pair of components.  The same formula applied to inverses
shows that inversion is regular.  The finite disjoint union is a
separated reduced finite-type algebraic group, and its identification
with $D$ is piecewise regular in both directions.  This proves essential
surjectivity and the equivalence.  Characteristic zero follows from
Lemma~\ref{C:piecewise-graph-category}.
\end{proof}

\paragraph{Positive-characteristic distinction.}
For $\operatorname{char}k=p>0$, the inverse of Frobenius on
$\mathbb G_a(k)$ is an additive homomorphism with closed graph
$\{(x,y):y^p=x\}$.  It is not regular: a polynomial $P(x)$ satisfying
$P(x)^p=x$ cannot exist by degree.  It is not even rational on a dense
open, since $x$ is not a $p$th power in $k(x)$ (use its valuation at
$x=0$).  It therefore does not belong to $\mathcal P_k$.  Its graph
projection has generic residue-field extension $k(x)\subset k(y)$,
$x=y^p$, of degree $p$, so it also fails the residue-field-trivial
convention.  The ordinary functor from algebraic groups to group objects
in $\mathcal C_k$ is consequently not full in characteristic $p$.
The theorem above makes no unproved assertion that all objects of that
larger characteristic-$p$ category are algebraizable; the point here
is the proved equivalence for $\mathcal P_k$ and the precise obstruction
to replacing its morphisms by arbitrary constructible graphs.

\paragraph{Historical limitation.}
Hrushovski's \emph{Valued fields, metastable groups}, introduction,
pp.~1--2\cite[pp.~1--2 and p.~2, footnote 3]{CHrushovskiMetastableDraft}, records both the characteristic-zero equivalence and the
residue-field-trivial convention, attributing the latter to Grothendieck's
letter to Poizat of 17 June 1979.  Poizat's \emph{Logique et g\'eom\'etrie},
Sections 2.B and 3.B\cite[Sections 2.B and 3.B]{CPoizat2025}, spells out the characteristic distinction.  These
opened sources justify the specified category.  They do not identify
it with every scheme, pro-affine object, sheaf or partition construction
that might have been intended by the 1964 expression ``sch\'emas
morcel\'es''.  That historical identification remains a separate task.

% C-owned proof fragment, 2026-11-19.  No historical identity is assumed.
\paragraph{Universal absolute flatness and finite piecewise morphisms.}

Let $\operatorname{Abs}_k$ be the full category of $k$-schemes whose
local rings are fields.  For a scheme $X$, let $t_X:T(X)\to X$ be its
universal absolutely flat scheme.  The functor $T$ is right adjoint to
the inclusion of absolutely flat schemes into schemes.  Thus
\[
 \operatorname{Hom}_{\operatorname{Abs}_k}(Z,T(X))
   =\operatorname{Hom}_{\operatorname{Sch}_k}(Z,X)
 \quad (Z\in\operatorname{Abs}_k).
\]
The scheme construction and its relation to the constructible topology
are recorded in Ferrand, Section 8.3, explicitly identifying it with
Grothendieck's $X^{\mathrm{cons}}$ in EGA I, 7.2.14
\cite{ROOT_Ferrand2003,ROOT_EGAI1971}.  On an affine it is
\[
 T(\operatorname{Spec}A)=\operatorname{Spec}
 \left(A[t_a:a\in A]/(a^2t_a-a,\,at_a^2-t_a:a\in A)\right).
\]
The two relations specify the unique pointwise inverse of $a$.
The universal property, residue fields and constructible topology follow
from Olivier's explicit construction
\cite[expos\'e 6]{ROOT_Olivier1968}\cite[Sections 1--3]{ROOT_Olivier1978}.
The morphism $t_X$ is affine, surjective and a monomorphism.

\begin{theorem}[Proved: comparison with finite piecewise morphisms]\label{C:absolute-flat-piecewise}
Let $k$ be algebraically closed, and let $X,Y$ be separated reduced
finite-type $k$-schemes.  There are natural bijections
\[
 \operatorname{Hom}_{\mathcal P_k}(X,Y)
 \simeq \operatorname{Hom}_{\operatorname{Sch}_k}(T(X),Y)
 \simeq \operatorname{Hom}_{\operatorname{Abs}_k}(T(X),T(Y)).
\]
Consequently $T$ gives a fully faithful realization of $\mathcal P_k$
as the full subcategory of $\operatorname{Abs}_k$ with objects $T(X)$
arising from separated finite-type reduced schemes (and their finite stratifications).
Its categorical product is
\[
 T(X)\mathbin{\times_{\operatorname{Abs}_k}}T(Y)
   \simeq T(X\times_kY)
   \simeq T\bigl(T(X)\times_kT(Y)\bigr).
\]
The product occurring inside the last $T$ is the ordinary scheme product.
It need not itself be absolutely flat.
\end{theorem}

\begin{proof}
We give the finite-partition limit argument, including the reduction
and finiteness points.  Let $\mathfrak P(X)$ be the directed system of
finite partitions of $X$ into affine locally closed reduced subschemes.
Such partitions exist by noetherian induction.  Common refinements
exist: intersect two partitions and, if necessary, subdivide the
intersections into finitely many affine locally closed reduced pieces.
Write $X_P=\coprod_{S\in P}S$.  Each $X_P$ is affine, of finite type,
and reduced.  For a refinement $P'$ of $P$ there is a natural morphism
$X_{P'}\to X_P$.  These transition morphisms are affine: the inverse
image of each affine component is a finite disjoint union of affine
components.  Let $L=\varprojlim_P X_P$; it is an affine scheme.

Its coordinate ring is absolutely flat.  A function $b$ in that ring
comes from some $\Gamma(X_P,\mathcal O)$, since the ring is a filtered
colimit.  Subdivide every component of $P$ into $D(b)$ and $V(b)_{
\mathrm{red}}$.  These are affine locally closed reduced subschemes,
so they form a refinement.  On the resulting disjoint union define
$b^{[-1]}$ to be $b^{-1}$ on $D(b)$ and zero on $V(b)_{
\mathrm{red}}$.  It satisfies
$b^2b^{[-1]}=b$ and $b(b^{[-1]})^2=b^{[-1]}$.  Thus every element of
the limit ring has a pointwise inverse.  Reduction of the closed
strata is essential here: no nilpotent thickening has been retained.

We next identify $L$ with $T(X)$ by the universal property.  If $Z$ is
absolutely flat and $f:Z\to X$ is a morphism, the inverse image of every
constructible subset of $X$ is clopen in $Z$.  This can be checked
locally on affine opens: every principal open in the spectrum of an
absolutely flat ring is clopen, and each constructible subset of the
noetherian scheme $X$ is a finite Boolean combination of such opens.
Moreover, a morphism from a reduced scheme whose image is contained
in a locally closed reduced subscheme factors through that subscheme;
on a closed affine piece, its defining ideal maps into all primes and
hence to zero.  It follows that $f$ factors uniquely through $X_P$
for every $P$, using the clopen inverse-image pieces of $Z$.  The
factorizations are compatible with refinement.  Hence $f$ factors
uniquely through $L$.  This proves $L\simeq T(X)$.

Since $Y$ is of finite presentation over $k$, the standard limit
theorem gives
\[
 \operatorname{Hom}_k(T(X),Y)
 =\varinjlim_P\operatorname{Hom}_k(X_P,Y).
\]
Here all hypotheses are present: the stages are quasi-compact and
quasi-separated, and the transition maps are affine.  One may apply
Stacks, Lemma 32.10.1, parts (2) and (3), tag 01ZM\cite[tag 01ZM]{BStacks}, to the section of
$Y\times_k T(X)\to T(X)$ corresponding to a morphism; both this
product and the identity source descend to each stage.  This proves
both existence of a finite stage and equality after a common
refinement.  A morphism $X_P\to Y$ is exactly a map regular on every
piece of $P$.  Two such maps agree as maps of $k$-points exactly when
their restrictions to a common refinement agree as scheme morphisms,
because the pieces are reduced and of finite type over an algebraically
closed field.  The colimit is therefore
$\operatorname{Hom}_{\mathcal P_k}(X,Y)$.  The second displayed
bijection is the defining adjunction for $T(Y)$.

For products, test against an arbitrary $Z\in\operatorname{Abs}_k$:
\[
\begin{aligned}
 \operatorname{Hom}(Z,T(X\times_kY))
  &\simeq\operatorname{Hom}(Z,X\times_kY)\\
  &\simeq\operatorname{Hom}(Z,T(X))\times
          \operatorname{Hom}(Z,T(Y)).
\end{aligned}
\]
This proves the internal product formula.  Applying the same argument
to the ordinary product $T(X)\times_kT(Y)$ proves the second formula.
The terminal object is $\operatorname{Spec}k$, so this fully faithful
realization preserves the finite products used to define group objects.
\end{proof}

\paragraph{Why the ordinary product cannot be substituted.}
The fields $k(t)$ and $k(s)$ are absolutely flat rings, but
$k(t)\otimes_k k(s)$ has both the zero prime and the nonzero prime
defined by $t-s$.  Indeed it is the localization of $k[t,s]$ obtained
by inverting all nonzero polynomials in $t$ alone and in $s$ alone;
none of these meets the prime $(t-s)$.  Thus this tensor product is
not absolutely flat.  The same issue already occurs for
$T(\mathbb A^1_k)\times_k T(\mathbb A^1_k)$: its coordinate ring
has the displayed tensor product as a quotient, obtained by taking
the generic residue-field quotient on both factors.  A quotient of
an absolutely flat ring would be absolutely flat, so the ordinary
product cannot be absolutely flat in this example.

\begin{corollary}[Proved: algebraization in the absolutely flat category]\label{C:absolute-flat-group-algebraization}
With the internal product just specified, group objects in the full
subcategory formed by the $T(X)$ for separated finite-type reduced $k$-schemes
are equivalent to algebraic $k$-groups.  In characteristic zero this
also realizes group objects with arbitrary constructible graphs.
\end{corollary}

\begin{proof}
Apply Theorem~\ref{C:absolute-flat-piecewise} to the multiplication,
inverse and unit morphisms and to their commutative diagrams.  Full
faithfulness makes the group identities equivalent.  Then apply
Theorem~\ref{C:piecewise-group-algebraization}.
\end{proof}

\paragraph{Exact historical scope.}
The reference to EGA I, 7.2.14 is an explicit published identification
of the universal absolutely flat construction with Grothendieck's
constructible-topology scheme.  The theorem above proves its comparison
with the stated \emph{finite-type, reduced, finite-partition} category.
It does not by itself identify the complete category intended by the
1964 subject, or extend the algebraization assertion to all absolutely
flat schemes, pro-affine schemes, or arbitrary sheaves.

% C-written, B-owned labels. Prospective fragment, not included at creation.
\paragraph{A finite Greenberg comparison and its exact infinite-level limit.}
We now give a statement directly about the finite-partition construction,
with its hypotheses separated from the wider historical sheaf programme.
Throughout the following proposition $k$ is algebraically closed,
$R_m=k[t]/(t^m)$, and $m\geq1$.  For a separated reduced finite-type
$k$-scheme $Y$, let $P$ run over its finite partitions into affine locally
closed reduced subschemes and write
\[
 A_P=\prod_{S\in P}\Gamma(S,\mathcal O_S).
\]
Thus $Y_P=\operatorname{Spec}A_P=\coprod_{S\in P}S$.

\begin{proposition}[Proved under stated hypotheses: finite Greenberg comparison]\label{B:finite-Greenberg-piecewise}
Let $X$ be a separated finite-type $R_m$-scheme, not necessarily reduced.
Its finite Greenberg realization $\operatorname{Gr}_m(X)$ is a separated
finite-type $k$-scheme characterized by
\[
 \operatorname{Gr}_m(X)(A)=X(R_m\otimes_k A)
\]
for every $k$-algebra $A$.  Define the finite-partition presheaf by
\[
 N_m(X)(Y)=\varinjlim_P\prod_{S\in P}
 X\bigl(R_m\otimes_k\Gamma(S,\mathcal O_S)\bigr).
\]
There are natural bijections
\[
\begin{aligned}
 N_m(X)(Y)
 &\simeq\operatorname{Hom}_k\bigl(T(Y),\operatorname{Gr}_m(X)\bigr)\\
 &\simeq\operatorname{Hom}_{\operatorname{Abs}_k}
       \bigl(T(Y),T(\operatorname{Gr}_m(X))\bigr).
\end{aligned}
\]
In particular this is a realization of the finite-partition construction
in the absolutely flat category, using its internal products.
The Greenberg functor preserves products and group objects.
\end{proposition}

\begin{proof}
We first supply the elementary equal-characteristic construction.
For an affine $X$, its coordinate ring has a finite presentation
\[
 R_m[z_1,\ldots,z_d]/(f_1,\ldots,f_r).
\]
Replace every $z_i$ by $\sum_{0\leq j<m}z_{ij}t^j$ and expand every
$f_\ell$ modulo $t^m$.  The $m$ coefficient equations for each
$f_\ell$ are polynomial equations over $k$.  Their common zero scheme
represents the required functor on $k$-algebras.  The universal property
makes the construction independent of the chosen presentation.

It glues for a separated finite-type $X$.  Indeed, for every $k$-algebra
$A$, the nilpotent immersion
$\operatorname{Spec}A\to\operatorname{Spec}A[t]/(t^m)$ is a
homeomorphism on underlying spaces.  Consequently a morphism into $X$
factors through an open $U\subset X$ exactly when its reduction does.
Thus the realization of $U$ is the inverse image of $U_k$ under the
reduction morphism.  Applying this to a finite affine cover and its
intersections glues the affine realizations, with precisely the required
functor of points.  The same coefficient construction preserves closed
immersions.  Since the construction preserves products by its universal
property, the realization of the diagonal is closed, proving separatedness.
Finite type follows from the finite affine cover and the explicit finite
polynomial presentations.

This is the equal-characteristic Greenberg functor, equivalently Weil
restriction along $\operatorname{Spec}R_m\to\operatorname{Spec}k$.
The general finite-artinian construction and this identification are
recorded in Proposition--Definition~7.10 and Remarks~7.17(b)--(d) of the
corrected version of Bertapelle--Gonz\'alez-Avil\'es
\cite[pp.~50--54, especially (7.19)]{BBertapelleGonzalezAviles2022}.
Its adjunction is
\[
 \operatorname{Hom}_k(Z,\operatorname{Gr}_m(X))
 \simeq\operatorname{Hom}_{R_m}(Z\times_k\operatorname{Spec}R_m,X).
\]
This proves preservation of products and transports group operations
and their identities.

We have already proved $T(Y)=\varprojlim_PY_P$.  The target
$\operatorname{Gr}_m(X)$ is of finite presentation over $k$, so
finite-presentation descent and equality after refinement give
\[
 \operatorname{Hom}_k(T(Y),\operatorname{Gr}_m(X))
 =\varinjlim_P\operatorname{Hom}_k(Y_P,\operatorname{Gr}_m(X))
\]
by \cite[Lemma~32.10.1(2)--(3), tag 01ZM]{BStacks}.
The adjunction and the finite disjoint union defining $Y_P$ identify
this with $N_m(X)(Y)$.  The second bijection is the defining right
adjunction for $T$.  Although $\operatorname{Gr}_m(X)$ can be
nonreduced, every map from $T(Y)$ or $Y_P$ factors through its
reduction.  No reducedness of the Greenberg realization has been
silently assumed.
\end{proof}

\begin{proposition}[Refuted: interchange of uniform finite partitions with infinite levels]\label{B:uniform-partition-obstruction}
Put $R=k[[t]]$, $A=k[x]$, and let $P$ range over the same finite reduced
affine locally closed partitions of $\operatorname{Spec}A$.  The
natural homomorphism
\[
 D:=\varinjlim_P A_P[[t]]\longrightarrow T(A)[[t]]
\]
is injective but not surjective.  Nevertheless, for every $m\geq1$,
\[
 D/t^mD\simeq T(A)[t]/(t^m),\qquad
 \widehat D^{\,(t)}\simeq T(A)[[t]].
\]
In particular $D$ is a proper dense subring of its $t$-adic completion,
and
\[
 \varinjlim_P\varprojlim_m A_P[t]/(t^m)
 \;\longrightarrow\;
 \varprojlim_m\varinjlim_P A_P[t]/(t^m)
\]
is not surjective.  This is already the additive-group construction for
$\mathbb G_{a,R}$; it is not caused by singularities or non-affineness.
\end{proposition}

\begin{proof}
Each refinement homomorphism $A_P\to A_Q$ is injective: a function
vanishing on every piece of a finite partition vanishes at every prime,
and the original stratum is reduced.  Hence $A_P\to T(A)$ is
injective.  Taking coefficients proves that
$A_P[[t]]\to T(A)[[t]]$ is injective, and the asserted map from $D$
is injective as well.

Choose pairwise distinct $a_0,a_1,\ldots\in k$ and let
\[
 e_n=1-(x-a_n)(x-a_n)^{[-1]}\in T(A).
\]
Here the bracket denotes the unique pointwise inverse.  Thus $e_n$ is
the idempotent equal to $1$ at the closed point $a_n$ and $0$ at all
other primes.  Consider
\[
 z=\sum_{n\geq0}e_nt^n\in T(A)[[t]].
\]
If $z$ came from $A_P[[t]]$, the stratum of $P$ containing the generic
point of the affine line would be a dense open $U$.  Its complement
is finite, so $a_n\in U$ for some $n$.  The $n$th coefficient of the
representative would then be a regular function on $U$ equal to $e_n$.
It vanishes on the dense open $U\setminus\{a_n\}$ and therefore
vanishes identically, since $U$ is reduced.  This contradicts its value
$1$ at $a_n$.  Consequently $z$ is outside the image.

Every finite set of elements of $T(A)=\varinjlim_PA_P$ comes from
one common partition.  Thus $D\to T(A)[t]/(t^m)$ is surjective.
If a representative in $A_P[[t]]$ maps to zero modulo $t^m$, its first
$m$ coefficients are already zero in $A_P$ by injectivity; it is
therefore $t^m$ times another series in $A_P[[t]]$.  The kernel is
exactly $t^mD$.  The asserted quotient and completion formulas follow.
For the specific series $z$, every finite truncation is represented
by the partition isolating the finitely many points involved.  What
fails is the existence of one finite partition for all coefficients.
\end{proof}

\paragraph{Presheaves, sheafification, and the historical scope.}
The last proposition is an exact statement about the displayed
finite-partition presheaves.  It must not be silently promoted to a
statement about a different site or its sheafification.  Such a
promotion requires checking which covers are allowed and whether a
cover can make the series admit one uniform finite partition.  In
particular it does not identify the entire historical target
$\mathcal C'(k)$ or resolve a crossed-out assertion about injective
sheaves.  The finite theorem and the obstruction concern precisely
the stated equal-characteristic, finite-type test setting.

The issue is consistent with, but distinct from, the infinite-level
localization warning in the modern source:
\cite[Proposition~18.2 and Remark~18.3, pp.~99--100]{BBertapelleGonzalezAviles2022}
corrects an earlier proof which had used a finite-level localization
identity at infinite Witt level.  The infinite point formula there
remains true by a different algebraization theorem.  We use neither
that faulty proof nor an unproved interchange of limits.

\begin{proposition}[Refuted: surjectivity after finite-surjective sheafification]\label{B:finite-cover-partition-obstruction}
Let $k$ be algebraically closed.  On the category of reduced affine
finite-type $k$-schemes, equip fibre products with their reduced scheme
structures and take finite surjective morphisms as covers.  For
$Y=\operatorname{Spec}B$ define the presheaves of additive abelian groups
\[
 D(Y)=\varinjlim_P B_P[[t]],\qquad E(Y)=T(B)[[t]],
\]
where $P$ runs over finite affine locally closed reduced partitions of
$Y$.  The induced morphism of associated \emph{sheaves of abelian groups}
\[
 D^\#\longrightarrow E^\#
\]
is not an epimorphism.  The same is true for the associated sheaves of
sets.  In particular the series $z$ of
Proposition~\ref{B:uniform-partition-obstruction} has no local lift from
$D^\#$ on this specified site.  No conclusion about categorical
epimorphisms of sheaves of rings is asserted.
\end{proposition}

\begin{proof}
Finite surjective maps compose and remain finite surjective under base
change and reduction, so these maps form a pretopology.  A finite
jointly surjective family can be replaced by its disjoint union.
For a morphism between affine test schemes, inverse images of affine
locally closed strata are affine locally closed; taking reductions
therefore defines the indicated pullback homomorphisms and presheaves.
The refinement maps are injective by the reducedness argument already
given, so $D\to E$ is injective.

Moreover $E$ is separated for this topology.  If
$\operatorname{Spec}B\to\operatorname{Spec}A$ is finite surjective,
then $T(A)\to T(B)$ is injective.  For each prime $\mathfrak p$ of
$A$ choose a prime $\mathfrak q$ of $B$ above it.  The map
$\kappa(\mathfrak p)\to\kappa(\mathfrak q)$ is injective, and these
are also the corresponding residue fields of $T(A)$ and $T(B)$.
An element in the kernel thus vanishes in every residue field of the
reduced ring $T(A)$, hence is zero.  This argument applies coefficient
by coefficient to $E$, proving separatedness.  Consequently
$E\to E^\#$ is injective.  The same holds for $D$, since $D\subset E$.

We claim that for every finite surjective
$f:Y\to\mathbb A^1_k$, the pullback $f^*z$ lies outside $D(Y)$.
Choose a reduced irreducible component $C\subset Y$ dominating
$\mathbb A^1_k$.  It exists by surjectivity and finiteness.  The
restriction $C\to\mathbb A^1_k$ is finite and dominant, hence
surjective, and $C$ is an integral affine curve.

Suppose that $f^*z$ is represented on a finite partition $P$ of $Y$.
Let $S\in P$ contain the generic point of $C$.  The reduced intersection
$U=C\cap S$ is a dense open in $C$.  Since $C$ is a noetherian curve,
$C\setminus U$ is a finite set of closed points, and its image under
$f$ is finite.  Choose $a_n$ outside that image and a point $y\in C$
over $a_n$; then $y\in U$.  The coefficient of $t^n$ on the stratum
$S$, restricted to $U$, is regular.  Its value at the generic point of
$C$ is zero, since that is the value of $f^*e_n$.  It therefore
vanishes identically on the integral scheme $U$.  At $y$, however,
its residue equals $1$, a contradiction.  This proves the claim.

For completeness we pass from this pointwise claim to the precise
sheaf-theoretic conclusion.  A section of an associated sheaf locally
comes from the presheaf; this follows from the plus construction,
\cite[Lemma~7.10.8 and Theorem~7.10.10]{BStacks}.  If the image of $z$
in $E^\#(\mathbb A^1_k)$ had a local lift to $D^\#$, refining that
cover and composing covers would produce one finite surjective
$f:Y\to\mathbb A^1_k$ and $d\in D(Y)$ whose image is $f^*z$ in
$E^\#(Y)$.  Injectivity of $E(Y)\to E^\#(Y)$ would make the equality
hold already in $E(Y)$, contradicting the claim.  Thus the morphism
is not locally surjective.  This is equivalent to failure of
surjectivity for sheaves of sets, by
\cite[Definition~7.11.1 and Lemma~7.11.2]{BStacks}, and to failure of
epimorphy for sheaves of abelian groups.  The latter can also be seen
from the nonzero associated-sheaf cokernel: the class of $z$ cannot
vanish locally, precisely by the same argument.
\end{proof}

\paragraph{Scope after sheafification.}
This last proposition supplies the additional cover analysis for the
explicit ordinary finite-type site just defined.  It does not assert
that this is every test object or covering allowed in the historical
$\mathcal C'(k)$, or that Grothendieck asserted the failed interchange.
The original phrase specifying finite surjective descent is relevant
motivation, while the test-object adjectives and any larger
pro-affine or perfect setting must be checked separately.

% Root derivation; prospective B-owned label, not included at creation.
\paragraph{The complete Greenberg functor on its stated test rings.}
The complete construction also has a precise modern theorem.  We specify
the test rings because this theorem does not interchange a finite
partition colimit with an inverse limit.

\begin{theorem}[Known, with explicit derivation: complete Greenberg points]\label{B:complete-Greenberg-points}
Let $R$ be a complete discrete valuation ring with perfect residue field
$k$ and uniformizer $\pi$.  In equal characteristic fix an isomorphism
$R\simeq k[[t]]$ carrying $\pi$ to $t$.  In mixed characteristic use the
canonical Cohen inclusion $W(k)\hookrightarrow R$.  Put $R_n=R/(\pi^n)$
and let $\mathcal R_n$ be its finite Greenberg ring scheme over $k$.
For a $k$-algebra $A$, assume that $A$ is perfect in mixed characteristic,
and define
\[
 S(A)=\varprojlim_n\mathcal R_n(A).
\]
For every $R$-scheme $X$, without a finite-type or separatedness assumption,
the scheme
\[
 \operatorname{Gr}_R(X)=
 \varprojlim_n\operatorname{Gr}_{R_n}(X\times_R R_n)
\]
exists.  There are natural bijections
\[
 \operatorname{Gr}_R(X)(A)
 \simeq\varprojlim_n X\bigl(\mathcal R_n(A)\bigr)
 \simeq X\bigl(S(A)\bigr),
\]
where all points on the right are morphisms over $R$.  Here
\[
 S(A)\simeq
 \begin{cases}
 A[[t]],&R\simeq k[[t]],\\
 R\otimes_{W(k)}W(A),&\operatorname{char}R=0,\ \operatorname{char}k>0.
 \end{cases}
\]
The construction preserves products and group objects.  It is generally
not of finite type: for $X=\mathbb A^1_R$ its underlying scheme is
$\operatorname{Spec}k[x_0,x_1,\ldots]$.
\end{theorem}

\begin{proof}[Derivation from the opened primary theorems]
The finite ring schemes and their transition maps are those of
\cite[Sections~4,6--7]{BBertapelleGonzalezAviles2022}.
The transition morphisms in the displayed system of Greenberg schemes
are affine; Proposition--Definition~17.5 and its proof give the inverse
limit and its functor of points
\cite[pp.~93--94; see also (18.1)]{BBertapelleGonzalezAviles2022}.
This proves the first bijection using the finite-level point formula
(7.16).  No exchange with a filtered colimit has been made.

In equal characteristic, the ring formula and
$S(A)/(t^n)=\mathcal R_n(A)$ are immediate from
$\mathcal R_n(A)=A[t]/(t^n)$.  In mixed characteristic,
\cite[Remark~6.6(b), (6.7)]{BBertapelleGonzalezAviles2022}
identifies $S(A)$ with $R\otimes_{W(k)}W(A)$ and its quotient by
$\pi^n$ with $\mathcal R_n(A)$.  We check completeness, which is the
essential hypothesis of the next theorem.  Write
$R=W(k)[T]/(f)$ for an Eisenstein polynomial $f$ of degree $e$,
using the chosen uniformizer.  Then
$S(A)=W(A)[T]/(f)$ is finite free over the $p$-adically complete ring
$W(A)$; here perfection of $A$ gives
$W(A)/(p^n)=W_n(A)$.  It is consequently $p$-adically complete.
The Eisenstein relation gives $\pi^e=p u$ with $u$ a unit:
the relevant factor has a unit constant term and all its other terms
are divisible by $\pi$, which is topologically nilpotent since
$\pi^e\in pS(A)$.  Thus the $p$-adic and $\pi$-adic topologies are
cofinal.  We obtain in both cases
\[
 S(A)\simeq\varprojlim_n S(A)/(\pi^n),
 \qquad S(A)/(\pi^n)\simeq\mathcal R_n(A).
\]

Bhatt's Theorem~4.1 gives
$X(S(A))\simeq\varprojlim_n X(S(A)/(\pi^n))$ for qcqs algebraic
spaces; his Remark~4.6 extends it to every scheme
\cite[Section~4, pp.~16--17]{ROOT_BhattAlgebraization2014}.
Apply that theorem to the underlying scheme of $X$ and the complete
ring $S(A)$.  The bijection respects the condition of being over $R$:
if the compatible maps at every level are over $R$, the structural map
$R\to S(A)$ obtained after algebraization agrees with the prescribed
one modulo every $\pi^n$, hence agrees with it because $S(A)$ is
separated.  This proves the second bijection, including its relative
typing.  This is the corrected justification of
\cite[Proposition~18.2]{BBertapelleGonzalezAviles2022}.

The finite adjunctions preserve products, as do inverse limits; applying
these identities to multiplication, inverse, identity and their diagrams
gives the assertion about group objects.  Finally, the underlying scheme
of $\mathcal R_n$ is affine $n$-space with the compatible coordinate
projections of the construction, so its inverse limit is the infinite
affine scheme displayed above, as also recorded in
\cite[Section~6 and the paragraph after (18.1)]{BBertapelleGonzalezAviles2022}.
\end{proof}

% Prospective B-owned label.  Original source convention and complete proof checked.
\begin{lemma}[Proved under stated hypotheses: a Witt-localization witness]\label{B:Witt-localization-witness}
Use standard $p$-typical Witt coordinates and write $[f]=(f,0,0,\ldots)$.
For every commutative ring $A$, every $f\in A$, every $N\geq0$, and every
$w=(w_i)_{i\geq0}\in W(A)$, one has
\[
 ([f]^Nw)_i=f^{Np^i}w_i.                                      \tag{\(*\)}
\]
Consequently, for $A=\mathbb F_p[f]$, with $f$ an indeterminate, the
canonical map
\[
 W(A)_{[f]}\longrightarrow W(A_f)
\]
is injective and is not surjective.  An explicit element outside its
image has coordinates $v_i=f^{-ip^i}$.  The same conclusions and witness
hold for the perfect domain
$A=\mathbb F_p[f^{1/p^r}:r\geq0]$.

In contrast, the vector with coordinates $w_i=f^{-i}$ does lie in
$W(A)_{[f]}$.  Thus this particular illustrative vector in
\cite[Remark~18.3]{BBertapelleGonzalezAviles2022} must be replaced when
using the standard-coordinate convention explicitly fixed in that
source, Section~2.3, pp.~8--9.  The localization inequality and the
corrected Greenberg theorem remain valid.
\end{lemma}

\begin{proof}
Work first over the universal torsion-free polynomial ring
$B=\mathbb Z[f,w_0,w_1,\ldots]$.  The $n$th ghost component of the
vector on the right of $(*)$ is
\begin{align*}
 \sum_{i=0}^n p^i(f^{Np^i}w_i)^{p^{n-i}}
 &=f^{Np^n}\sum_{i=0}^n p^i w_i^{p^{n-i}}\\
 &=\operatorname{gh}_n([f]^N)\operatorname{gh}_n(w).
\end{align*}
The ghost map is injective over $B$: induction on $n$, using its
triangular formula and the absence of $p$-torsion, recovers the $n$th
Witt coordinate from the first $n+1$ ghost components.  Hence $(*)$
holds over $B$.  Since Witt operations are universal integral
polynomials, specialization proves $(*)$ over every ring, including
rings of characteristic $p$ where the ghost map itself need not be
injective.

For either domain $A$ in the statement, the injection $A\hookrightarrow
A_f$ gives a coordinatewise injection $W(A)\hookrightarrow W(A_f)$.
The Teichm\"uller vector $[f]$ becomes a unit, with inverse $[f^{-1}]$.
Formula $(*)$ also shows that multiplication by $[f]$ on $W(A)$ is
injective.  Thus localization embeds $W(A)_{[f]}$ in $W(A_f)$.
If $v=(f^{-ip^i})_{i\geq0}$ belonged to its image, there would be an
$N\geq0$ such that $[f]^Nv\in W(A)$.  Its $i$th coordinate would be
\[
 f^{(N-i)p^i}.
\]
Taking $i>N$ gives a negative power of $f$, which belongs to neither
$\mathbb F_p[f]$ nor its stated perfection.  This proves non-surjectivity.

Finally, for $w_i=f^{-i}$, formula $(*)$ gives
$([f]w)_i=f^{p^i-i}$.  Since $p\geq2$ implies $p^i\geq i$ for every
$i\geq0$, all these coordinates belong to $A$.  Hence
$w=[f]^{-1}([f]w)$ is in the localized Witt ring.  This directly checks
the source's displayed tuple; no assertion of priority is made.
\end{proof}

The accompanying exact certificate
\texttt{code/C\_Witt\_localization\_certificate.py} records 147 passing
checks in \texttt{logs/C\_Witt\_localization\_certificate.json}
(elapsed time $0.06913506400451297$ seconds).  Its finite checks
supplement the all-coordinate proof above.

\paragraph{A precise perfect-test relative-cohomology theorem.}
The following gives a modern derived construction directly relevant to
the relative-cohomology programme on cote~43, folios~95--96.
Its sites and comparison hypotheses are specified independently; we
do not identify them with an incompletely reconstructed historical site.

\begin{theorem}[Known: perfect Greenberg realization and relative cohomology]\label{B:perfect-Greenberg-cohomology}
Let $R$ be a complete discrete valuation ring with fraction field $K$,
uniformizer $\pi$, and perfect residue field $k$ of characteristic $p>0$.
Equal and mixed characteristic are both allowed. Fix a coefficient map
$W(k)\to R$. For a finite product $L$ of perfect field extensions of $k$,
put
\[
 R(L)=\varprojlim_n
 \bigl(W_n(L)\otimes_{W_n(k)}R/\mathfrak m_R^n\bigr),
 \qquad K(L)=R(L)[1/\pi].
\]
For $R=k[[t]]$ with its specified coefficient field, these are $L[[t]]$
and $L((t))$. For $L$ a field, $R(L)$ is a complete discrete valuation
ring with residue field $L$.

Let $S_{\mathrm{pa}}$ be the étale site on finite products of perfect
field extensions of $k$, $S_{\mathrm{pf}}$ the pro-fppf site on
quasi-compact quasi-separated perfect $k$-schemes, and
$S_{\mathrm{ir}}$ the pro-étale site on ind-rational $k$-algebras.
Here rational algebras are finite products of perfections of finitely
generated field extensions, and ind-rational algebras are filtered
unions of rational subalgebras. In the pro-fppf site, coverings are
finite jointly surjective families flat of profinite presentation in
the perfect sense. The inclusions induce premorphisms
\[
 h:S_{\mathrm{pf}}\to S_{\mathrm{pa}},\qquad
 \widetilde h:S_{\mathrm{pf}}\to S_{\mathrm{ir}},\qquad
 \mathsf a=R\widetilde h_*Lh^*:
 D(S_{\mathrm{pa}},\mathbf{Ab})\to D(S_{\mathrm{ir}},\mathbf{Ab}).
\]
These premorphisms are not assumed to have exact inverse image on
arbitrary sheaves of sets.
\begin{enumerate}
\item Let $\mathcal P'_{\mathrm{fc}}$ consist of commutative
proalgebraic groups representable as extensions of perfections of
abelian varieties by perfect affine groups, with finite étale component
group. Restriction to $S_{\mathrm{pa}}$ is fully faithful on this
category, and its derived counit is an isomorphism:
\[
 Lh^*(P|_{S_{\mathrm{pa}}})\simeq P
 \quad\text{in }D(S_{\mathrm{pf}},\mathbf{Ab}).
\]
\item The functors $R(-),K(-)$ induce premorphisms
$\pi_R:R_{\mathrm{fppf}}\to S_{\mathrm{pa}}$ and
$\pi_K:K_{\mathrm{fppf}}\to S_{\mathrm{pa}}$. Define
\[
 \mathcal R_R(F)=\mathsf aR\pi_{R,*}F,\qquad
 \mathcal R_K(F)=\mathsf aR\pi_{K,*}F,\qquad
 \mathcal R_x(F)=
 \operatorname{fib}\bigl(\mathcal R_R(F)\to\mathcal R_K(F|_K)\bigr).
\]
They give a functorial localization triangle.
Call $Q\in D(S_{\mathrm{pa}},\mathbf{Ab})$ $h$-acyclic if its unit
$Q\to Rh_*Lh^*Q$ is an isomorphism. If $R\pi_{R,*}F$ is
$h$-acyclic, then
\[
 R\Gamma(L_{\mathrm{pro\acute et}},\mathcal R_R(F))
 \simeq R\Gamma(R(L)_{\mathrm{fppf}},F).
\]
The analogous assertion holds over $K$. If both unit conditions hold
for $F$ and $F|_K$, the analogous comparison holds with support on
the closed fibre. Under this unit hypothesis, if both $F$ and
$\mathcal R_R(F)$ are bounded below, there is the spectral sequence
\[
 E_2^{ij}=
 H^i(L_{\mathrm{pro\acute et}},\mathcal H^j\mathcal R_R(F))
 \Longrightarrow H^{i+j}(R(L)_{\mathrm{fppf}},F).
\]
For algebraically closed $L$, this reduces to
$\mathcal H^q\mathcal R_R(F)(L)=H^q(R(L)_{\mathrm{fppf}},F)$.
\item If $G/R$ is a smooth commutative group scheme of finite type,
write $G_n=G\times_RR/\mathfrak m_R^n$ and
\[
 \mathbf G=\varprojlim_n
 \bigl(\operatorname{Gr}_{R/\mathfrak m_R^n}(G_n)\bigr)^{\mathrm{perf}}.
\]
Then $\mathbf G\in\mathcal P'_{\mathrm{fc}}$ and
\[
 0\longrightarrow U_G\longrightarrow\mathbf G
 \longrightarrow G_k^{\mathrm{perf}}\longrightarrow0
\]
is exact, with $U_G$ connected pro-unipotent.
Consequently $\pi_0(\mathbf G)=\pi_0(G_k)$, and
$\mathcal R_R(G)\simeq\mathbf G[0]$.
The unit condition in the preceding clause holds.
\item If $N/R$ is finite flat and commutative, then
$\mathcal H^0\mathcal R_R(N)$ is finite étale and agrees with
$\mathcal H^0\mathcal R_K(N_K)$,
$\mathcal H^1\mathcal R_R(N)$ is connected pro-unipotent, and
$\mathcal H^q\mathcal R_R(N)=0$ for $q\geq2$.
The object $R\pi_{R,*}N$ is $h$-acyclic.
\end{enumerate}
\end{theorem}

\begin{proof}[Derivation from the opened sources]
The sites, ring functors and premorphisms are those of
\cite[Definitions~3.2, 3.4, 3.7--3.8, 7.4--7.7, 8.1--8.4;
Proposition~8.3]{BSuzuki2021}; all locators here refer to its opened
7~October~2019 author version. Finite étale residue-field extensions
lift to finite étale extensions of Witt rings. Tensoring and completing
with $R$ preserves these finite free extensions and their covering
fibre products. This proves the asserted premorphism property, rather
than asserting exactness of an arbitrary inverse-image functor.

For the first clause, the specified category is Definitions~3.12--3.13
of that source. Theorem~3.15 applies to the larger category $E_k$, and
therefore gives the displayed derived counit. Its degree-zero part
and adjunction give, for $P,Q\in\mathcal P'_{\mathrm{fc}}$,
\[
 \operatorname{Hom}_{S_{\mathrm{pa}}}(P,Q)
 =\operatorname{Hom}_{S_{\mathrm{pf}}}(h^*P,Q)
 =\operatorname{Hom}_{S_{\mathrm{pf}}}(P,Q).
\]
This is precisely Proposition~7.1, or Theorem~1.1.
The source proves the derived assertion in its complete
Sections~4--6 using the Mac Lane resolution and pro-fppf covers of
generic points. Proposition~3.9 explicitly rules out replacing that
argument by exactness on arbitrary sheaves of sets.

For the second clause, $\mathsf a$ and the relative functors are
Definitions~7.7 and~8.4. Derived functors preserve the defining fibre
triangle, giving (8.1). Factor
$h=\alpha\epsilon\widetilde h$ as in Definition~7.5, with
$\epsilon$ the change from pro-étale to étale topology on ind-rational
algebras and $\alpha_*$ exact. The \emph{unit} hypothesis gives
\[
 R\pi_{R,*}F
 \simeq Rh_*Lh^*R\pi_{R,*}F
 \simeq\alpha_*R\epsilon_*\mathcal R_R(F).
\]
Derived sections at $L$ give the claimed comparison, exactly as in
Proposition~8.5. The displayed spectral sequence is the bounded-below hypercohomology
spectral sequence of the relative object. We explicitly retain a lower
bound on that object; no general t-exactness of $Lh^*$ is inferred.
This extra bound is automatic in the third and fourth clauses. Pro-étale covers of an algebraically
closed field split, so sections there are exact. Applying the argument
over $K$ and then taking the localization fibre gives the supported
version. The counit of the first clause and this unit are different
maps; the second assertion does not discard its acyclicity hypothesis.

For the third clause, finite Greenberg representability identifies
the points of its $n$th term on perfect residue-field extensions with
$G(R(L)/\mathfrak m_R^n)$
\cite[Proposition--Definition~7.10, (7.16)]{BBertapelleGonzalezAviles2022}.
Smoothness gives surjectivity of reduction. The successive transition
kernels become vector groups after perfection, so the kernel of the
map to $G_k^{\mathrm{perf}}$ is connected pro-unipotent. The first
stage has finite component group because $G_k$ is of finite type.
These are exactly the argument and conclusions of
\cite[Proposition~(3.4.2)(a), pp.~48--49]{BSuzuki2020},
whose locators refer to the opened 22~March~2018 author version.
The infinite point formula follows by taking the complete compatible
system of points; the general algebraization justification is
\cite[Proposition~18.2]{BBertapelleGonzalezAviles2022}, using
\cite[Theorem~4.1, Remarks~4.3 and~4.6]{ROOT_BhattAlgebraization2014}.

Smooth fppf cohomology agrees with étale cohomology. The exactness
statement in \cite[Proposition~(3.4.1)]{BSuzuki2020} says that the
relative étale direct image has no higher cohomology: étale covers of
the complete valuation-ring tests are locally refined by lifted
residue-field covers. Thus the old relative object is
$\mathbf G[0]$. Propositions~10.8--10.9 of
\cite{BSuzuki2021} identify it with the new functor above and give
the unit condition. Proposition~7.10(b) also verifies that condition
for the tower with connected unipotent transition kernels. The
vanishing is a statement about relative cohomology \emph{sheaves};
ordinary cohomology for non-algebraically-closed $L$ is still computed
by the spectral sequence in the second clause.

For the fourth clause, the standard smooth resolution
\cite[Proposition~2.2.1, pp.~25--27]{BBegueri1980}, applied to the
Cartier dual of $N$, gives
\[
 0\longrightarrow N\longrightarrow H\longrightarrow J\longrightarrow0
\]
with $H,J$ smooth affine commutative and with connected fibres.
The first smooth group is a Weil restriction of $\mathbb G_m$ from
the finite flat Cartier dual; its geometric fibres are unit groups of
finite algebras and are connected. Its quotient by the finite
subgroup is connected as well. Applying the third clause gives
\[
 0\to\mathcal H^0\mathcal R_R(N)\to\mathbf H\to\mathbf J
 \to\mathcal H^1\mathcal R_R(N)\to0
\]
and vanishing in degrees at least two.
Proposition~(3.2.10) of \cite{BSuzuki2020} identifies the degree-zero
object with the finite étale group defined by $N(K^{\mathrm{ur}})$
and its residue-field Galois action. Since $N$ itself is the locally
finite part, the $R$ and $K$ point functors agree.
The cokernel is connected proalgebraic and is annihilated by the
order of $N$. Its semiabelian part is therefore zero, so it is
pro-unipotent. These are the complete steps of
\cite[Proposition~(3.4.2)(b)]{BSuzuki2020}. The comparison and the
unit condition are again Propositions~10.8--10.9 of \cite{BSuzuki2021}.
\end{proof}

\paragraph{The perfection qualification and the original programme.}
Bégueri's own Section~3.1 fixes mixed characteristic and an
algebraically closed residue field. Its Section~4.1 and the remarks
after Proposition~6.2.3, pp.~76--77, construct perfect Greenberg
realizations and a derived functor on groups of multiplicative type.
They provide a direct precursor for the present relative construction.
The nonperfect kernel formula of Lemma~4.1.1(2) is nevertheless false:
for $R=W(k)$ and $G=\mathbb G_a$, the proposed kernel map has kernel
$\alpha_p$
\cite[Remark~14.22]{BBertapelleGonzalezAviles2022}.
Perfection removes that defect, as Remark~15.9 explains and as
\cite[footnote~13, p.~48]{BSuzuki2020} explicitly requires.
We use perfect groups throughout. Neither this theorem nor the
finite-partition comparison identifies every historical test object
or proves the crossed-out injective-preservation argument on
cote~43/96. Those are scope limits on the historical interpretation,
not silently assumed steps in the proved modern theorem.


\subsubsection{12. Globalization and extensions of N\'eron models}
\begin{theorem}[Known: the Dedekind-base abelian theorem]\label{B:subject12}
Let $S$ be a connected Dedekind scheme, $K$ its function field and $A/K$ an abelian variety. There is a smooth separated finite-type $S$-group scheme $\mathcal A$ with generic fibre $A$, satisfying
\[
 \operatorname{Hom}_S(T,\mathcal A)=\operatorname{Hom}_K(T_K,A)
\]
for every smooth $S$-scheme $T$. On an open set on which $A$ extends to an abelian scheme, that abelian scheme is its N\'eron model.
\end{theorem}
\begin{proof}[Derivation]
This is \cite[\S1.4, Theorem 3]{BBLR1990}. The local models of subject~11 are uniquely identified on their generic overlaps by the mapping property, hence glue with the abelian scheme on its good-reduction open. Uniqueness makes the gluing cocycles automatic. The proof in the book supplies existence and finite type, not merely uniqueness. This is the precise extension from a valuation ring to a one-dimensional base that the catalogue uses. A higher-dimensional base or an arbitrary smooth algebraic group changes the existence problem and is not included by the word ``extensions'' alone.
\end{proof}

\subsubsection{13. Greenberg asymptotics and approximation}
\begin{proposition}[Known, with the approximation bounds made explicit]\label{B:subject13}
Let $R$ be a henselian discrete valuation ring with uniformizer $t$, and assume $\widehat K/K$ is separable if the fraction field $K$ has positive characteristic. For a finite-type $R$-scheme $X$ there is a linear function $\beta(n)=an+b$ such that any point modulo $t^{\beta(n)}$ has the same reduction modulo $t^n$ as some point of $X(R)$. In particular
\[
 X(R)\ne\varnothing\quad\Longleftrightarrow\quad
 X(R/t^n)\ne\varnothing\text{ for every }n.
\]
For an excellent henselian noetherian local ring $(R,\mathfrak m)$, every finite system of polynomial equations over $R$ with a solution in $\widehat R$ has solutions in $R$ congruent to it modulo any prescribed $\mathfrak m^n$.
\end{proposition}
\begin{proof}[Derivation]
Greenberg \cite[Theorem 1 and Corollaries 1--2]{BGreenberg1966} gives constants $N,c,s$ so that a solution of the equations modulo $t^\nu$, $\nu\ge N$, is congruent to an exact solution modulo $t^{\lfloor\nu/c\rfloor-s}$. Taking $\nu\ge\max(N,c(n+s+1))$ gives the first assertion; increasing the constant term gives one linear $\beta$ valid for all $n$. Applying it for one $n$ proves the nonemptiness implication, without requiring the approximate points to form a compatible inverse system. The converse is reduction of an exact point.

For the second assertion, excellence makes $R\to\widehat R$ regular. Popescu's theorem, in the opened proof at Stacks Theorem 16.12.1, tag 07GC, expresses $\widehat R$ as a filtered colimit of smooth $R$-algebras. The algebra defined by the finitely many equations is finitely presented, so its solution map into $\widehat R$ factors through one such smooth algebra $B$. The resulting map $B\to\widehat R/\mathfrak m^n=R/\mathfrak m^n$ lifts to $B\to R$: smoothness gives an \'etale-local system of coordinates at the residue point, the coordinates lift to $R$, and the resulting \'etale section lifts by henselianity. Composing with the original finitely presented algebra gives the desired exact solution with the prescribed congruence. A general henselian noetherian ring need not have a regular completion map; that hypothesis is not suppressed.
\end{proof}

\begin{theorem}[Known: the exact noetherian local approximation criterion]\label{B:strong-approximation}
For a noetherian local ring $(R,\mathfrak m)$, the following three conditions are equivalent: $R$ is excellent and henselian; $R$ has the approximation property; and $R$ has the strong approximation property. Here approximation means that every solution in $\widehat R$ of a finite polynomial system over $R$ can be approximated to every finite order by a solution in $R$. Strong approximation means that, for each $f=(f_1,\ldots,f_r)\in R[Y_1,\ldots,Y_N]^r$, there is a function $\beta:\mathbb N_{>0}\to\mathbb N_{>0}$ such that
\[
 f(a)\in(\mathfrak m^{\beta(n)})^r
 \quad\Longrightarrow\quad
 \exists b\in R^N:\ f(b)=0,\quad b-a\in(\mathfrak m^n)^N.
\]
In particular every complete noetherian local ring has strong approximation. The function need not admit a linear bound, even for $R=\mathbb Q[[s,t]]$.
\end{theorem}
\begin{proof}[Derivation]
The complete-ring assertion is the Pfister--Popescu strong approximation theorem, in the full statement of \cite[Theorem 9]{BPopescu2017}; it does not assume regularity or equicharacteristic. For a ring with approximation, apply that complete-ring theorem to an approximate solution viewed in $\widehat R$, and then approximate the resulting exact formal solution by one in $R$. The congruence descends because $\mathfrak m^n\widehat R\cap R=\mathfrak m^n$. Conversely, enlarge $\beta$ so that $\beta(n)\ge n$. Given a formal solution and $n$, choose an element of $R^N$ with the same reduction modulo $\mathfrak m^{\beta(n)}$. Strong approximation gives the desired exact solution modulo $\mathfrak m^n$.

For necessity of the ring-theoretic condition, a simple root modulo $\mathfrak m$ of a monic polynomial lifts to $\widehat R$, and approximation modulo $\mathfrak m$ supplies the required root in $R$. This is the simple-root criterion for henselianity. Rotthaus \cite[Theorem 1, p.~455; proof \S3, pp.~464--466]{BRotthaus1990} proves that every noetherian semilocal ring with the approximation property is excellent, with no characteristic restriction. Her Definition~1.3 uses existence of an actual solution whenever a formal solution exists. This is equivalent to our approximation formulation: to impose congruence to $a$ modulo $\mathfrak m^n=(c_1,\ldots,c_e)$, adjoin the equations $Y_i-a_i=\sum_jc_jZ_{ij}$ and apply that existence property to the enlarged system. Finally, Proposition~\ref{B:subject13} gives approximation for excellent henselian rings. This proves all three equivalences. The next proposition proves the failure of linearity.
\end{proof}

\begin{proposition}[Refuted: linear approximation over all complete local rings]\label{B:nonlinear-Artin}
In $R=\mathbb Q[[s,t]]$, with maximal ideal $\mathfrak m=(s,t)$, put $f(U,V,W)=U^2-V^2W$. For every integer $m\ge5$, set
\[
 \begin{split}
 u_m&=\sum_{j=0}^{m-1}{\binom{1/2}{j}}
                 s^{2m-2-2j}t^{(m-2)j},\\
 v_m&=s^{2m-3},\qquad w_m=s^2+t^{m-2}.
 \end{split}
\]
Then $\operatorname{ord}_{\mathfrak m}f(u_m,v_m,w_m)=m^2-4$, but no exact zero of $f$ is congruent to $(u_m,v_m,w_m)$ modulo $\mathfrak m^{2m-2}$. Hence every function $\beta$ in Theorem~\ref{B:strong-approximation} for this $f$ satisfies
\[
 \beta(2m-2)>m^2-4\qquad(m\ge5),
\]
so no linear bound exists.
\end{proposition}
\begin{proof}
This is Rond's counterexample \cite[Theorem 1.2 and \S3]{BRond2005}, with parameters specialized and the error order made explicit. Write $p=m-2$ and $H_{m-1}(T)=\sum_{j=0}^{m-1}\binom{1/2}{j}T^j$. The first nonzero term of $H_{m-1}(T)^2-(1+T)$ has degree $m$ and coefficient $-2\binom{1/2}{m}\ne0$. We have
\[
 f(u_m,v_m,w_m)=s^{4m-4}
       \bigl[H_{m-1}(t^p/s^2)^2-1-t^p/s^2\bigr].
\]
Every displayed monomial has nonnegative exponents, and the first has total degree $4m-4+m(p-2)=m^2-4$. Subsequent terms have strictly larger total degree since $p>2$. This proves the error order.

We next show that $\operatorname{ord}_{\mathfrak m}(s^2+t^p-z^2)\le p$ for every $z\in R$. If the order were greater than $p$, comparison of quadratic terms would force the linear initial form of $z$ to be $s$ or $-s$; changing its sign, write $z=s+h$ with $\operatorname{ord}h\ge2$. If the lowest homogeneous part of $h$ has degree $d<p-1$, the term $2sh_d$ has degree $d+1<p$ and cannot cancel either $h^2$ or $t^p$. Hence $h$ has order at least $p-1$. But then its degree-$p$ contribution $2sh_{p-1}$ is divisible by $s$, and cannot equal $t^p$. This contradiction proves the bound, attained by $z=s$.

For an exact solution $(u,v,w)$ with $v\ne0$, we have $w=(u/v)^2$ in the fraction field. Since the regular local ring $R$ is integrally closed and $u/v$ solves the monic equation $T^2-w=0$, it lies in $R$. Thus $w$ is a square, and the preceding bound prevents congruence even modulo $\mathfrak m^{p+1}$. If $v=0$, the equation gives $u=0$, and $\operatorname{ord}(v_m-v)=2m-3$. These two cases exclude congruence modulo $\mathfrak m^{2m-2}$. If $\beta(2m-2)\le m^2-4$, the constructed approximate solution would contradict its defining property. Exact rational polynomial checks for $5\le m\le12$ accompany the symbolic proof, not a claim that those finite checks prove the infinite family.
\end{proof}

\subsubsection{14. Points of topoi and algebraic presentations of stacks}
\begin{proposition}[Known: two precise representability endpoints]\label{B:subject14}
A topos admitting a site with finite limits and with a finite subcover refining every covering has enough points. Also, if $R\rightrightarrows U$ is a smooth groupoid in algebraic spaces, its quotient stack $[U/R]$ is algebraic. Conversely an algebraic stack with a smooth atlas $U$ has the presentation $[U/(U\times_{\mathcal X}U)]$.
\end{proposition}
\begin{proof}[Derivation]
The first assertion is Deligne's enough-points theorem, in the opened proof of Stacks Theorem 7.39.3, tag 00YQ, tracing to SGA~4, Expos\'e VI, Appendix, Proposition 9.0. A finite-limit site with finite covering refinements defines a coherent topos, exactly the theorem's hypothesis; a conservative family of its point functors is the asserted family of fibre functors.

For the second, Stacks Theorem 94.17.3, tag 04TK, proves representability of the diagonal of $[U/R]$ and that $U\to[U/R]$ is smooth and surjective. Choosing an \'etale scheme cover of the algebraic space $U$ gives a smooth scheme atlas, which is the definition of algebraicity. Conversely descent for objects and arrows of the stack identifies an object with an $R$-torsor and its equivariant map to $U$; this is precisely the quotient-stack construction. Neither claim assumes that every topos has points or that an arbitrary stack is algebraic.
\end{proof}

\begin{proposition}[Refuted: arbitrary nondegenerate topoi need not have points]\label{B:pointless}
Let $B$ be the complete Boolean algebra of regular open subsets of $\mathbb R$, where $U$ is regular open if $U=\operatorname{int}\overline U$. Regard $B$ as a site whose covers of $U$ are families with join $U$. The Grothendieck topos $\operatorname{Sh}(B)$ is nondegenerate and has no geometric point $\operatorname{Set}\to\operatorname{Sh}(B)$. In particular it has no set-valued fibre functor that preserves finite limits and all colimits.
\end{proposition}
\begin{proof}
Barr \cite[Introduction, p.~265]{BBarr1974} gives the construction for every complete atomless Boolean algebra. We verify all features for this explicit $B$. Meets and joins are
\[
 \bigwedge_\alpha U_\alpha=\operatorname{int}\!\bigcap_\alpha U_\alpha,
 \qquad
 \bigvee_\alpha U_\alpha=\operatorname{int}\overline{\bigcup_\alpha U_\alpha},
 \qquad \neg U=\operatorname{int}(\mathbb R\setminus U).
\]
These are regular open, and give the complete Boolean operations; finite meets distribute over arbitrary joins. Every nonempty regular open contains a nonempty bounded interval properly, so $B$ has no atoms. It is a set and $\varnothing\ne\mathbb R$, so its sheaf topos exists and is nondegenerate: its subterminal-object frame is $B$.

A geometric point's inverse-image functor would induce a frame homomorphism $h:B\to\{0,1\}$, preserving arbitrary joins and finite meets. Boolean complements are consequently preserved. Put $b=\bigvee\{U:h(U)=0\}$. Join preservation gives $h(b)=0$, so $a=\neg b$ is nonzero. If $0<c<a$, either $h(c)=0$, forcing $c\le b$ and hence $c=0$, or $h(c)=1$. In the latter case $h(\neg c)=0$, hence $\neg c\le b$ and $a\le c$, again a contradiction. Thus $a$ would be an atom, which does not exist. This excludes a point. A finite-limit- and colimit-preserving functor from a Grothendieck topos to sets has the right adjoint required for a geometric point, and the same subterminal argument directly excludes it as well. The obstruction is infinite: no finite atomless Boolean algebra exists, so no finite enumeration is presented as a certificate for this example.
\end{proof}

\begin{theorem}[Known: Boolean localization for every Grothendieck topos]\label{B:Boolean-localization}
For every Grothendieck topos $\mathcal E$, there are a complete Boolean algebra $B$ and a surjective geometric morphism
\[
 p:\operatorname{Sh}(B)\longrightarrow\mathcal E.
\]
Equivalently, $p^*:\mathcal E\to\operatorname{Sh}(B)$ preserves finite limits and all colimits and is conservative. It is comonadic, so $\mathcal E$ is recovered as the category of coalgebras for its induced comonad. This gives a universal Boolean-valued replacement for ordinary fibre functors even when $\mathcal E$ has no points.
\end{theorem}
\begin{proof}[Derivation]
Barr \cite[Theorem 4.1]{BBarr1974} states this as existence of an exact cotripleable functor $\mathcal E\to\operatorname{Sh}(B)$. His Corollaries 4.5 and 4.7 factor its construction through a complete Heyting algebra and its Boolean enlargement. The adjunction and conservativity are established in Theorem 3.1 and Corollary 3.4. Thus the functor is the inverse image of a geometric morphism, and its conservativity is exactly surjectivity of that morphism. ``Cotripleable'' is the comonadic assertion, which supplies the claimed reconstruction. Unlike the sufficient coherent-topos hypothesis in Proposition~\ref{B:subject14}, this theorem applies to every Grothendieck topos.
\end{proof}

\subsubsection{15. Homotopy types and infinitesimal homotopy groups}
\begin{theorem}[Known: schematization and a nonabelian comparison]\label{B:subject15}
For every field $k$ and pointed connected simplicial set $T$, there is a schematization $(T\otimes k)^{\mathrm{sch}}$, characterized by the universal map from $T$ to a pointed schematic homotopy type. Its fundamental group is the affine proalgebraic completion of $\pi_1(T)$; for simply connected finite-type $T$ in characteristic zero its higher homotopy groups are the additive group schemes of $\pi_i(T)\otimes_{\mathbb Z}k$. Separately, for a henselian pair $(A,I)$ and a proper $X/A$, the closed fibre $Z=X_{A/I}$ has the same profinite \'etale homotopy type as $X$.
\end{theorem}
\begin{proof}[Derivation]
To\"en \cite[Theorem 3.3.4 and Corollaries 3.3.7, 3.3.9]{BToen2006} constructs the schematization and gives these identifications. Its adjunction, applied to classifying stacks of affine groups, identifies the linear representations of its $\pi_1$ with those of $\pi_1(T)$, which is precisely proalgebraic completion. In the simply connected characteristic-zero finite-type case, the comparison with rational homotopy identifies each higher group with the indicated vector group. The finite-type hypothesis is retained.

The second assertion is Haine--Holzschuh--Wolf \cite[Corollary 2.34]{BHHW2024}. Their theorem says that restriction preserves sections of every torsion sheaf of spaces. Apply it to constant sheaves with any finite homotopy type as value. These section spaces are the mapping spaces out of the profinite \'etale homotopy type; equality for all finite targets identifies the representing pro-objects. Thus an actual homotopy comparison, not only equality of cohomology dimensions, follows. These two constructions make precise two different parts of the handwritten programme. They do not identify an unspecified ``infinitesimal'' homotopy theory of every topos with either construction.
\end{proof}

% Insertable scoped result for Part B, subject 15.  This does not identify the
% unspecified 1964 construction with the later Pursuing Stacks candidate.
\paragraph{A precise later candidate and its limitation.}
In \emph{Pursuing Stacks}, \S111, original manuscript
pp.~447--449, the first proposed category of unipotent bundles is a full
subcategory of fpqc sheaves; maps between vector bundles may therefore be
nonlinear polynomial maps.  In \S118, equations (9)--(10), original
pp.~496--497, the proposed comparison is
\(\pi_i(X_\bullet(k))\simeq H_i(\operatorname{Lie}X_\bullet)\).
The conditions explicitly listed there are \(X_0=X_1=\operatorname{Spec}k\),
smooth terms, and possibly a Kan condition.
The following example tests these stated conditions, without assigning a
meaning to further unspecified restrictions on the proposed model category
\cite{A567_GrothendieckPursuing1983}.

\begin{proposition}[Refuted: smooth Kan conditions do not imply the Lie comparison]\label{B:Lie-comparison-counterexample}
There is a pointed simplicial affine scheme \(X_\bullet\) over
\(k=\mathbf F_2\) such that
\[
 X_0=X_1=\operatorname{Spec}k,
 \qquad X_n\simeq\mathbf A_k^{\,2\binom n2},
\]
all simplicial structure maps are pointed polynomial maps, and
\(X_\bullet(R)\) is Kan for every commutative \(k\)-algebra \(R\), but
\[
 \pi_2\bigl(X_\bullet(k)\bigr)\simeq\mathbf Z/4\mathbf Z,
 \qquad
 H_2\bigl(N\operatorname{Lie}X_\bullet\bigr)
       \simeq(\mathbf Z/2\mathbf Z)^2.
\]
Here \(N\) denotes the normalized chain complex.  In particular, the left
group admits no \(k\)-module structure, and the comparison cannot hold
even as an isomorphism of abelian groups.
\end{proposition}

\begin{proof}
Let \(G=\mathbf A_k^2\) with the commutative group law
\[
 (a,b)+(c,d)=(a+c,\,b+d+ac).
 \tag{*}
\]
For every \(k\)-algebra \(R\), identify \((a,b)\) with
\(1+at+bt^2\) in the multiplicative group of units of \(R[t]/(t^3)\).
Multiplication gives exactly (*), so associativity and commutativity hold
over every such \(R\).  The identity is \((0,0)\), and the inverse is
\((a,b+a^2)\).  Thus \(G\) is a smooth commutative affine group scheme.
On \(k\)-points, \((1,0)\) has successive multiples
\[
 (1,0),\quad(0,1),\quad(1,1),\quad(0,0).
\]
Consequently \(G(k)\) is cyclic of order four.

We give an explicit simplicial construction, so no representability or
infinite-degree assertion is inferred from a finite computation.  For any
abelian group \(A\), let \(K(A,2)_n\) be the group of normalized
\(2\)-cocycles on the ordered simplex with vertices \(0,\ldots,n\).
Its elements are arrays \((u_{ijk})_{i<j<k}\) satisfying
\[
 u_{jkl}-u_{ikl}+u_{ijl}-u_{ijk}=0
 \quad(i<j<k<l).
 \tag{**}
\]
An array is uniquely determined by the coordinates
\(u_{0ij}\), \(1\leq i<j\leq n\): all remaining coordinates are
\[
 u_{ijk}=u_{0jk}-u_{0ik}+u_{0ij}\quad(0<i<j<k).
 \tag{***}
\]
Conversely (***) satisfies (**) by cancellation.  Hence
\(K(A,2)_n\simeq A^{\binom n2}\).
For an order-preserving map between finite ordered sets, pull back a
cochain along the induced vertex map and set a value with repeated
vertices equal to zero.  Pullback preserves the cocycle equation and
composition.  It therefore defines all faces and degeneracies and proves
the simplicial identities in every degree.

Apply this construction functorially to \(A=G(R)\).  Formulas (**)--(***)
use only the group law, inverse, and zero, which are polynomial maps on
\(G\).  Thus the result is represented by a simplicial affine group scheme
\(X_\bullet=K(G,2)\), with
\(X_n=G^{\binom n2}\simeq\mathbf A_k^{2\binom n2}\) and the asserted
first two terms.  Every \(X_\bullet(R)\) is a simplicial group and hence a
Kan complex.  The usual proof fills horns by finite successive
applications of group operations and degeneracies; it is natural in
\(R\), so here it even supplies polynomial sections of the horn maps
\cite[Tag~08NZ]{A567_StacksSimplicial}.

For completeness, the normalization can be computed in all degrees.
Use the convention
\[
 N_nK(A,2)=\bigcap_{i=1}^n\ker(d_i),\qquad\partial=d_0.
\]
Its groups in degrees zero and one vanish, and its degree-two group is
\(A\).  In degree three, the conditions \(d_1=d_2=d_3=0\) successively
kill the free coordinates \(u_{023},u_{013},u_{012}\), so the entire
cocycle vanishes.  If \(n\geq4\), each triple of vertices omits at least
one positive vertex; the corresponding vanishing face kills its
coordinate.  Thus the normalized complex has only \(A\) in degree two.
The homotopy groups of a simplicial abelian group are the homology of its
normalized complex; the group operation agrees with the pointwise
abelian operation by the interchange law
\cite[Chapter~III, Corollary~2.7, pp.~152--153]{A567_GoerssJardine1999}.
Equivalently in degree two, homotopies of cycles change a cycle by the
boundary of a normalized three-simplex, and that normalized group is zero
here.  We obtain
\(\pi_2(X_\bullet(k))=G(k)\simeq\mathbf Z/4\mathbf Z\).

Finally, insert \((a,b)=(\epsilon u,\epsilon v)\) and
\((c,d)=(\epsilon u',\epsilon v')\) into (*), with
\(\epsilon^2=0\).  The cross term vanishes, and the resulting tangent law
is \((u,v)+(u',v')=(u+u',v+v')\).  Therefore
\(\operatorname{Lie}G=k^2\) with its ordinary additive law.
Differentiation preserves finite products and sends every sum and inverse
in the construction above to the corresponding linear operation.
Consequently \(\operatorname{Lie}X_\bullet=K(k^2,2)\).
The same all-degree normalization calculation gives
\(H_2(N\operatorname{Lie}X_\bullet)=k^2\).
Its additive group has exponent two, whereas \(G(k)\) has an element of
order four.  This proves every assertion.
\end{proof}

\paragraph{Exact valid replacement and historical scope.}
If the terms are vector group schemes and all faces and degeneracies are
linear maps, the comparison does hold: rational points and tangent spaces
are then the same simplicial \(k\)-module, and normalized homology computes
its homotopy groups.  The example shows why allowing arbitrary polynomial
structure maps changes that conclusion.  Grothendieck also considers
defining weak equivalences directly by Lie homology in \S118; that is a
different choice of invariant, and the proposition does not deny its
existence.  No identification of either later choice with the undefined
1964 higher infinitesimal construction has been proved here.

\paragraph{Finite certificate and its limits.}
The program \texttt{code/A567\_Lie\_comparison.py} checks all four rational
points, all 64 associativity triples, and all 16 dual-number tangent-law
pairs, with 117 assertions in total.  Its recorded output is
\texttt{logs/A567\_Lie\_comparison.json}; the elapsed time is measured in
that run.  It certifies the finite arithmetic distinction
\(C_4\not\simeq C_2^2\).  The simplicial identities, smoothness, Kan
condition and all-degree homotopy calculation are supplied by the proof,
not extrapolated from this finite check.

\begin{proposition}[Refuted: the same comparison over the integers]\label{B:Lie-comparison-integral-counterexample}
There is a pointed simplicial affine scheme \(Y_\bullet\) over
\(\mathbf Z\) with
\[
 Y_0=Y_1=\operatorname{Spec}\mathbf Z,
 \qquad
 Y_n\simeq\mathbf A_{\mathbf Z}^{\,2\binom n2+\binom n3},
\]
polynomial structure maps, and a Kan simplicial set \(Y_\bullet(R)\) for
every commutative ring \(R\), such that
\[
 \pi_2\bigl(Y_\bullet(\mathbf Z)\bigr)\simeq\mathbf Z,
 \qquad
 H_2\bigl(N\operatorname{Lie}Y_\bullet\bigr)
       \simeq\mathbf Z\oplus\mathbf Z/2\mathbf Z.
\]
Thus restricting the displayed comparison of \S118 to \(k=\mathbf Z\)
does not repair its stated smooth Kan sufficient conditions.
\end{proposition}

\begin{proof}
Define the smooth commutative group scheme \(G=\mathbf A_{\mathbf Z}^2\)
by the same truncated-unit multiplication law
\[
 (a,b)+(c,d)=(a+c,\,b+d+ac).
\]
The identification with \(1+at+bt^2\) modulo \(t^3\) works over every
commutative ring; the inverse is now \((-a,a^2-b)\).
The polynomial map
\[
 \delta:\mathbf G_a\longrightarrow G,
 \qquad t\longmapsto(2t,2t^2)
\]
is a homomorphism: its second coordinate at \(s+t\) is
\(2s^2+2t^2+4st\), which is precisely the second coordinate of
\(\delta(s)+\delta(t)\).

In the abelian category of fpqc sheaves of abelian groups on affine
\(\mathbf Z\)-schemes, form the nonnegative chain complex \(C_\bullet\)
with \(C_3=\mathbf G_a\), \(C_2=G\), differential \(\delta\), and all
other terms zero.  Apply the Dold--Kan construction
\cite[Theorem~14.24.3, Tag~019G]{A567_StacksSimplicial}.
Its degree \(n\) is the direct sum of \(C_r\) over order-preserving
surjections \([n]\twoheadrightarrow[r]\).  There are \(\binom nr\)
such surjections.  Here all direct sums in a fixed degree are finite,
and hence are products, so the resulting simplicial sheaf is represented
by
\[
 Y_n=G^{\binom n2}\times\mathbf G_a^{\binom n3}.
\]
The Dold--Kan structure maps use only group sums, inverses, identities,
zero maps, and \(\delta\); each is therefore a polynomial morphism of
these affine schemes.  The construction proves all simplicial identities
and all-degree normalization \(NC_\bullet^{\mathrm{DK}}\simeq C_\bullet\).
Every term is smooth over \(\mathbf Z\), the first two are points, and
every ring-valued simplicial group is Kan.

Taking \(\mathbf Z\)-points commutes with this construction: it
preserves finite products and every group operation and differential
entering the displayed formulas.  Thus the normalized complex of
\(Y_\bullet(\mathbf Z)\) is
\[
 \mathbf Z\xrightarrow{\delta}G(\mathbf Z)
\]
in degrees three and two.  The integer-valued map
\[
 \theta:G(\mathbf Z)\longrightarrow\mathbf Z^2,
 \qquad (a,b)\longmapsto\left(a,b-\frac{a(a-1)}2\right)
\]
is an isomorphism of abelian groups.  Indeed \(a(a-1)\) is even for
every integer \(a\); the inverse sends
\((a,c)\) to \((a,c+a(a-1)/2)\); and
\[
 \binom{a+c}{2}=\binom a2+\binom c2+ac
\]
checks additivity.  This is an isomorphism on integer points; no
polynomial morphism over \(\mathbf Z\) is asserted.
Under \(\theta\), the differential becomes
\[
 t\longmapsto(2t,t).
\]
The quotient of \(\mathbf Z^2\) by this primitive vector is
\(\mathbf Z\): the map \((a,c)\mapsto a-2c\) is surjective and its
kernel is exactly \(\mathbf Z(2,1)\).  Normalized homology therefore
gives the first asserted homotopy group.

On tangent spaces, \(\operatorname{Lie}G=\mathbf Z^2\) with the
ordinary additive law, while
\[
 \operatorname{Lie}\delta:\mathbf Z\longrightarrow\mathbf Z^2,
 \qquad t\longmapsto(2t,0).
\]
The Lie functor preserves finite products, group sums, and the
compositions in the Dold--Kan formulas.  Hence
\(\operatorname{Lie}Y_\bullet\) is the Dold--Kan simplicial module of
the two-term Lie complex just displayed.  No exactness of the Lie
functor is needed or claimed.  Its degree-two normalized homology is
\(\mathbf Z^2/\mathbf Z(2,0)\simeq\mathbf Z\oplus\mathbf Z/2\mathbf Z\).
The torsion subgroup distinguishes it from \(\mathbf Z\), proving the
claim.
\end{proof}

The exact certificate \path{code/A567_Lie_comparison_integral.py}
verifies the polynomial homomorphism and additivity identities, the
derivative, and the two Smith normal forms.  Its seven assertions and
measured runtime are recorded in
\path{logs/A567_Lie_comparison_integral.json}.  These are symbolic
identities and exact integer matrix operations, not a finite-range
experiment.  The all-degree simplicial construction and the distinction
between integer-valued and integral polynomial maps are proved above.


\begin{theorem}[Known, with derivation: integral affinization and its exact comparison]\label{B:integral-affinization}
Let $T$ be a pointed simply connected space with every $\pi_i(T)$,
$i\ge2$, finitely generated. Work in hypercomplete fpqc stacks on
affine schemes over $\mathbf Z$. Let affine stacks mean the full
subcategory generated by the $K(\mathbf G_a,n)$ under small homotopy
limits, and put
\[
 F_T=(T\otimes\mathbf Z)^{\mathrm{uni}}
     =\mathbf R\operatorname{Spec}C^*(T,\mathbf Z),\qquad
 \mathbb H=\operatorname{Spec}\operatorname{Int}(\mathbf Z).
\]
Here $C^*$ is the cosimplicial commutative cochain ring, and
$\operatorname{Int}(\mathbf Z)=\{P\in\mathbf Q[x]:P(\mathbf Z)\subset\mathbf Z\}$,
with the additive Hopf structure. Then $F_T$ is universal for maps from
$T$ to integer-valued points of affine stacks, and
\[
 \pi_i(F_T)\simeq\pi_i(T)\otimes_{\mathbf Z}\mathbb H\quad(i\ge2)
\]
as fpqc abelian sheaves; these sheaves have natural $\mathbb H$-module
structures. The unit $u_T:T\to F_T(\mathbf Z)$ has a natural retraction.
The induced map of mapping spaces
$\operatorname{Map}(T,S)\to\operatorname{Map}(F_T,F_S)$ also has a
retraction for every $S$ in this class, and $F_T\simeq F_S$ if and only
if $T\simeq S$. At the induced basepoint,
\[
 \pi_i(F_T(\mathbf Z))\simeq
 \pi_i(T)\oplus\prod_{p\text{ prime}}\widehat{\pi_{i+1}(T)}_p
 \qquad(i\ge1).
\]
The homotopy sheaves need not be $\mathcal O$-modules, and the unit need
not be a weak equivalence. These assertions concern this specified
integral affine-stack construction, without identifying it with the
unspecified infinitesimal enrichment in the 1964 sheet.
\end{theorem}
\begin{proof}
To\"en \cite[\S1, Definition~2.1, Theorem~4.1 and Corollary~4.4]{BToen2020}
defines this adjunction, the Hilbert group $\mathbb H$, and proves the
homotopy-sheaf and module assertions with exactly the stated hypotheses.
Theorem~4.1 follows from the Eilenberg--Mac Lane calculation in Lemma~4.2,
finite Postnikov induction using Eilenberg--Moore, and the countable fpqc
limit argument. We use this proved theorem, not an ordinary spectrum of
the graded cohomology ring.

For the more explicit point comparison, set
$A_f=(\prod_p\mathbf Z_p)\otimes\mathbf Q$.
The arithmetic square and the comparisons in \S5 give
\[
 F_T(\mathbf Z)\simeq
 F_T(\mathbf Q)\mathbin{\times^h}_{F_T(A_f)}\prod_pF_T(\mathbf Z_p).
\]
The complete-local terms are the free-loop spaces $L(T_p)$ of the
$p$-completions. Lemma~5.3 factors their maps through loop evaluation;
Remark~5.6, together with the arithmetic fracture square, identifies
this functorially as
\[
 F_T(\mathbf Z)\simeq
 T\mathbin{\times^h}_{\prod_pT_p}\prod_pL(T_p).
\]
Projection and constant loops give $r_Tu_T=\mathrm{id}_T$, as in
Corollary~5.4. Evaluation on each free-loop space has fibre $\Omega T_p$
and a constant-loop section. The split homotopy sequence and finite
generation give exactly the displayed formula of Corollary~5.1.
For $i=1$, simple connectivity removes a possible nonabelian action:
the remaining group is the product of the completed second homotopy
groups.

A map $g:F_T\to F_S$ is sent to $r_S\circ g(\mathbf Z)\circ u_T$.
Naturality makes this a retraction of the map of mapping spaces.
Corollary~5.5 proves compatibility with homotopy-sheaf maps; since
Corollary~5.2 gives $\Gamma(\mathbf Z,M\otimes\mathbb H)=M$ for finitely
generated abelian $M$, these sheaves detect their underlying groups.
An equivalence $F_T\simeq F_S$ therefore returns a weak equivalence
$T\simeq S$. This argument does not assume that the retraction preserves
composition of arbitrary affine-stack maps.

Finally take $T=K(\mathbf Z,2)$. Corollaries~2.3 and~3.4 give
$\pi_2(F_T)=\mathbb H$ and $\mathbb H(\mathbf F_2)=\mathbf Z_2$.
An $\mathcal O$-module evaluated at $\mathbf F_2$ would be annihilated
by two, whereas $\mathbf Z_2$ is not. The point formula also gives
$\pi_1(F_T(\mathbf Z))=\prod_p\mathbf Z_p\ne0$, while $\pi_1(T)=0$.
Thus the two stated limitations are explicit.
\end{proof}

\begin{theorem}[Known, with derivation: unipotent homotopy of schemes and its comparisons]\label{B:unipotent-homotopy}
Let $k$ be a field. A higher stack here is a sheaf of spaces for the
fpqc topology on affine $k$-schemes. An affine stack is
$\operatorname{Spec}B$, where $B$ belongs to the homotopy theory of
cosimplicial commutative $k$-algebras. For a pointed higher stack $(X,x)$
such that $k\xrightarrow{\sim}H^0(X,\mathcal O_X)$, put
\[
 \mathbf U(X)=\operatorname{Spec}R\Gamma(X,\mathcal O_X),\qquad
 \pi_i^{\mathrm U}(X,x)=\pi_i(\mathbf U(X),x)\quad(i\ge1).
\]
The derived algebra structure, including its cosimplicial commutative
operations, is part of this definition. Then the following hold.
\begin{enumerate}
\item The map $X\to\mathbf U(X)$ is universal for maps to affine stacks.
The stack $\mathbf U(X)$ is connected and hypercomplete, and its
$\pi_i^{\mathrm U}$ are represented by unipotent affine group schemes,
commutative for $i\ge2$. Here unipotent means that every nonzero
finite-dimensional representation has a nonzero invariant vector.
Finite type of these group schemes is not required.
\item If $X$ is a scheme of finite type over $k$ and $x\in X(k)$, then
$\pi_1^{\mathrm U}(X,x)$ is naturally the Tannaka group of the category
of unipotent vector bundles on $X$, with fibre functor $x^*$.
Here a unipotent vector bundle has a finite filtration by subbundles
whose successive quotients are trivial line bundles. This is the
unipotent Nori fundamental group in this specified sense.
\item For every integer $r\ge1$,
\[
 H^j(X,\mathcal O_X)=0\ (1\le j\le r)
 \quad\Longleftrightarrow\quad
 \pi_j^{\mathrm U}(X,x)=1\ (1\le j\le r).
\]
Under these equivalent conditions there are natural maps
\[
 H^{r+1}(X,\mathcal O_X)
   \simeq\operatorname{Hom}(\pi_{r+1}^{\mathrm U}(X,x),\mathbf G_a),
 \qquad
 \operatorname{Ext}^1(\pi_{r+1}^{\mathrm U}(X,x),\mathbf G_a)
   \hookrightarrow H^{r+2}(X,\mathcal O_X).
\]
The extension group is for commutative fpqc group sheaves; the
restriction $r\ge1$ ensures that its first argument is commutative.
\item If $X$ is a proper $k$-scheme and $k$ has characteristic $p>0$, every
$\pi_i^{\mathrm U}(X,x)$ is profinite. If moreover $k$ is algebraically
closed, the comparison induced by $\mathbf F_p\to\mathcal O_X$ identifies
its maximal pro-\'etale quotient with the affine group scheme associated
to the pro-$p$-finite group
\[
 \pi_i\bigl(\operatorname{Et}(X)^{\wedge}_p,x\bigr).
\]
Here $\operatorname{Et}(X)^{\wedge}_p$ is the pro-$p$-finite completion
of the entire \'etale homotopy type. The displayed group retains its
profinite topology. Still assuming $k$ algebraically closed, if
Frobenius acts bijectively on every $H^j(X,\mathcal O_X)$, these
comparison maps are themselves isomorphisms.
\end{enumerate}
These conclusions concern a specified theory of schemes and higher stacks.
They make no identification with the unspecified infinitesimal enrichment
of the 1964 thesis sheet.
\end{theorem}
\begin{proof}[Derivation from the opened primary theorems]
All Mondal--Reinecke locators below refer to their opened 95-page final
author version of 18 August 2025; the publication is in the January 2026
volume of the journal \cite{BMondalReinecke2026}.

\emph{Affinization and degree one.}
Definitions~2.1.8--2.1.12 and Remark~2.1.16 give the adjunction
between derived global sections and $\operatorname{Spec}$, with
$R\Gamma(\operatorname{Spec}B,\mathcal O)\simeq B$.
Thus Definition~3.0.1 and Remark~3.0.2 give exactly the stated
universal map. Theorem~2.1.17 and Remark~2.1.18, which invoke
To\"en's original Theorems~2.4.1 and~2.4.5
\cite[pp.~51--61]{BToen2006}, give connectedness, hypercompleteness
and unipotent representability from the augmentation and $H^0=k$.
They do not impose finite type on the representing groups.

For the degree-one identification, a pointed $G$-torsor is equivalently
a pointed map $X\to BG$. If $G$ is a unipotent affine group scheme,
$BG$ is an affine stack and is $1$-truncated. Hence each such map
factors uniquely through
\[
 X\longrightarrow\mathbf U(X)\longrightarrow
 \tau_{\le1}\mathbf U(X)=B\pi_1^{\mathrm U}(X,x).
\]
The unipotent-bundle Tannaka group has precisely this universal pointed
torsor property, as stated and derived in
\cite[Definition~3.1.1, Lemmas~3.1.2--3.1.6]{BMondalReinecke2026}.
Applying the two universal properties to one another produces inverse
homomorphisms, proving Proposition~3.1.7 and assertion~(2).
This description fixes which Nori group is meant; the finite-flat
torsor group in Theorem~\ref{B:nori-finite-flat} has a separately
specified category of representations.

\emph{Hurewicz.}
Put $F=\mathbf U(X)$. Proposition~2.1.19 says that its Postnikov
truncations are affine, that
\[
 H^j(\tau_{\le m}F,\mathcal O)\longrightarrow H^j(F,\mathcal O)
\]
is an isomorphism for $j\le m$ and injective for $j=m+1$, and that
$F\simeq\varprojlim_m\tau_{\le m}F$.
The adjunction identifies $R\Gamma(F,\mathcal O)$ with
$R\Gamma(X,\mathcal O_X)$. If the first $r$ homotopy groups vanish,
$\tau_{\le r}F$ is the point, so the asserted cohomology groups vanish.
Conversely, proceed by induction. In degree one,
\[
 H^1(F,\mathcal O)
  =H^1(B\pi_1(F),\mathcal O)
  =\operatorname{Hom}(\pi_1(F),\mathbf G_a).
\]
A nontrivial unipotent affine group admits a nonzero morphism to
$\mathbf G_a$, so degree-one vanishing kills $\pi_1(F)$.
For $2\le j\le r$, the preceding induction steps identify
$\tau_{\le j}F$ with $K(\pi_j(F),j)$. Lemma~3.2.10 then gives
\[
 H^j(F,\mathcal O)=\operatorname{Hom}(\pi_j(F),\mathbf G_a),
\]
so the same detection argument kills $\pi_j(F)$.
Finally $\tau_{\le r+1}F=K(\pi_{r+1}(F),r+1)$.
Lemma~3.2.10 identifies its degree $r+1$ cohomology with the stated
Hom group and its degree $r+2$ cohomology with the stated Ext group,
since $r+1\ge2$. The Postnikov isomorphism and injection supply
assertion~(3). This is the positive-connectivity scope of
\cite[Proposition~3.2.11 and Corollary~3.2.12]{BMondalReinecke2026}.

\emph{Profiniteness.}
In characteristic $p$, write $k_\sigma[F]$ for the Frobenius skew
polynomial ring, with $Fa=a^pF$. Call a module torsion if every vector
lies in a finite-dimensional Frobenius-stable subspace.
Proposition~4.2.31 detects profiniteness of a unipotent affine group
$G$ by torsion of $\operatorname{Hom}(G,\mathbf G_a)$.
The torsion modules form a Serre subcategory after extending to an
algebraic closure, and Corollary~4.2.20 descends profiniteness.
Proposition~4.2.44 and Lemma~4.3.6 show that a truncated affine stack
whose homotopy groups are profinite has torsion cohomology with any
torsion Frobenius representation. Proposition~4.2.45 supplies the
converse detection needed in the induction: for a profinite unipotent
group, if the invariant vectors of an F-representation are torsion,
then the whole representation is torsion.

Here is the precise induction used by Proposition~4.3.7. Degree one
follows from $H^1(F,\mathcal O)=\operatorname{Hom}(\pi_1(F),\mathbf G_a)$.
Assume the groups through degree $m$ are profinite. The Leray sequence
for the next Postnikov fibration has
\[
 E_2^{a,b}=
 H^a\bigl(\tau_{\le m}F,
 H^b(K(\pi_{m+1}(F),m+1),\mathcal O)\bigr).
\]
If $V=\operatorname{Hom}(\pi_{m+1}(F),\mathbf G_a)$ were not torsion,
Proposition~4.2.45 would make $E_2^{0,m+1}$ non-torsion.
The intervening rows $0<b\le m$ vanish. Its only possible outgoing
differential is therefore on page $m+2$, with target a subquotient
of the torsion module $H^{m+2}(\tau_{\le m}F,\mathcal O)$.
Since torsion modules are closed under extensions, the kernel remains
non-torsion; no differential enters from a negative column.
It survives to $E_\infty$, contradicting torsion of
$H^{m+1}(F,\mathcal O)$. This proves the induction whenever all
$H^j(X,\mathcal O_X)$ are torsion Frobenius modules.
For proper $X/k$ they are finite dimensional, so the hypothesis holds.
This proves Proposition~4.3.1 with exactly its stated generality.

\emph{The \'etale quotient.}
Now let $k$ be algebraically closed and put
$A=R\Gamma(X,\mathcal O_X)$. The Artin--Schreier sequence identifies
$R\Gamma_{\acute et}(X,\mathbf F_p)$ with $A^{F=1}$.
Proposition~4.1.9, applied to the finite-dimensional Frobenius
cohomology modules, yields
\[
 R\Gamma_{\acute et}(X,k)\simeq A^{\mathrm{perf}}.
\]
Frobenius twists are taken over the perfect field $k$ as in that
proposition. Since $\operatorname{Spec}$ converts colimits of these
algebras to limits of affine stacks,
\[
 \operatorname{Spec}R\Gamma_{\acute et}(X,k)
  \simeq\mathbf U(X)^{\mathrm{perf}}
  =\varprojlim_F\mathbf U(X).
\]
The natural map
\[
 \phi_i:\pi_i(\mathbf U(X)^{\mathrm{perf}})
       \longrightarrow(\pi_i^{\mathrm U}(X))^{\mathrm{perf}}
\]
is injective: its composite with the perfection map and the
\'etale comparison is an isomorphism by the preceding equivalence.
It is also surjective by the Milnor sequence for pointed objects
in the replete fpqc topos
\cite[Proposition~3.22, author pp.~10--11]{BMondalReineckePostnikov2025}:
\[
 1\longrightarrow\varprojlim{}^1\pi_{i+1}(F_n)
 \longrightarrow\pi_i(\varprojlim F_n)
 \longrightarrow\varprojlim\pi_i(F_n)\longrightarrow1.
\]
Applying this to the Frobenius tower proves that $\phi_i$ is an
isomorphism, without assuming that homotopy groups commute with
arbitrary limits. By Proposition~4.2.17, perfection of a profinite
group scheme over a perfect field is its maximal pro-\'etale quotient.
This proves Proposition~4.4.7.

Definition~4.4.2 names the target
$\pi_i(\operatorname{Spec}R\Gamma_{\acute et}(X,k))$.
Proposition~4.4.3 identifies it with the affine group scheme of
$\pi_i(\operatorname{Et}(X)^\wedge_p)$: cochains with
$\mathbf F_p$ coefficients are unchanged by pro-$p$-finite completion,
because their Eilenberg--Mac Lane targets are $p$-finite, and the
$p$-adic affine-stack comparison identifies the resulting homotopy
groups. The source uses To\"en's Corollary~2.5.3
\cite[p.~64]{BToen2006}, with its nilpotent finite-type hypotheses,
and the pro-object argument of its own Remark~4.4.4.
This conclusion concerns completion of the shape; no interchange
with completion of individual uncompleted homotopy groups is used.
Finally, if Frobenius is bijective on all coherent cohomology, then
$A\to A^{\mathrm{perf}}$ is already an equivalence. The same cochain
comparison proves the last assertion, as in Remark~4.4.8.
\end{proof}


\begin{theorem}[Known: formal duals and the Artin--Mazur comparison]\label{B:formal-homotopy}
Let $k$ be a field, let $n\geq2$, and let $X$ be a pointed higher stack
for the fpqc topology on affine $k$-schemes.  Use the affine-stack
construction
\[
 \mathbf U(X)=\operatorname{Spec}R\Gamma(X,\mathcal O),
 \qquad G=\pi_n^{\mathrm U}(X)=\pi_n(\mathbf U(X)),
\]
where $R\Gamma(X,\mathcal O)$ retains its cosimplicial commutative
algebra structure, in the category used in
Theorem~\ref{B:unipotent-homotopy}.
Suppose that
\[
 \begin{gathered}
 H^0(X,\mathcal O)=k,\qquad
 H^i(X,\mathcal O)=0\quad(0<i<n),\\
 g=\dim_k H^n(X,\mathcal O)<\infty,
 \qquad H^{n+1}(X,\mathcal O)=0.
 \end{gathered}
\]
Then $G$ is a commutative affine unipotent group scheme and its dual
$G^\vee$ is a commutative formal Lie group of dimension $g$.
Here the dual is defined by the linearly compact topological Hopf algebra
$\mathcal O(G)^*$, with the topology dual to the finite-dimensional
subcoalgebras of $\mathcal O(G)$; the conclusion includes an isomorphism
of its underlying topological algebra with $k[[t_1,\ldots,t_g]]$.

If, in addition, $k$ is algebraically closed of characteristic $p>0$
and $X$ is a proper $k$-scheme with a chosen $k$-point, let
$\operatorname{Art}_k$ be the category of local Artinian $k$-algebras
with specified residue identification $A/\mathfrak m_A=k$.  The covariant functor $\Phi_X^n:\operatorname{Art}_k\to\mathrm{Ab}$
is defined by
\[
 \Phi_X^n(A)=\ker\!\left(
 H^n_{\mathrm{\acute et}}(X\times_k\operatorname{Spec}A,\mathbf G_m)
 \longrightarrow H^n_{\mathrm{\acute et}}(X,\mathbf G_m)\right)
 \qquad(A\in\operatorname{Art}_k)
\]
and is represented by the formal Lie group $G^\vee$.  In particular,
\[
 \Phi_X^n\simeq (\pi_n^{\mathrm U}(X))^\vee,
 \qquad
 \mathbf M(\Phi_X^n)
 \simeq \varprojlim_r H^n_{\mathrm{Zar}}(X,W_r),
\]
where $W_r$ denotes the length-$r$ $p$-typical Witt-vector group and
$\mathbf M$ is the formal-group Dieudonn\'e functor in the convention
specified below.  Smoothness, projectivity, ordinarity and finite
formal-group height are not needed for this comparison.
\end{theorem}

\begin{proof}[Derivation from the opened sources]
The first assertion is Mondal--Reinecke
\cite[Proposition~6.2.5 and Construction~6.2.6, p.~76]{BMondalReinecke2026}.
We give the mechanism, including the criterion for the formal dual.
The augmentation and $H^0=k$ make $\mathbf U(X)$ a pointed connected
affine stack.  Their Hurewicz theorem
\cite[Proposition~3.2.11 and Corollary~3.2.12]{BMondalReinecke2026}
and the lower cohomology vanishings give
\[
 \pi_i^{\mathrm U}(X)=0\quad(0<i<n),
 \qquad \tau_{\leq n}\mathbf U(X)\simeq K(G,n).
\]
Their affine-stack representability theorem gives $G$ as an affine
unipotent group scheme; it is commutative since $n\geq2$.
The Postnikov cohomology comparison
\cite[Proposition~2.1.19]{BMondalReinecke2026}
identifies degree-$n$ cohomology and injects degree-$(n+1)$ cohomology.
Consequently
\[
 \begin{aligned}
 \operatorname{Hom}(G,\mathbf G_a)
  &\simeq H^n(K(G,n),\mathcal O)
   \simeq H^n(X,\mathcal O),\\
 \operatorname{Ext}^1(G,\mathbf G_a)
  &\simeq H^{n+1}(K(G,n),\mathcal O)=0.
 \end{aligned}
\]
The Hom and Ext identifications are their Lemma~3.2.10; the Ext group
is taken in the abelian category of fpqc sheaves of abelian groups.

The criterion needed here is precisely
\cite[Proposition~6.1.13]{BMondalReinecke2026}:
a commutative affine unipotent group $G$ with finite-dimensional
$\operatorname{Hom}(G,\mathbf G_a)$ has a smooth formal dual if and
only if this Ext group vanishes.  In characteristic zero a commutative
affine unipotent group is an inverse limit of vector groups; finite
$\dim\operatorname{Hom}(G,\mathbf G_a)=g$ makes it $\mathbf G_a^g$,
whose dual is the additive formal group of dimension $g$.
In positive characteristic its proof in \S6.1 makes the
algebraic content explicit.  The topological dual $D=\mathcal O(G)^*$
is a complete noetherian local algebra with residue field $k$ and
cotangent dimension $g$, hence has a minimal presentation
\[
 D=k[[t_1,\ldots,t_g]]/J,\qquad J\subseteq(t_1,\ldots,t_g)^2.
\]
Duality identifies $H^2(BG,\mathcal O)$ with
$\operatorname{Ext}^2_D(k,k)$.  The calculation in their
Lemma~6.1.14, including its characteristic-two argument, gives
\[
 \dim_k H^2(BG,\mathcal O)
 =\binom g2+\dim_k\operatorname{Ext}^1(G,\mathbf G_a).
\]
For the displayed minimal presentation, the degree-two Betti-number
calculation used there gives
\[
 \dim_k\operatorname{Ext}^2_D(k,k)
 =\binom g2+\dim_k J/(t_1,\ldots,t_g)J.
\]
Thus Ext vanishing forces $J=0$ by Nakayama's lemma.  This proves the
formal-dual assertion, with the Hopf structure inherited from $G$.
If $g=0$, the same argument gives $D=k$ and the trivial formal group.

Now impose the additional hypotheses for the comparison.  We first
check the classical deformation functor.  On an Artinian test algebra
$A$, the nilpotent-units sheaf $1+\mathfrak m_A\mathcal O_{X_A}$ has
a finite filtration with quotients of the form
$\mathcal O_X\otimes_k(\mathfrak m_A^j/\mathfrak m_A^{j+1})$.
Since $H^{n-1}(X,\mathcal O)=0$, induction through this filtration
shows that the previous-degree deformation functor $\Phi_X^{n-1}$
is zero, hence formally smooth.  For a small Artinian extension with
square-zero kernel $I$, the obstruction to lifting a degree-$n$
deformation lies in $H^{n+1}(X,\mathcal O)\otimes_k I$, which is zero.
The tangent space is $H^n(X,\mathcal O)$ and is finite-dimensional by
properness.  Artin--Mazur's Schlessinger argument
\cite[II, Corollaries~2.4, 2.5, 2.11--2.12; II, Corollary~4.2,
pp.~106--109]{CArtinMazur1977}
therefore represents $\Phi_X^n$ by a commutative formal Lie group of
dimension $g$.  The residue identification in
$\operatorname{Art}_k$ supplies the restriction map used in its
definition.

We now identify that group with $G^\vee$, following the complete
proof of \cite[Proposition~6.3.6, pp.~77--78]{BMondalReinecke2026}.
The Postnikov map $X\to\mathbf U(X)\to K(G,n)$ and the universal
property of $\mathbf U(X)$ give compatible natural maps
\[
 v_r:\operatorname{Hom}(G,W_r)\longrightarrow H^n(X,W_r).
\]
Indeed, $K(W_r,n)$ is an affine stack since $W_r$ is unipotent, and
it is $n$-truncated, so maps to it factor through
$\tau_{\leq n}\mathbf U(X)$.
The finite filtration of $W_r$ by additive groups identifies its
Zariski and fpqc cohomology here, so either topology computes the
right side.

We prove simultaneously, for every $r\geq1$, that
\[
 H^{n+1}(X,W_r)=0,
 \qquad\operatorname{Ext}^1(G,W_r)=0,
 \qquad v_r\text{ is an isomorphism}.
\]
For $r=1$, these follow from $W_1=\mathbf G_a$, the hypotheses and
the Hom--Ext calculation above.  The short exact sequence
\[
 0\longrightarrow W_r\longrightarrow W_{r+1}
 \longrightarrow\mathbf G_a\longrightarrow0
\]
and its long exact cohomology and Ext sequences propagate the first
two assertions from $r$ to $r+1$.  They also give the compatible
short exact rows
\[
\begin{aligned}
0&\longrightarrow\operatorname{Hom}(G,W_r)
 \longrightarrow\operatorname{Hom}(G,W_{r+1})
 \longrightarrow\operatorname{Hom}(G,\mathbf G_a)
 \longrightarrow0,\\
0&\longrightarrow H^n(X,W_r)
 \longrightarrow H^n(X,W_{r+1})
 \longrightarrow H^n(X,\mathcal O)
 \longrightarrow0.
\end{aligned}
\]
On the lower left use $H^{n-1}(X,\mathcal O)=0$; on the right use
the two vanishings just proved.  The maps between these rows are
$v_r,v_{r+1},v_1$.  The short five lemma proves the induction step.

For a commutative formal Lie group $E$ over the perfect field $k$,
use the convention
\[
 \mathbf M(E)=\operatorname{Hom}(\widehat W,E)
 \simeq\operatorname{Hom}(E^\vee,W),
 \qquad W=\varprojlim_rW_r.
\]
Here $\widehat W$ is the formal completion of the Witt group at zero;
its dual is $W$ by the Artin--Hasse exponential, as proved in
\cite[Lemma~6.3.4]{BMondalReinecke2026}.
Taking inverse limits of the $v_r$ gives
\[
 \mathbf M(G^\vee)
 =\operatorname{Hom}(G,W)
 \simeq\varprojlim_rH^n(X,W_r)
 \simeq\mathbf M(\Phi_X^n).
\]
The last identification is the original
Artin--Mazur \cite[II, Corollary~4.3, p.~110]{CArtinMazur1977}.
These maps are natural in the Witt-vector coefficients, hence
compatible with the Frobenius and Verschiebung operators and the
Dieudonn\'e module structures.  Full faithfulness of the formal-group
Dieudonn\'e functor, applied in
\cite[Proposition~6.3.6]{BMondalReinecke2026}, yields
$G^\vee\simeq\Phi_X^n$ as asserted.
\end{proof}

\paragraph{Scope in degree one.}
The same source identifies $\Phi_X^1$ with
$(\pi_1^{\mathrm U}(X)^{\mathrm{ab}})^\vee$, not with the dual of the
whole generally noncommutative group.  The theorem above and the
corresponding equivalence
$\tau_{\leq n}\mathbf U(X)\simeq K((\Phi_X^n)^\vee,n)$ are deliberately
stated only for $n\geq2$.  Thus the displayed equivalence in
\cite[Construction~6.3.7]{BMondalReinecke2026} is not imported in
degree one without its necessary abelianization qualification.


\paragraph{Earlier higher homotopy and the remaining historical identification.}
An earlier higher construction is explicit in Grothendieck's letter to
Serre of 18 July 1959 \cite[pp.~88--90]{A567_GrothendieckSerre2004}:
apply connected components to the \v Cech nerves of pointed unramified
covering families and form the resulting pro-homotopy invariants. The
letter expressly says that this construction ignores inseparability.
SGA~1, X, \S1.6 \cite[retyped p.~203; original p.~264]{A567_SGA1_2003}
also acknowledges a provisional definition of higher groups. The
remaining reconstruction question concerns the exact infinitesimal
enrichment and comparison intended in 1964, not the existence of any
earlier higher topological construction.


\subsubsection{16. Specialization of the full fundamental group}
\begin{corollary}[Refuted, with the valid specialization theorem]\label{B:subject16}
The blanket assertion that full \'etale fundamental groups are invariant in smooth proper curve families is false, even for the explicit genus-one family of Proposition~\ref{B:specialization}. For such families over a trait of residue characteristic $p$, specialization is surjective and induces an isomorphism on maximal prime-to-$p$ quotients. At the homotopy level, after completion away from the residue characteristics, the geometric fibres of a smooth proper morphism of quasi-compact quasi-separated schemes have equivalent \'etale homotopy types under specialization.
\end{corollary}
\begin{proof}
The failure and the group-level replacement were proved in Proposition~\ref{B:specialization}, using Stacks Theorem 58.30.3, tag 0C0R. The stronger homotopy statement is Haine--Holzschuh--Wolf \cite[Proposition 2.49]{BHHW2024}: compare each geometric fibre with its strict-local ``Milnor ball'' by proper base change, and compare the balls along a specialization by their nonabelian smooth base-change theorem. Both comparisons use only finite homotopy groups of orders invertible on the base. Their composite gives the equivalence after that completion. In characteristic five our computed $\mathbb Z/5$-characters show exactly why the omitted primary part cannot be restored.
\end{proof}

% Insertable historical clarification for B subjects 15--16.
% Prepared by /root/geometry_a567; B owns integration and final labels.
\paragraph{Historical scope of the infinitesimal part (source-verified).}
The quotation marks around the ``true'' fundamental group in subject~16
are mathematically significant.  In SGA~1, Expos\'e~X, Remarques~2.5,
p.~208 of the retyped edition (original p.~271), Grothendieck describes
the ``composante continue'' of the ``vrai'' fundamental group through
torsors whose structure groups are infinitesimal group schemes.
He then proposes a finite-flat group-scheme formulation for smooth
proper families.  Footnote~(2) on that same page explicitly records
that this proposal is false, citing an unpublished example of M.~Artin
with curve fibres of a fixed genus $g\geq2$
\cite{A567_SGA1_2003}.  This is a primary-source report of a
counterexample, not a reconstruction of Artin's unpublished example.
In particular, variation of the full \emph{\'etale} fundamental group
alone does not disprove the proposed finite-flat replacement.

The letter to Serre of 18~August~1959, published pp.~91--95, already
constructs a candidate degree-one pro-group from pointed torsors under
finite group schemes, including their infinitesimal parts.  Its
pp.~93--94 distinguish the rational-basepoint/descent issue and leave
the desired homotopy exact sequence and higher groups for future work
\cite{A567_GrothendieckSerre2004}.  Thus these sources identify a
substantial degree-one antecedent of subjects~15--16.  They do not
define all the $\pi_i$ of the arbitrary-topos theory requested in
subject~15.  Nor should that theory be identified without further
evidence with To\"en's later construction: To\"en explicitly distinguishes
his schematization from the construction in \emph{Pursuing Stacks},
which he considers closer to his affination
\cite[p.~11 and Appendix~A]{BToen2006}.

\begin{theorem}[Known modern degree-one replacement: Nori]\label{B:nori-finite-flat}
Let $k$ be algebraically closed, let $X/k$ be a smooth projective
connected curve, and fix $x\in X(k)$.  There is a profinite affine
$k$-group scheme $\pi^{\mathrm N}(X,x)$ such that, for every finite
$k$-group scheme $G$, homomorphisms
\[
  \pi^{\mathrm N}(X,x)\longrightarrow G
\]
classify isomorphism classes of pointed fppf $G$-torsors $(P,p)$ on
$(X,x)$, with the structure group $G$ fixed.  Restricting this
classification to finite \'etale $G$ recovers the usual pointed finite
\'etale torsors, so the maximal pro-\'etale quotient is the constant
pro-group scheme associated with $\pi_1^{\mathrm{\acute et}}(X,x)$.
\end{theorem}
\begin{proof}[Proof and source locator]
Use the neutral Tannakian category of essentially finite vector
bundles with the fibre functor $x^*$.  Nori constructs this category
and its dual on pp.~37--39 of \cite{A567_Nori1976}.
His Lemma~3.9 and the concluding remarks on p.~41 show that the
dual is an inverse limit of finite group schemes.  A pointed
$G$-torsor gives its associated-bundle tensor functor from
$\operatorname{Rep}_k(G)$ to this category; the point $p$ identifies
its fibre at $x$ with the forgetful functor.  Tannakian duality gives
the indicated homomorphism.  Nori's Proposition~2.9 identifies the
reconstructed torsor and its change of structure group, and
Proposition~3.11, pp.~40--41, proves existence and uniqueness of
the homomorphism with the specified point.  These are precisely
the hypotheses of the statement above, a restriction of Nori's
proper reduced connected perfect-field setting.

For the final assertion, a torsor under a finite \'etale group scheme
is finite \'etale because this is true after its fppf trivializing
base change and finite \'etale morphisms descend.  Conversely a
pointed finite \'etale Galois cover is its usual deck-group torsor.
Over algebraically closed $k$ finite \'etale group schemes are
constant.  The two classification properties therefore identify
all finite quotients and their transition maps, and hence their
inverse limits.
\end{proof}

\paragraph{What remains to identify.}
The finite-flat degree-one object is now specified.  What remains
historically unidentified in subject~15 is the intended category
and construction of the higher infinitesimal $\pi_i$, together with
the precise Hurewicz, exact-sequence and comparison statements.
This is a limitation of the recovered historical formulation,
not a claim that the modern degree-one theorem or To\"en's stated
theorems are unresolved.

% Insertable precise counterexample; B owns integration and final labels.
\begin{proposition}[Refuted base-change claim; explicit smooth counterexample]\label{B:nori-base-change-counterexample}
Let $k=\overline{\mathbf F}_2$, choose
$\alpha\in\mathbf F_4\subset k$ with $\alpha^2+\alpha+1=0$,
and let $X/k$ be the smooth projective model of
\[
 y^2+x(x+1)y=x(x+1)(x^3+\alpha x+\alpha).
\]
Choose any $x_0\in X(k)$ and put $K=\overline{k(t)}$.
Then $X$ is an ordinary curve of genus $2$, but the canonical map
\[
 \pi^{\mathrm N}(X_K,x_0)\longrightarrow
 \pi^{\mathrm N}(X,x_0)\times_k K
\]
is not an isomorphism.
\end{proposition}
\begin{proof}
In the notation of Laszlo--Pauly's ordinary genus-two normal form,
our parameters are $(a,b,c)=(1,1,\alpha)$.
They have $abc=\alpha\ne0$.  Section~2.3 and Lemma~2.1 of
\cite{A567_LaszloPauly2004}, read together with the displayed
description of the ordinary moduli space there, therefore give a
smooth ordinary genus-two curve.  Its function field may equivalently
be written
\[
 z^2+z=x+\alpha/x+1/(x+1),
 \qquad z=y/[x(x+1)]+\alpha;
\]
this equality follows from the defining relation of $\alpha$.
The parameter condition in Pauly's Proposition~4.1 is satisfied:
\[
 a^2+b^2+c^2+a+c
 =1+1+(\alpha+1)+1+\alpha=0.
\]
Pauly constructs an open subset $\Delta_0$ of a projective line in
the stable rank-two moduli space and, after an \'etale cover
$S\to\Delta_0$, a family of stable bundles $E$ with trivial
determinant and $(F_X^4)^*E\simeq\mathcal O_X^2$.
The classifying map $S\to\Delta_0$ is dominant
\cite[Proposition~4.1 and its complete proof, pp.~2--4]{A567_Pauly2007}.

Take an integral component of $S$ dominating $\Delta_0$.
Its function field is a finitely generated extension of $k$ of
transcendence degree one, so it embeds into $K$: choose a
transcendence basis and then embed its finite algebraic extension
in the algebraic closure $K$.  The corresponding $K$-valued
generic point yields a stable $F^4$-trivial bundle $E_K$ whose
coarse moduli point is not $k$-valued.  Consequently $E_K$ cannot
be the scalar extension of a bundle on $X$, since scalar extension
of such a bundle would have a $k$-valued moduli point.

We spell out the descent implication used in Pauly's Proposition~2.2.
The bundle $E_K$ is essentially finite: push its frame
$\mathrm{GL}_2$-torsor out along the entrywise fourth Frobenius
$F^4:\mathrm{GL}_2\to\mathrm{GL}_2$.  Its transition matrices
become those of $(F_{X_K}^4)^*E_K$, which is trivial.  A chosen
trivialization reduces the original frame torsor to the finite
infinitesimal group scheme $\ker(F^4)$.  Nori's
Proposition~3.8 in \cite{A567_Nori1976} therefore applies.

Suppose the displayed homomorphism were an isomorphism.  The
representation associated with $E_K$ would then be a representation
of $\pi^{\mathrm N}(X,x_0)_K$.  Its finitely many matrix coefficients
factor through $G_K$ for some finite quotient group scheme $G/k$,
since $\pi^{\mathrm N}(X,x_0)$ is an inverse limit of finite group
schemes.  This representation is simple: a proper nonzero
subrepresentation would give an essentially finite subbundle,
hence a degree-zero subbundle by Nori's Proposition~3.8,
contradicting stability of the degree-zero bundle $E_K$.

Put $H=\mathcal O(G)^\vee$, the finite-dimensional dual Hopf
algebra.  Since $k$ is algebraically closed, its semisimple quotient
is a finite product of matrix algebras over $k$.  If $J$ denotes
its Jacobson radical, then $J\otimes_kK$ is nilpotent and
\[
 (H\otimes_kK)/(J\otimes_kK)
       \simeq\prod_i M_{n_i}(K).
\]
Consequently $J\otimes_kK$ is exactly the radical of
$H\otimes_kK$, and every simple $H\otimes_kK$-module is the
extension of a simple $H$-module.  The simple representation above
therefore descends to $G/k$.  Its associated bundle on $X$
pulls back to $E_K$, by the defining compatibility of the
canonical base-change homomorphism.  This contradicts the
nonconstant coarse moduli point and proves the assertion.
The field and polynomial substitutions are additionally checked
exactly by \texttt{code/A567\_Nori\_curve.py}, with complete output
\texttt{logs/A567\_Nori\_curve.json}; the computation does not
replace the geometric theorems just invoked.
\end{proof}

\paragraph{Precise historical consequence.}
This is a complete counterexample to arbitrary algebraically closed
base-field extension for the finite-flat torsor fundamental group,
even for smooth proper ordinary curves.  It is a different statement
from prime-to-characteristic \'etale specialization and from the
unpublished Artin family example mentioned in SGA~1.  No identification
of these counterexamples is asserted.  The explicit transcendental
extension above also avoids the overbroad wording ``any'' extension
in Pauly's Corollary~4.2: that printed proof requires a parameter
not defined over $k$, and cannot include the identity extension $K=k$.

\subsubsection{17. Separation and representability of relative Picard functors}
\begin{theorem}[Known: Picard representability over an arbitrary base]\label{B:Picard-general}
Let $f:X\to S$ be a proper flat finitely presented morphism of algebraic spaces. Assume that $\mathcal O_T\to f_{T*}\mathcal O_{X_T}$ is an isomorphism for every scheme $T\to S$. Then the fppf sheaf $\operatorname{Pic}_{X/S}$ is an algebraic space, quasi-separated and locally of finite presentation over $S$. It is separated if all geometric fibres of $f$ are integral, and smooth if $f$ has relative dimension at most one. No noetherian or one-dimensional hypothesis on the base is required.
\end{theorem}
\begin{proof}[Derivation]
Stacks Proposition~99.11.8, tag~0D2C, has exactly the properness, flatness, finite-presentation and universal-functions hypotheses stated here. Its construction is worth making explicit: after the faithfully flat base change $X\to S$, the family has a section. Rigidify each invertible sheaf along that section. Universal equality of global functions makes every automorphism preserving the rigidification the identity, so the rigidified Picard stack is an algebraic space. Descent of algebraic spaces removes the chosen section. This proves representability and explains why its output is an algebraic space.

Quasi-separatedness and local finite presentation are Stacks Lemma~108.9.3; Lemma~108.9.4 applies the valuative criterion under geometric integrality to give separatedness. For smoothness, Lemmas~108.8.5 and~108.9.5 use the following obstruction calculation. For a square-zero thickening of affine base schemes, an invertible sheaf lifts if its obstruction in $H^2(X_T,\mathcal I)$ vanishes. Relative dimension at most one gives $R^qf_{T*}\mathcal I=0$ for $q>1$, and the affine base annihilates higher cohomology of the remaining quasi-coherent direct images. Thus $H^2=0$. The Picard stack is formally smooth and locally of finite presentation, hence smooth; its smooth surjective $\mathbf G_m$-gerbe over the Picard space transfers smoothness to that space. Every condition used remains valid over the arbitrary base allowed in the statement.
\end{proof}

\begin{theorem}[Known: Raynaud's precise curve criterion]\label{B:subject17}
Let $R$ be a strictly henselian discrete valuation ring with perfect residue field, and let $f:X\to\operatorname{Spec}R$ be a proper flat regular model of a smooth projective geometrically connected curve. Write the special fibre as $\sum_jm_jC_j$, and put $\delta=\gcd_jm_j$. Then
\[
 \delta=1\quad\Longleftrightarrow\quad
 \operatorname{Pic}^0_{X/R}\text{ is representable by a scheme}
 \quad\Longleftrightarrow\quad
 \operatorname{Pic}^0_{X/R}\text{ is separated}.
\]
Moreover the appropriate separated quotient of the closure of the generic Picard identity component is the N\'eron model of the generic Jacobian.
\end{theorem}
\begin{proof}[Derivation]
Raynaud \cite[Theorem 8.2.1]{BRaynaud1970} lists eight conditions, including the three displayed, and proves all eight equivalent when its total multiplicity $\delta$ equals $d'$. His sufficient hypotheses for that equality are local factoriality of $X$ and perfect residue field. Regularity gives local factoriality; over the strictly henselian perfect residue field, total multiplicity is the displayed ordinary multiplicity. His condition $(N)^*$, defined in 6.1.4, requires no embedded components in the special fibre, normality of the \emph{total space} at the generic points of that fibre, and $f_*\mathcal O_X=\mathcal O_R$. The first two follow from regularity and flatness over a trait; no normality of a nonreduced special fibre is asserted. The last follows from geometric connectedness of the smooth generic fibre, properness, and normality of $R$ via Stein factorization. Thus every hypothesis of 8.2.1 is checked.

For the final assertion use Raynaud's Theorem 8.1.4. The separated quotient $Q$ is formed by dividing by the schematic closure of the identity of the generic fibre, and $Q'$ is the closure of its generic identity component. The same local factoriality and perfect-residue assumptions make $Q'$ a finite-type N\'eron group. Its generic fibre is the Jacobian, so uniqueness from subject~11 identifies it with the N\'eron model. The theorem thereby gives both a representability criterion and the correction obtained when separation fails.
\end{proof}

\begin{proposition}[Refuted: scheme representability for every regular curve family]\label{B:Picard-double-fibre}
Let $R=\mathbb C[[t]]$, $R'=\mathbb C[[s]]$ with $t=s^2$, and let $E/\mathbb C$ be the elliptic curve $y^2=x^3-x$, with origin $O$ and two-torsion point $T=(0,0)$. The involution
\[
 (P,s)\longmapsto(P+T,-s)
\]
on $Y=E\times\operatorname{Spec}R'$ has a scheme quotient $X=Y/C_2$. The morphism $X\to\operatorname{Spec}R$ is proper, flat and finitely presented, its total space is regular, and its generic fibre is a smooth geometrically connected genus-one curve. Nevertheless $\operatorname{Pic}^0_{X/R}$ is not separated and is not representable by a scheme. The full Picard functor is represented by an algebraic space.
\end{proposition}
\begin{proof}
The involution fixes $t=s^2$ and squares to the identity. Translation by the nonzero point $T$ has no fixed geometric point, so the action is free. Since $Y$ is projective over $R$, Stacks Lemma~66.14.2, tag~07S5, constructs the scheme quotient and makes $q:Y\to X$ a finite \'etale $C_2$-torsor. The quotient is separated by Lemma~66.14.3. Finite presentation and flatness descend along $q$, and properness follows because the image of every closed subset of a base change of $X$ is also the image of its inverse image under the finite surjection from the proper scheme $Y$. Regularity is \'etale-local and $Y$ is smooth over the regular ring $R'$. After the field extension $\mathbb C((t))\subset\mathbb C((s))$, Galois descent identifies the generic fibre with $E_{\mathbb C((s))}$, proving its asserted geometric properties.

Put $D=E\times\{s=0\}$. Its quotient is the smooth irreducible curve $C=E/\langle T\rangle$. Since $q$ is \'etale, the equality $q^*(\operatorname{div}_X(t))=2D$ gives $\operatorname{div}_X(t)=2C$: there is a single special-fibre component, with multiplicity two. Consequently the integer $\delta$ of Theorem~\ref{B:subject17} is $2$. Its checked hypotheses hold because $R$ is strictly henselian with perfect residue field and $X$ is regular. That theorem proves both failure of separation and failure of scheme representability of $\operatorname{Pic}^0$.

For the final assertion, take any $R$-algebra $A$. Functions on $Y_A$ are $A\otimes_R R'=A\oplus As$, since a proper geometrically connected elliptic curve has only base functions, universally. Descent identifies functions on $X_A$ with the $C_2$-invariants; they are exactly $A$, as $2$ is invertible and the involution negates $s$. This verifies the universal-functions condition of Theorem~\ref{B:Picard-general}, which gives the Picard algebraic space. Finally, on the chart $x\ne0$ translation is $(x,y)\mapsto(-1/x,y/x^2)$. The program verifies preservation of the curve equation, the involution identity and the empty fixed-point ideal; the two omitted points $O,T$ are exchanged. These exact symbolic checks supplement the quotient and multiplicity proof.
\end{proof}

\subsubsection{18. Effective flat descent and quotients}
\begin{proposition}[Known: effective affine descent, including Artinian bases]\label{B:subject18}
For any fpqc covering $\{S_i\to S\}$, every descent datum on schemes affine over $S_i$ is effective, with no noetherian assumption. In particular this applies over an Artinian base. If a finite locally free group scheme $G/S$ acts freely on an $S$-scheme $X$, the fppf quotient is an algebraic space; when that quotient is represented by a scheme, the descended $G$-torsor and its affine descent data are effective.
\end{proposition}
\begin{proof}[Derivation of the affine assertion]
Stacks Lemma 35.37.1, tag 0245, reduces to a faithfully flat ring map $R\to A$ and an $A$-algebra $B$ equipped with a cocycle isomorphism between its two pullbacks. Faithfully flat module descent gives an $R$-module $C$ with $A\otimes_RC\simeq B$. The equalizer description of $C\subset B$ shows it is closed under multiplication and contains $1$, because the cocycle is an algebra map. Thus $C$ is an $R$-algebra and $\operatorname{Spec}C$ supplies the descent. This proof works unchanged when $R$ is Artinian.

For the quotient assertion use Stacks Theorem 80.10.1, tag 04S6. The free action identifies the relation with $G\times_SX\rightrightarrows X$; the fppf quotient of a flat locally finitely presented equivalence relation is an algebraic space, and its pullback to $X$ is the torsor relation. The scheme conclusion also follows if $X$ is quasi-projective over $S$, by Stacks Lemma 66.14.2, tag 07S5; the displayed algebraic-space statement does not assume quasi-projectivity. The original word on the scan is ``artinien''; the affine result is also valid if that word were read as ``arbitraire''.
\end{proof}

\begin{proposition}[Refuted: unrestricted flat descent over an Artinian base]\label{B:descent-counterexample}
Let $k=\mathbb Q$, $L=\mathbb Q(i)$, and let $X/L$ be the affine line with doubled origin: glue two copies $U_+$ and $U_-$ of $\operatorname{Spec}L[t]$ by the identity on $D(t)$. The involution that conjugates $i$, fixes $t$, and interchanges the two charts gives a descent datum for the finite \'etale cover $\operatorname{Spec}L\to\operatorname{Spec}k$. This datum is not effective in schemes, although $X$ is smooth and of finite type over $L$ and the base $k$ is Artinian. Its quotient is an algebraic space that is not a scheme.
\end{proposition}
\begin{proof}
The two chart maps agree on the gluing open, and their composite squares to the identity. Thus they define a semilinear action of $\operatorname{Gal}(L/k)=C_2$. The decomposition $L\otimes_kL\simeq L\times L$ translates this action into a descent datum; its cocycle equation is exactly the group law. The program checks the four products in $C_2$, the permutation of the two charts, and the coefficient action $a+bi\mapsto a-bi$. Smoothness and finite type are local on $X$, where both follow from the two affine-line charts.

Suppose $Y/k$ were a scheme with this descent datum as its base change. The two origins of $X$ would map to the same point $y\in Y$, since they lie in one Galois orbit. An affine open neighborhood $W$ of $y$ would pull back to an affine open $W_L\subset X$ containing both origins. No such open exists. Indeed the two chart inclusions give maps
\[
 \operatorname{Spec}L[t]_{(t)}\rightrightarrows X
\]
that agree at the generic point and differ at the closed point. Both maps factor through every open containing both origins, since this trait has only its generic and closed points. An affine $L$-scheme is separated, so uniqueness in the valuative criterion forbids this pair of maps into an affine open. This contradiction proves noneffectivity.

The graph relation of the semilinear action is free as a relation over $k$: the nontrivial Galois automorphism fixes no geometric point of $\operatorname{Spec}L$. Its two projections are finite \'etale, so Stacks Theorem 80.10.1 supplies an algebraic-space quotient. Descent along $\operatorname{Spec}L\to\operatorname{Spec}k$ identifies its base change with $X$. The preceding contradiction proves that this quotient is not a scheme. Nonseparatedness of $X$ is part of this example and is not suppressed.
\end{proof}

\begin{proposition}[Proved: an exact criterion for finite Galois descent]\label{B:descent-criterion}
Let $L/k$ be a finite Galois extension and let an $L$-scheme $X$ carry a semilinear action of $G=\operatorname{Gal}(L/k)$. Its descent datum is effective in schemes if and only if $X$ has an open covering by $G$-stable affine subschemes.
\end{proposition}
\begin{proof}
Necessity follows by pulling back an affine cover of a descended scheme: finite field extension is affine and every such inverse image is $G$-stable. For sufficiency, descend each stable affine $U_\alpha$ by the algebra descent proved in Proposition~\ref{B:subject18}, obtaining $V_\alpha/k$. Each intersection $U_\alpha\cap U_\beta$ is $G$-stable. Its image in $V_\alpha$ is open, because the finite \'etale surjection $U_\alpha\to V_\alpha$ is open; its full inverse image is the original intersection, because the fibres are Galois orbits. Thus the intersection descends as an open subscheme of $V_\alpha$ and of $V_\beta$. The identity between the two pullbacks descends uniquely by fpqc descent of morphisms (Stacks Lemma 35.35.4); the cocycle on triple intersections descends by the same uniqueness. Gluing these open subschemes gives a scheme $Y/k$ whose base change, with its action, is $X$. This criterion specifies the obstruction in the preceding example without asserting that arbitrary descent data are effective.
\end{proof}

\subsubsection{19. Positivity for cycles of arbitrary codimension}
\begin{theorem}[Known: a positive cone of higher-codimension classes]\label{B:subject19}
For a projective variety $X$ over an algebraically closed field, let $N^k(X)$ be the space of numerical codimension-$k$ classes. The pliant cone $\operatorname{PL}^k(X)$, generated and closed by products of Schur classes of globally generated vector bundles, is full dimensional and salient; ample complete intersections lie in its interior. These cones are stable under pullback, lie in the pseudoeffective cones and are stable under products. The pseudoeffective cone of $k$-cycles is full dimensional and salient, and degree against an ample class extends to a norm on the numerical cycle space.
\end{theorem}
\begin{proof}[Derivation]
These are Fulger--Lehmann \cite[Theorems 1.3 and 1.4]{BFulgerLehmann2017}. Globally generated bundles give maps to Grassmannians, where the Schur classes are effective; their products explain pullback stability and positivity. The theorem proves the nontrivial assertions of interior and salience, giving a well-defined higher-codimension positivity theory with the stated functorial properties. Since $k=1$ recovers the usual nef-divisor setting and arbitrary $k$ is permitted, it is a concrete answer to the dimensional extension requested in subject~19. The sheet does not specify axioms selecting a unique meaning of ``amplitude'' in higher codimension; the theorem does not claim that all possible positive cones coincide.
\end{proof}

\subsubsection{20. Vanishing on homogeneous spaces}
\begin{theorem}[Known: Kempf vanishing, with the Frobenius derivation]\label{B:subject20}
Let $G$ be a split semisimple simply connected algebraic group over a perfect field, $B$ a Borel subgroup, and $\lambda$ a dominant weight. With the convention that dominant weights define globally generated line bundles on $G/B$,
\[
 H^q(G/B,\mathcal L_\lambda)=0\qquad(q>0).
\]
The assertion holds in every characteristic.
\end{theorem}
\begin{proof}[Derivation from the opened proof]
Haboush \cite[Theorems 2.1 and 3.1]{BHaboush1980} gives the following characteristic-$p$ proof. Write $\rho$ for the half-sum of positive roots and $\operatorname{St}_r$ for the Steinberg module. The Frobenius pushforward identity is
\[
 F^r_*\mathcal L_{(p^r-1)\rho+p^r\lambda}
 \simeq\operatorname{St}_r\otimes\mathcal L_\lambda.
\]
Its proof identifies the fibre module with the Steinberg representation and checks the rank. The line bundle on the left is
$\mathcal L_{-\rho}\otimes\mathcal L_{\lambda+\rho}^{\otimes p^r}$.
Because $\lambda+\rho$ is strictly dominant, $\mathcal L_{\lambda+\rho}$ is ample, and Serre vanishing kills its higher cohomology for sufficiently large $r$. Frobenius is finite, hence has no higher direct images. Taking cohomology of the identity gives $\operatorname{St}_r\otimes H^q(G/B,\mathcal L_\lambda)=0$. The Steinberg module is nonzero, so the latter cohomology group is zero. Haboush's Remark 3.2 obtains characteristic zero by a Chevalley model over $\mathbb Z$ and semicontinuity: vanishing at a positive-characteristic fibre gives vanishing at the generic fibre. This verifies the precise homogeneous-space vanishing statement rather than invoking an unspecified general vanishing principle.
\end{proof}

\subsubsection{21. Affine base change and tubular neighborhoods}
\begin{corollary}[Known, with the two comparisons distinguished]\label{B:subject21}
For a quasi-compact quasi-separated closed pair, degree-zero comparison for every \'etale sheaf of sets is equivalent to geometric henselianity. In the affine case it implies comparison in every degree for every torsion abelian \'etale sheaf; the same conclusion holds in the proper situation of Theorem~\ref{B:henselian-general}, and for every separated finite-type scheme over a henselian discrete valuation ring by Theorem~\ref{C:henselian-trait-comparison}. If $Z$ is affine and $X$ is quasi-separated, Proposition~\ref{C:henselian-affine-closed-fibre} forces $X$ to be affine and gives the same conclusion. Theorem~\ref{C:henselian-noetherian-pair-comparison} proves the torsion comparison for every separated finite-type scheme over any noetherian henselian ring pair $(A,I)$ with $Z=X\times_A A/I$, assuming only that every closed point of $X$ lies in $Z$. Corollary~\ref{C:henselian-noetherian-closedpoint-criterion} proves that, in this class, the closed-point condition is equivalent to universal degree-zero comparison. No flatness, locality, catenarity, excellence or properness of $X$ is required. Theorem~\ref{C:henselian-horizontal-proper} identifies the proper horizontal components; Theorem~\ref{B:henselian-local-flat-proper} remains a separate proof for flat positive-dimensional local bases. Completeness and henselization give the two affine cases named on scan~4. Completion of a punctured henselian pair with finitely generated ideal gives a second, distinct comparison. The affine and punctured statements extend to torsion sheaves of spaces, and consequently to the associated profinite \'etale homotopy types. The assertion with arbitrary abelian coefficients is false even for the constant integer sheaf on an affine henselian pair, by Proposition~\ref{B:constant-integer-rigidity}; Proposition~\ref{B:nontorsion-rigidity} supplies a separate coherent-sheaf example; the torsion implication for arbitrary closed pairs is also false: Theorem~\ref{C:subject21-torsion-counterexample} constructs a qcqs nonseparated nonnoetherian pair with universal degree-zero comparison but failure of degree-one comparison for every constant $\mathbf Z/n\mathbf Z$, $n>1$. It realizes the supported-degree-two obstruction of Proposition~\ref{B:H1-gate} and the cyclic-torsor failure characterized by Proposition~\ref{C:torsion-degree-two-reduction}.
\end{corollary}
\begin{proof}
Theorems~\ref{B:henselian} and~\ref{B:henselian-general}, Propositions~\ref{B:H0-counterexample}, \ref{B:constant-integer-rigidity}, \ref{B:nontorsion-rigidity}, \ref{B:H1-gate}, Theorem~\ref{C:henselian-trait-comparison}, Corollary~\ref{C:henselian-proper-closed-fibre}, Proposition~\ref{C:henselian-affine-closed-fibre}, Theorems~\ref{B:henselian-local-flat-proper} and~\ref{C:henselian-noetherian-pair-comparison}, Corollary~\ref{C:henselian-noetherian-closedpoint-criterion}, and Theorem~\ref{B:Fujiwara} prove the degree-zero, affine, proper, separated finite-type henselian-base, coefficient, and punctured assertions. Lemma~\ref{B:henselized-affine-cover} and Proposition~\ref{B:henselized-Cech-obstruction} explain why the natural affine-henselization descent cover does not automatically prove the general assertion. Theorem~\ref{C:subject21-torsion-counterexample} supplies its actual torsion counterexample. Haine--Holzschuh--Wolf \cite[Corollaries 2.34 and 2.36]{BHHW2024} establish the homotopical extension by reducing to groupoid-valued sheaves and torsion abelian cohomology. Applying it to constant finite homotopy types gives the conclusion about profinite shapes, as in subject~15. Their proof identifies the exact original Gabber locator as \cite[\S5, Corollary 1]{BGabber1994}, and Fujiwara's as \cite[Corollary 6.6.4]{BFujiwara1995}. We obtained the complete modern proofs, but not the required theorem passages in the specified original articles; the access distinction is recorded in the bibliography. An adequate original preview would suffice for the source check, but none obtained contains the required statements and hypotheses.
\end{proof}

\subsubsection{22. Tensor categories, signs, and gerbes}
\begin{theorem}[Known: the precise Tannakian reconstruction]\label{B:subject22}
Let $k$ be a field and $\mathcal C$ a rigid abelian $k$-linear symmetric tensor category with $\operatorname{End}(\mathbf1)=k$. If there is an exact faithful $k$-linear symmetric tensor functor $\omega:\mathcal C\to\operatorname{Vec}_k$, then
\[
 G=\operatorname{Aut}^{\otimes}(\omega)
\]
is an affine group scheme and $\mathcal C\simeq\operatorname{Rep}_k(G)$. Its other fibre functors over a $k$-algebra are classified by $G$-torsors. If such a fibre functor exists only after a field extension, the stack of fibre functors is an affine gerbe and its representation category recovers $\mathcal C$.
\end{theorem}
\begin{proof}[Derivation]
The neutral theorem is Deligne--Milne \cite[Theorem 2.11]{BDeligneMilne1982}: reconstruction first gives the coalgebra of natural endomorphisms of the fibre functor; tensor product makes it a bialgebra, and rigidity gives the antipode. Thus the representing affine monoid is a group, and the reconstructed comodules are precisely the original objects. Theorem 3.2 identifies $\operatorname{Isom}^{\otimes}(\omega,\eta)$ with the torsor corresponding to $\eta$, proving the second assertion.

The gerbe statement is Theorem 3.9 in the opened corrected author edition, which explicitly invokes Deligne (1990), 1.12--1.13. Its equivalence is obtained by evaluation at fibre functors and descent along their fpqc-local isomorphisms. That edition expressly says this theorem was unavailable in the original 1982 chapter and points out the defect in Saavedra's earlier formulation. We therefore do not attribute the corrected general theorem solely to the marginal note ``Saavedra en train''.
\end{proof}
\begin{proposition}[Refuted: rigidity alone gives ordinary group representations]\label{B:super}
If $\operatorname{char}k\ne2$, finite-dimensional super vector spaces with the Koszul symmetry form a rigid abelian symmetric tensor category with endomorphism ring $k$ of the unit, but admit no faithful symmetric tensor functor to ordinary vector spaces. Thus the sign and fibre-functor hypotheses cannot be removed from subject~22.
\end{proposition}
\begin{proof}
Let $L$ be the one-dimensional odd super vector space. Its tensor square is the unit and its self-symmetry is $c_{L,L}=-\operatorname{id}$. Any tensor functor to ordinary vector spaces must carry $L$ to a one-dimensional vector space, because its dimension has square one. The ordinary symmetry on its tensor square is $+\operatorname{id}$. Symmetric compatibility would identify these two maps. Faithfulness and $2\ne0$ rule this out. The remaining categorical hypotheses hold componentwise on the even and odd summands; duals are the usual graded duals and the Koszul sign gives their evaluation symmetry. The attached exact calculation checks the two signs. The corrected theorem is the preceding reconstruction with a symmetric fibre functor; using super fibre functors gives a different reconstruction problem.
\end{proof}

\subsubsection{23. Chow groups and the regular-target issue}
\begin{proposition}[Known: integral smooth-target and rational regular-scheme results]\label{B:subject23}
For schemes of finite type over a field, ordinary Chow homology admits proper pushforward and flat pullback. If $Y$ is smooth quasi-projective of pure dimension $e$, $X$ is quasi-projective and pure of dimension $d$, and $f:X\to Y$ is any morphism, there is an integral pullback on classes
\[
 f^*: \mathrm{CH}^r(Y)\longrightarrow\mathrm{CH}_{d-r}(X),
 \qquad a\longmapsto f^*a\cap[X].
\]
There is no smoothness assumption on $X$. For a separated regular noetherian scheme $Y$ of finite Krull dimension, rational Chow groups with supports have a canonical intersection product, agreeing with the Serre Tor formula when cycles meet properly. In particular $\mathrm{CH}^*(Y)_{\mathbb Q}$ is a graded commutative ring.
\end{proposition}
\begin{proof}[Derivation and exact scope]
Fulton \cite[\S\S1--3, especially the moving lemma in \S2.3 and the cap product in \S3.1]{BFulton1975} constructs the integral theory. Move a representative of $a$ to intersect the finitely many fibre-dimension strata of $f$ properly, then intersect it with $[X]$ along $f$. The moving lemma and the identities of \S2.2 make the result independent of the moved representative and of rational equivalence. Equivalently the graph of $f$ is a regular embedding of codimension $e$ in $X\times Y$, as a section of a smooth projection, so its refined Gysin operation on $[X]\times a$ has degree $d-r$. When $X$ is also smooth this is the usual contravariant Chow-ring map.

For the second assertion, the original Gillet--Soul\'e \cite[Theorems 8.2 and 8.3 (Theorem C), pp.~274--276]{BGilletSoule1987}, also recalled in \cite[\S4.1.1]{BGilletSoule1990}, supplies the general regular-scheme theorem: for every closed support $Z\subset Y$, the cycle class gives
\[
 \mathrm{CH}^r_Z(Y)_{\mathbb Q}
 \simeq\operatorname{Gr}^r_\gamma K_0^Z(Y)_{\mathbb Q}.
\]
Multiplication of perfect complexes with supports maps
$K_0^{Z_1}(Y)\otimes K_0^{Z_2}(Y)\to K_0^{Z_1\cap Z_2}(Y)$ and respects the gamma filtration. Transporting the product through these isomorphisms gives the asserted product; associativity, commutativity and the unit are inherited from tensor product. Their theorem identifies proper intersections with alternating Tor lengths. This addresses the arbitrary regular-scheme direction with rational coefficients. It does not establish the sheet's strongest possible integral formulation for every regular target over every base: smoothness over a field and tensoring with $\mathbb Q$ are distinct scope restrictions, explicitly retained.
\end{proof}

\begin{proposition}[Proved: rational pullback from a regular target to an arbitrary source]\label{B:regular-pullback}
Let $Y$ be a separated regular noetherian scheme of finite Krull dimension. Give it the dimension function $\delta_Y(y)=-\operatorname{codim}_Y(y)$, component by component, and give finite-type $Y$-schemes the induced dimension function. There is a canonical graded ring injection
\[
 \mathrm{CH}^*(Y)_{\mathbb Q}\longrightarrow A^*(Y)_{\mathbb Q},
 \qquad \alpha\longmapsto\widetilde\alpha,
\]
into operational Chow cohomology, satisfying $\widetilde\alpha\cap[Y]=\alpha$. Consequently every morphism $f:X\to Y$ of finite type, with no regularity assumption on $X$, has canonical operations
\[
 f^*\widetilde\alpha\cap-:
 \mathrm{CH}_j(X)_{\mathbb Q}\longrightarrow
 \mathrm{CH}_{j-p}(X)_{\mathbb Q}
 \quad(\alpha\in\mathrm{CH}^p(Y)_{\mathbb Q}).
\]
They commute with proper pushforward, flat pullback and refined Cartier-divisor intersection, and are functorial under further pullback as operational classes. If $X$ is pure of dimension $d$ for the chosen dimension function, evaluating on $[X]$ gives a rational class in $\mathrm{CH}_{d-p}(X)$. Thus the regular-target assertion has an arbitrary-source rational form, including the ample-line-bundle setting of the manuscript.
\end{proposition}
\begin{proof}
A regular noetherian scheme is Cohen--Macaulay, so Stacks Example~42.7.4, tag~0F91, proves universal catenarity and verifies this dimension function. Separatedness gives an affine diagonal and finite Krull dimension makes the function bounded. Stacks Lemma~42.58.1, tag~0FEY, therefore says that
\[
 \theta:K_0(\operatorname{Vect}(Y))_{\mathbb Q}
 \xrightarrow{\operatorname{ch}}A^*(Y)_{\mathbb Q}
 \xrightarrow{\cap[Y]}\mathrm{CH}^*(Y)_{\mathbb Q}
\]
is an isomorphism of vector spaces. For $\alpha$ of codimension $p$, let $v=\theta^{-1}(\alpha)$ and define $\widetilde\alpha=\operatorname{ch}(v)$. By Lemma~42.56.6, the Adams operator $\psi^2$ multiplies $\operatorname{ch}_j$ by $2^j$. Hence $\theta(2^{-p}\psi^2v)=\alpha$, so uniqueness implies $\psi^2v=2^pv$. Therefore
\[
 (2^j-2^p)\operatorname{ch}_j(v)=0,
\]
and over $\mathbb Q$ the operational class has only degree $p$. Its evaluation on $[Y]$ is $\alpha$, proving injectivity. The product in Stacks Section~42.58 is precisely the product transported from $K_0$ by $\theta$; multiplicativity of the Chern character proves that our injection preserves it.

For an arbitrary $f:X\to Y$, pull back the vector bundles representing $v$. The resulting operational class is $\operatorname{ch}(f^*v)$, so it is canonical and has the stated degree. Operational Chow classes are defined as families of homomorphisms on Chow groups of schemes over their base, compatible with the three listed operations; these compatibilities and their pullback follow directly from that definition and from the Chern-class construction. Since pullback of vector bundles is associative, further pullback of this class is associative as well. No inverse Chern-character map on the potentially singular scheme $X$ is used. All inverses and the division by $2^j-2^p$ occur over $\mathbb Q$; the proof makes no integral conclusion.
\end{proof}

\subsubsection{24. The two equivalence questions}
\begin{corollary}[Refuted for the first question; proved in the requested fourfold classes]\label{B:subject24}
Numerically trivial one-cycles on a smooth projective connected threefold need not be $\tau$-trivial. Numerical and rational-coefficient homological equivalence do agree for codimension-two cycles on every smooth projective fourfold in characteristic zero, and for abelian fourfolds in positive characteristic for every permitted $\ell$, using the 2026 Tate theorem stated in Theorem~\ref{B:abelian}.
\end{corollary}
\begin{proof}
The precise counterexample is Theorem~\ref{B:tau}. The fourfold proofs are Theorems~\ref{B:betti}, \ref{B:charzero} and \ref{B:abelian}. Theorem~\ref{B:criterion} identifies the exact positivity criterion and a further class dominated by abelian fourfolds. No claim is made that these arguments establish standard conjecture $D$ for all positive-characteristic fourfolds. Ancona's numerical positivity alone does not remove that gate, as Proposition~\ref{B:ancona-scope} proves.
\end{proof}

\subsubsection{25. Products and the Albanese kernel}
\begin{corollary}[Refuted for both questions, with an explicit replacement]\label{B:subject25}
The product of two algebraically trivial cycles need not be rationally trivial, and an Albanese-trivial zero-cycle on a surface need not be rationally trivial. Both failures occur for the single explicit surface and cycles in Theorem~\ref{B:square}; the zero-cycle has infinite order. On a product of elliptic curves, the Albanese kernel is generated by the exterior products of degree-zero divisors.
\end{corollary}
\begin{proof}
The square $D^2=2z\ne0$, the nontorsion argument from Mumford's trace forms and the zero Albanese image were proved in Theorem~\ref{B:square}. Generation and the valid algebraic-equivalence ideal statement are Proposition~\ref{B:alb-kernel}. Thus the two questions and the negative marginal note have each been checked, rather than inferred from the title of Mumford's paper.
\end{proof}

\subsubsection{26. Relative Rosenlicht theory}
\begin{theorem}[Known: the relative generalized Jacobian universal property]\label{B:subject26}
Let $S$ be a scheme, $X\to S$ a smooth projective curve with geometrically connected fibres, and $D\subset X$ a relative effective Cartier divisor finite and faithfully flat over $S$. Put $D_n=V(\mathcal I_D^{n+1})$, and let $J_n$ be the relative Picard sheaf of the pinching $X\amalg_{D_n}S$, equivalently line bundles on $X$ trivialized along $D_n$. Let $J=\varprojlim_nJ_n$. The diagonal ideal supplies an Abel map $a:X\setminus D\to J$, and for every smooth commutative $S$-group scheme $G$,
\[
 \operatorname{Hom}_{S\text{-}\mathrm{grp}}(J,G)
 \xrightarrow{\ \circ a\ }\operatorname{Hom}_S(X\setminus D,G)
\]
is bijective. There is an exact sequence of fpqc sheaves
\[
 1\longrightarrow\mathbb G_{m,S}
 \longrightarrow (D_n/S)_*\mathbb G_{m,D_n}
 \longrightarrow J_n\longrightarrow\operatorname{Pic}_{X/S}
 \longrightarrow1.
\]
\end{theorem}
\begin{proof}[Derivation]
These are Contou-Carr\`ere \cite[\S1.2, (1.2.1)--(1.2.4)]{BContouCarrere2013}, which explicitly restate the global relative theorem of \cite{BContouCarrere1990}. The 2013 full text was opened; the 1990 chapter itself was not. Pinching gives a distinguished section, so the Picard functor is described by pairs $(L,\alpha:L|_{D_n}\simeq\mathcal O_{D_n})$. Changing a trivialization multiplies it by a unit of $D_n$, and the automorphisms induced by global units of $X$ are exactly units of $S$ because the smooth proper fibres are geometrically connected. This gives the first three terms and the kernel in the exact sequence. A line bundle becomes trivial on the finite $S$-scheme $D_n$ fpqc locally on $S$, giving the final sheaf surjectivity.

The diagonal ideal on $X\times_S(X\setminus D)$ is canonically trivial on $D_n\times_S(X\setminus D)$, and hence defines compatible Abel maps. The cited universal property applies for arbitrary $S$, with the displayed group hypotheses, giving the bijection. The use of the whole Picard object, including its degree components, records constant maps as well as pointed ones; replacing it silently by its degree-zero part would change the statement.
\end{proof}

\subsubsection{27. Higher ad\`eles and local Jacobians}
\begin{theorem}[Known: the two precise constructions]\label{B:subject27}
For a noetherian scheme $X$ and a quasi-coherent sheaf $F$, the Parshin--Beilinson ad\`ele sheaves over flags give a functorial flasque resolution
\[
 0\longrightarrow F\longrightarrow\mathcal A_X^0(F)
 \longrightarrow\mathcal A_X^1(F)\longrightarrow\cdots.
\]
Its global-sections complex computes $R\Gamma(X,F)$. For a noetherian ring $A$, a formal curve $\mathfrak X=\operatorname{Spf}B$ with $B\simeq A[[t]]$, and $U=\operatorname{Spec}B[1/t]$, define the local Jacobian as the Zariski group sheaf with
\[
 \mathcal J(R)=\bigl(R[[t]][t^{-1}]\bigr)^\times
 \qquad(R\text{ an }A\text{-algebra}).
\]
There is an Abel map $a:U\to\mathcal J_{\mathrm{omb}}$ to its formal shadow such that, for every smooth separated commutative $A$-group scheme $G$,
\[
 \operatorname{Hom}_{A\text{-}\mathrm{grp}}(\mathcal J,G)
 \xrightarrow{\ h\mapsto h_{\mathrm{omb}}\circ a\ }G(U)
 \quad\text{is bijective}.
\]
\end{theorem}
\begin{proof}[Derivation]
The ad\`ele resolution is Beilinson's theorem, stated precisely in Braunling--Groechenig--Wolfson \cite[Theorem 2.1]{BBGW2016}, with its detailed proof attributed there to Huber. The terms are formed recursively by completing and localizing along chains of specializations, with the restricted-product conditions; the differential is the alternating sum of face maps. Flasque sheaves are acyclic for global sections, so the asserted computation follows immediately from the resolution. For a smooth integral proper curve over a field and $F=\mathcal O_X$, the global-sections complex is
\[
 K(X)\oplus\prod_x\widehat{\mathcal O}_{X,x}
 \longrightarrow\prod_x'K_x,
\]
recovering the ordinary finite-place ad\`eles. If $X$ is reduced, purely $n$-dimensional and of finite type over a field $k$, the opened paper's Theorem 4.17 identifies the Tate and Beilinson $n$-fold cubical endomorphism algebras of a saturated flag. When $k$ is perfect, it also identifies these with the Yekutieli algebra on each higher-local-field factor. These are the additional hypotheses and factor conventions for that higher local structure.

The local Jacobian statement is Contou-Carr\`ere \cite[Theorems 1.4.4 and 1.6.6]{BContouCarrere2013}. We specify the shadow to fix the type of every arrow. Lemmas 0.7--0.8 and (0.9)--(0.13) uniquely factor a local unit into a power of $t$, a unit of $R[[t]]$, and a polynomial
\[
 1+a_{-1}t^{-1}+\cdots+a_{-m}t^{-m}
 \quad\text{with nilpotent coefficients }a_{-i}.
\]
Thus the last factor is the increasing union of formal schemes
$W_m=\operatorname{Spf}A[[a_{-1},\ldots,a_{-m}]]$; the two other factors form $J=\mathbb Z_A\times(R\mapsto R[[t]]^\times)$. The transition $W_m\to W_{m+1}$ sets the last coefficient equal to zero. In \S1.5, the shadow replaces each $W_m$ by $\operatorname{Spec}A[[a_{-1},\ldots,a_{-m}]]$, takes their inductive-limit sheaf, and extends the $J$-torsor $\mathcal J\to\varinjlim W_m$ to these shadows. The splitting used to describe it depends on $t$, but the resulting torsor and its Abel map are intrinsic. Definition 1.5.18 constructs that map from the diagonal ideal, trivialized off the diagonal. Equations (1.6.3)--(1.6.5) extend a homomorphism $h$ uniquely to the required map $h_{\mathrm{omb}}$, independently of the chosen parameter. This makes the displayed composition well-defined.

Theorem 1.4.4 first identifies homomorphisms from power-series units with $G(U)/G(\operatorname{Spec}B)$; Theorem 1.6.6 adds the remaining components and obtains the unrestricted section $G(U)$. The noetherian hypothesis is part of this theorem, and is not removed by the arbitrary-base assertion for the distinct global object in subject~26. The full 2013 author paper supplies the cited global restatement and the complete local proof. An unopened 1990 chapter or Huber's original ad\`ele proof is not being substituted for those precise opened theorem inputs.
\end{proof}

\subsubsection{28. Motives over arithmetic bases}
\begin{theorem}[Known: rational triangulated motives and six operations]\label{B:subject28}
For finite-type schemes over $\mathbb Z$ there are the triangulated categories of rational Beilinson motives $\mathrm{DM}_{\mathrm B}(X)$ and their constructible subcategories. For separated finite-type morphisms the six operations preserve constructibility. If $a:X\to\operatorname{Spec}\mathbb Z$ is separated of finite type, then $a^!\mathbf1$ is dualizing on constructible motives: with
\[
 D_X(M)=\underline{\operatorname{Hom}}(M,a^!\mathbf1)
\]
the canonical map $M\to D_XD_XM$ is an isomorphism, and duality exchanges the expected pairs of six operations.
\end{theorem}
\begin{proof}[Derivation]
Cisinski--D\'eglise \cite[Theorem 15.2.1]{BCisinskiDeglise2019} gives constructible stability under $f^*$ and tensor product without a base restriction, under $f_!$ for separated finite-type maps, and under $f_*,f^!$ and internal Hom with the quasi-excellence hypotheses stated there. Schemes of finite type over $\mathbb Z$ are excellent, so these hypotheses hold. Their Theorem 15.2.4 applies with both its regular base $S$ and excellent base $B$ equal to $\operatorname{Spec}\mathbb Z$, of dimension one. Its clauses (a)--(d) then give precisely the dualizing object, biduality and exchange identities displayed here.

This is an unconditional categorical construction over the base requested on the sheet. It is a rational triangulated theory; it does not assert the existence of an unconditional motivic $t$-structure or of the full abelian category of mixed motives. Such an assertion would be an additional conjecture, not a consequence of these theorem locators.
\end{proof}

\subsection{The nine items in the ``Pour m\'emoire'' diagram}
\label{B:pourmemoire}

The diagram is on cote~66, scan~4. We distinguish its nine named items: resolution, purity, regular base change, local duality, finiteness, affine cohomological dimension, Grauert finiteness, rigid analytic theories, and the questions in SGA~1962. The original scan resolves the final year as 1962. The diagram expresses dependencies, rather than nine propositions with all hypotheses already specified. The following are the precise proved endpoints, including the restrictions needed to use them.

\subsubsection{PM1. Resolution of singularities}
\begin{theorem}[Known: characteristic-zero resolution and general prime-to-prime alterations]\label{B:PM1}
Every generically reduced noetherian quasi-excellent scheme $X$ over $\mathbb Q$ admits a finite sequence of blowups along regular centres, all disjoint from the inverse image of $X_{\mathrm{reg}}$, whose final scheme is regular. The sequence can be chosen functorially for regular morphisms, omitting empty blowups after base change.

For a noetherian quasi-excellent scheme $X$, a nowhere dense closed subset $Z$, and a prime $\ell$ invertible on $X$, there is a finite family $\{X_i\to X\}$ covering in the $\ell'$-alteration topology such that each $X_i$ is regular and integral and the inverse image of $Z$ is the support of a strict normal-crossings divisor. This latter assertion does not assert a birational resolution in positive characteristic.
\end{theorem}
\begin{proof}[Derivation and scope]
The first statement is exactly Temkin \cite[Theorem 1.2.1]{BTemkin2012}; its four clauses give regular centres, preservation of the regular locus, termination in a regular scheme and functoriality. A blowup is proper, and its restriction away from its centre is an isomorphism, so the composite is proper and is an isomorphism over $X_{\mathrm{reg}}$. The generically reduced assumption makes this a dense open, giving the birational conclusion intended by resolution.

The second statement is Gabber's prime-to-$\ell$ uniformization theorem in \cite[Expos\'e IX, Theorem 1.1; Introduction, Theorem 2]{BILO2014}. The topology is generated by Nisnevich covers and proper surjective generically finite maps whose generic residue degrees are prime to $\ell$. Thus the theorem permits changes in the function fields and multiple covering members. Those are the precise features that allow the cohomological applications below without assuming birational resolution in all characteristics. The general positive-characteristic resolution problem is not declared solved by this alteration theorem.
\end{proof}

\subsubsection{PM2. Absolute cohomological purity}
\begin{theorem}[Known: Gabber's absolute purity]\label{B:PM2}
Let $n\geq1$ be invertible on a noetherian regular scheme $X$, and let $i:Z\hookrightarrow X$ be a closed regular subscheme of pure codimension $c$. With $\Lambda=\mathbb Z/n\mathbb Z$, the cycle-class Gysin map is an isomorphism
\[
 \Lambda_Z(-c)[-2c]\xrightarrow{\sim}i^!\Lambda_X.
\]
Consequently $R^qi^!\Lambda_X=0$ for $q\ne2c$, and $R^{2c}i^!\Lambda_X=\Lambda_Z(-c)$.
\end{theorem}
\begin{proof}[Derivation]
This is \cite[Expos\'e XVI, Theorem 3.1.1, p.~222]{BILO2014}, which writes the same map as $\Lambda\to i^!\Lambda(c)[2c]$. Undoing its twist and shift gives the displayed formula. Taking cohomology sheaves gives the asserted concentration. Both the ambient and the closed scheme are regular; codimension and invertibility are hypotheses of the theorem, rather than consequences of the word ``purity'' on the sheet. The local-cohomology form follows by the identification $R\mathcal H_Z(\Lambda_X)=i_*i^!\Lambda_X$.
\end{proof}

\subsubsection{PM3. Regular base change}
\begin{theorem}[Known: the regular base-change theorem]\label{B:PM3}
Let $f:Y\to X$ be a regular morphism of locally noetherian $\mathbb Z[1/n]$-schemes, let $g:X'\to X$ be quasi-compact and quasi-separated, and form $Y'=Y\times_XX'$ with projections $g':Y'\to Y$ and $f':Y'\to X'$. For every $K\in D^+(X'_{\acute et},\mathbb Z/n\mathbb Z)$ the canonical map
\[
 f^*Rg_*K\xrightarrow{\sim}Rg'_*f'^*K
\]
is an isomorphism. In particular this applies to the completion map of an excellent noetherian local ring.
\end{theorem}
\begin{proof}[Derivation]
The first assertion is \cite[Expos\'e XIV, Theorem 2.5.3; Expos\'e XVII, Proposition 4.2.1, p.~264]{BILO2014}. It includes complexes bounded below, without a constructibility hypothesis on $K$. A regular morphism is flat with geometrically regular fibres. Excellence supplies exactly this property for the map from a noetherian local ring to its completion, proving the stated application. The proof in the cited exposition reduces locally to the expression of a regular algebra as a filtered colimit of smooth algebras, applies smooth base change, and passes to the limit in the quasi-compact, quasi-separated setting. Thus the diagram's arrow is a generalization of smooth base change, not a claim that arbitrary flat base change has this property.
\end{proof}

\subsubsection{PM4. Local duality and biduality}
\begin{theorem}[Known: the dualizing-complex and biduality theorem]\label{B:PM4}
Let $X$ be an excellent noetherian $\mathbb Z[1/n]$-scheme with a dimension function $\delta$, and put $\Lambda=\mathbb Z/n\mathbb Z$. There is a potential dualizing complex $K_{X,\delta}$, uniquely determined together with its compatible local normalizations. It belongs to $D^b_{ctf}(X_{\acute et},\Lambda)$, and
\[
 D_X(M)=R\mathcal Hom_\Lambda(M,K_{X,\delta})
\]
preserves $D^b_c(X_{\acute et},\Lambda)$. For every $M$ in that category, $M\to D_XD_X(M)$ is an isomorphism. If $X$ has finite Krull dimension, $K_{X,\delta}$ has finite quasi-injective dimension and is a dualizing complex in the sense of SGA~5.
\end{theorem}
\begin{proof}[Derivation]
Existence and uniqueness with the local normalizations are \cite[Expos\'e XVII, Theorem 5.1.1, p.~269]{BILO2014}. All four further assertions are the clauses of its Theorem 6.1.1, p.~275. In particular, the argument does not require the tested object $M$ itself to have finite Tor dimension: bounded constructibility suffices. Applying the fourth clause to each such $M$ is exactly the biduality requested on the diagram. The finite-dimensional specialization supplies the usual finite quasi-injective-dimension definition of a dualizing complex. Forgetting the local normalizations would invalidate the phrase ``unique up to unique isomorphism'', since an unnormalized dualizing complex has scalar automorphisms; we have retained that data.
\end{proof}

\subsubsection{PM5. Finiteness for finite-type morphisms}
\begin{theorem}[Known, with the necessary coefficient restriction]\label{B:PM5}
Let $X$ be noetherian and quasi-excellent, let $f:Y\to X$ be of finite type, and let $n$ be invertible on $X$. For a constructible $\mathbb Z/n\mathbb Z$-sheaf $F$ on $Y_{\acute et}$, every $R^qf_*F$ is constructible and these sheaves vanish for all sufficiently large $q$. Thus $Rf_*$ preserves bounded constructible complexes. Without invertibility, even finiteness of $H^1$ over an algebraically closed field can fail.
\end{theorem}
\begin{proof}[Derivation and the excluded case]
The sheaf assertions are \cite[Expos\'e XIII, Theorem 1.1.1, p.~179]{BILO2014}. For a bounded constructible complex use its finite cohomological filtration and the resulting hypercohomology spectral sequence: only finitely many constructible cohomology sheaves occur, each has bounded constructible direct image, so the total direct image is bounded constructible as well.

To verify the coefficient restriction concretely, let $k$ be algebraically closed of characteristic $p>0$ and $Y=\mathbb A^1_k$. Artin--Schreier and affine vanishing for quasi-coherent sheaves give
\[
 H^1_{\acute et}(Y,\mathbb F_p)
   \simeq k[t]/\{a(t)^p-a(t):a(t)\in k[t]\}.
\]
The classes of $t^m$ for distinct positive $m$ prime to $p$ are $\mathbb F_p$-linearly independent. Indeed, a nonzero finite linear combination has highest degree prime to $p$, whereas a nonconstant $a^p-a$ has highest degree divisible by $p$; a constant cannot produce the combination. Thus this cohomology group is infinite dimensional. For $f:Y\to\operatorname{Spec}k$, the sheaf $R^1f_*\mathbb F_p$ is not constructible. This identifies the exact failed hypothesis rather than extending the finite-type theorem to all torsion coefficients.
\end{proof}

\subsubsection{PM6. Cohomological dimension of affine morphisms}
\begin{theorem}[Known: the relative affine Lefschetz bound]\label{B:PM6}
Let $X$ be a quasi-excellent scheme equipped with a dimension function $\delta_X$, let $n$ be invertible on $X$, let $f:Y\to X$ be affine of finite type, and let $F$ be a constructible $\mathbb Z/n\mathbb Z$-sheaf on $Y$. Put
\[
 \delta_Y(y)=\delta_X(f(y))+operatorname{trdeg}_{\kappa(f(y))}\kappa(y),
 \qquad
 \delta_X(G)=\sup_{G_{\bar x}\ne0}\delta_X(x).
\]
Then $\delta_X(R^qf_*F)\leq\delta_Y(F)-q$. In particular, for an affine $d$-dimensional scheme of finite type over an algebraically closed field, cohomology with constructible prime-to-characteristic torsion coefficients vanishes above degree $d$.
\end{theorem}
\begin{proof}[Derivation]
This is \cite[Expos\'e XV, Theorem 1.1.2, p.~201]{BILO2014}, with its definition 1.1.1 of the dimension of a sheaf. For the specialization take $X=\operatorname{Spec}k$ and $\delta_X=0$. A nonzero sheaf on this point has dimension zero, whereas $\delta_Y(F)\leq\dim Y=d$. If $q>d$, the displayed inequality forces $R^qf_*F=0$. Over the algebraically closed point this is $H^q(Y,F)$. The formulation includes the support-sensitive relative bound, which is stronger than the absolute dimension estimate alone.
\end{proof}

\subsubsection{PM7. Grauert-type finiteness}
\begin{theorem}[Known: coherent proper direct image in complex analysis]\label{B:PM7}
For a proper holomorphic map $f:X\to S$ of complex spaces and a coherent analytic sheaf $F$ on $X$, every $R^qf_*F$ is coherent. In particular $H^q(X,F)$ is finite dimensional when $X$ is compact. This does not, without further hypotheses, say that $R^qf_*F$ is locally free or that arbitrary base change commutes with it.
\end{theorem}
\begin{proof}[Derivation]
The proper-direct-image assertion is Grauert \cite[\S6, Hauptsatz I, p.~59]{BGrauert1960}. If $S$ is a point, a coherent $\mathcal O_S$-module is a finite-dimensional complex vector space and its direct-image stalk is $H^q(X,F)$, giving the stated consequence. Local freeness is a stronger property than coherence, and is not included in the invoked theorem. Grauert's subsequent base-change results require their own flatness or fibre-dimension hypotheses; none is silently added here.
\end{proof}

\subsubsection{PM8. Rigid analytic theories}
\begin{theorem}[Known: an affinoid cohomological endpoint]\label{B:PM8}
Let $K$ be a complete algebraically closed nonarchimedean field with a nontrivial valuation, and let $A$ be a classical $K$-affinoid algebra of dimension $d$. For every torsion abelian sheaf $F$ on $(\operatorname{Spec}A)_{\acute et}$,
\[
 H^i(\operatorname{Spec}A,F)=0\quad(i>d+1).
\]
If $F$ is $\ell$-primary and $\ell$ is different from the residue characteristic of $K$, the stronger bound is $i>d$.
\end{theorem}
\begin{proof}[Derivation and site convention]
These are the two clauses of Bhatt--Mathew \cite[Theorem 7.3]{BBhattMathew2021}. Their notation 7.2 fixes the nontrivial valuation and a topologically nilpotent unit. A classical affinoid algebra is topologically of finite type over $K$, so the hypotheses agree exactly. The statement is deliberately on the algebraic \'etale site of $\operatorname{Spec}A$; their \S6.3 proves the affinoid comparison through descent and explains its identification with the corresponding rigid analytic \'etale theory. The theorem is a precise modern endpoint for the diagram's open-ended phrase ``rigid analytic theories''. Its second bound is not claimed for arbitrary residue-characteristic torsion by citing this theorem: the source expressly separates the two clauses.
\end{proof}

\subsubsection{PM9. The questions in SGA 1962}
\begin{theorem}[Known: the local-to-global homotopical-depth conjecture]\label{B:PM9}
Let $(X,Y)$ be a Hausdorff polyhedral pair, meaning a pair homeomorphic to a locally finite simplicial complex and a subcomplex, with $X$ second countable. Let $n\geq1$. Suppose every $y\in Y$ has a fundamental system of neighbourhoods $U$ for which $(U,U\setminus Y)$ is $(n-1)$-connected. Then $(X,X\setminus Y)$ is $(n-1)$-connected. The implication therefore applies to second-countable finite-dimensional complex analytic spaces and closed analytic subspaces whenever the stated local depth bound holds.
\end{theorem}
\begin{proof}[Derivation and identification of the seminar]
The seminar is SGA~2, whose title explicitly identifies the year 1962. Its Expos\'e XIII contains the problems and conjectures, and its \S6, Definition 1 and the following discussion formulate this local-to-global question \cite{BSGA2}. Eyral \cite[Theorem 2.4, p.~828]{BEyral2000} proves the displayed implication. His Proposition 2.2 identifies the neighbourhood condition with the homotopical depth definition in the polyhedral setting; his conventions specify local finiteness, Hausdorffness and second countability as retained above. The analytic application follows from the compatible triangulation used in his Proposition 1.9 and Theorem 3.11.

The source also supplies a concrete mechanism for the passage from local to global. Choose a triangulation $(K,L)$ of the pair. Intersections of finitely many open stars at vertices of $L$ are very good neighbourhoods and satisfy the same relative connectivity bound. The excision lemmas 4.1.3--4.1.4 show inductively that finite unions of those stars, and their intersections with any finite star intersection, retain the bound. Enumerate a countable star cover of $L$ and take its union $V$. Every relative sphere or disk used to test connectivity has compact image, hence lies in a finite partial union; the connectivity of those unions gives the bound for $(V,V\setminus |L|)$. Excision for the open covering $|K|=V\cup(|K|\setminus |L|)$ then gives the result. This also explains why local information alone is not being used without a gluing theorem.

The annotated SGA~2 edition further gives the exact later locators for its analytic Lefschetz programme: Hamm--L\^e, Theorem 3.2.1 for purity, 3.3.1 for local Lefschetz, 3.4.1 for global Lefschetz, and 3.5.1 for the fundamental group. Its notes also flag that Conjecture A, in its original local-Hurewicz form, is false and needs the corrected Theorem 3.1.4. Those annotations were opened; the full Hamm--L\^e chapter was not obtained and is not used as a proved input here. Proposition~\ref{B:local-Hurewicz} independently refutes the original topological statement and proves the exact correction for finite conical links. Thus the line ``Questions dans SGA 1962'' has been identified with its actual seminar, a fully sourced proved conjecture and a complete counterexample to another, without asserting that every sentence of that seminar is one theorem.
\end{proof}

\begin{proposition}[Refuted: the original local Hurewicz conjecture; exact conical repair]\label{B:local-Hurewicz}
There is a finite polyhedral pair $(X,Y)$, with $Y$ a single point $v$, such that
\[
 \mathcal H^i_Y(\mathbb Z_X)=0\quad(i<4),\qquad
 \pi_2^Y(X,v)=0,\qquad
 \pi_3^Y(X,v)=\mathbb Z/2.
\]
It refutes SGA~2, Expos\'e XIII, \S6, Conjecture A, which is stated for topological locally triangulable pairs and includes the local fundamental-group condition.

For an arbitrary nonempty connected finite polyhedron $L$, let $X=cL$ be its cone and $Y=\{v\}$ its vertex. For every $n\geq3$, the exact replacement is
\[
 \operatorname{prof\,hpt}_{Y}(X)\geq n
 \quad\Longleftrightarrow\quad
 \begin{cases}
  \pi_1(L)=0,\\
  \widetilde H^j(L,\mathbb Z)=0& (j<n-1),\\
  H^{n-1}(L,\mathbb Z)\text{ is torsion-free}.
 \end{cases}
\]
The last condition is equivalently torsion-freeness of
$\mathcal H^n_Y(\mathbb Z_X)_v$.
\end{proposition}
\begin{proof}
Here is a finite construction, including the integral torsion calculation. Let $P$ have vertices $1,\ldots,6$ and maximal simplices
\[
 123,\ 124,\ 135,\ 146,\ 156,\
 236,\ 245,\ 256,\ 345,\ 346.
\]
Its one-skeleton is the complete graph, so $P$ is connected. The five edges $1i$ form a spanning tree; projection onto the other edges identifies the integral cycle lattice with $\mathbb Z^{10}$. Denote those chord coordinates by
\[
 (a,b,c,d,e,f,g,h,i,j)=(23,24,25,26,34,35,36,45,46,56).
\]
The ten oriented face boundaries impose
\[
 a=b=f=i=j=0,\qquad
 a-d+g=b-c+h=c-d+j=e-f+h=e-g+i=0.
\]
They reduce to one generator $d$ with the sole relation $2d=0$. The boundary matrix in these cycle coordinates has determinant $-2$, so its kernel is zero. Thus
\[
 \widetilde H_k(P,\mathbb Z)=
 \begin{cases}\mathbb Z/2&k=1,\\0&k\ne1.\end{cases}
\]
The attached program constructs all edges and boundary matrices, checks that their composite is zero, and obtains Smith diagonal $(1,\ldots,1,2)$. This verifies the finite calculation; the preceding elimination also proves it directly.

Take the simplicial suspension $L=\Sigma P$, using two new vertices. Two thickened cones cover $L$ and have path-connected intersection homotopy equivalent to $P$. Van Kampen gives $\pi_1(L)=0$; Mayer--Vietoris gives
$\widetilde H_2(L,\mathbb Z)=\mathbb Z/2$ and zero reduced homology in all other degrees. The Hurewicz theorem gives $\pi_2(L)=\mathbb Z/2$; integral universal coefficients give
\[
 \widetilde H^j(L,\mathbb Z)=0\ (j\leq2),\qquad
 H^3(L,\mathbb Z)=\mathbb Z/2.
\]
The inputs and their proofs were opened in Hatcher \cite[Theorem 2.35, Example 2.40, Theorems 3.2 and 4.32]{BHatcher2002}.

Now take $X=cL$ and $Y=\{v\}$. Radially smaller cones give a cofinal neighborhood system at $v$; their punctures retract onto $L$ and all transition maps are homotopy equivalences. The original local homotopy definition therefore gives
$\pi_i^Y(X,v)=\pi_{i-1}(L)$ for $i\geq2$. In the polyhedral category, constant-sheaf and singular cohomology agree, as computed by open-star covers. Excision and the long exact sequence of the contractible cone and its puncture give
\[
 \mathcal H^i_Y(\mathbb Z_X)_v
   =\widetilde H^{i-1}(L,\mathbb Z)\quad(i\geq1),
 \qquad \mathcal H^0_Y(\mathbb Z_X)=0.
\]
Outside $v$ these supported sheaves are zero. All hypotheses of the original conjecture hold for $n=4$, including local triangulability and local fundamental-group vanishing, while $\pi_3^Y(X,v)$ is nonzero. No analytic realization of the pair is being claimed or required by that topological conjecture.

For the exact replacement, the left side means that $L$ is $(n-2)$-connected. If this holds, Hurewicz and universal coefficients give the conditions on the right. Conversely, suppose those conditions hold. For $n=3$, simple connectivity already proves the assertion, and $H^2(L,\mathbb Z)=\operatorname{Hom}(H_2(L,\mathbb Z),\mathbb Z)$ is automatically torsion-free. For $n\geq4$, all homology groups of the finite polyhedron are finitely generated. The short exact sequences
\[
 0\longrightarrow\operatorname{Ext}^1(H_{j-1}(L,\mathbb Z),\mathbb Z)
 \longrightarrow H^j(L,\mathbb Z)
 \longrightarrow\operatorname{Hom}(H_j(L,\mathbb Z),\mathbb Z)
 \longrightarrow0
\]
show, by applying them in consecutive degrees, that $H_j(L,\mathbb Z)=0$ for $1\leq j\leq n-3$, while $H_{n-2}(L,\mathbb Z)$ is finite torsion. In degree $n-1$, its Ext group is precisely the torsion subgroup of $H^{n-1}(L,\mathbb Z)$, since the Hom quotient is free. The last hypothesis therefore kills $H_{n-2}(L,\mathbb Z)$. Starting with $\pi_1(L)=0$, successive Hurewicz isomorphisms now kill every $\pi_j(L)$ for $j\leq n-2$. This proves the required connectivity and both directions of the criterion.
\end{proof}
